% concatenated from Regularity.tex
\documentclass[11pt,reqno]{amsart}

\usepackage[T1]{fontenc}
\usepackage{lmodern}
\usepackage{amsmath,amssymb,amsthm,mathtools,mathrsfs}
\usepackage{booktabs,longtable,array}
\usepackage{enumitem}
\usepackage[margin=1.08in]{geometry}
\usepackage{microtype}
\usepackage[hidelinks]{hyperref}
\usepackage{url}

\allowdisplaybreaks[2]
\numberwithin{equation}{section}
\setlist[itemize]{leftmargin=2em}
\setlist[enumerate]{leftmargin=2.5em}

\newtheorem{theorem}{Theorem}[section]
\newtheorem{proposition}[theorem]{Proposition}
\newtheorem{lemma}[theorem]{Lemma}
\newtheorem{corollary}[theorem]{Corollary}
\newtheorem{claim}[theorem]{Claim}
\theoremstyle{definition}
\newtheorem{definition}[theorem]{Definition}
\theoremstyle{remark}
\newtheorem{remark}[theorem]{Remark}

\newcommand{\C}{\mathbb C}
\newcommand{\R}{\mathbb R}
\newcommand{\Z}{\mathbb Z}
\newcommand{\Q}{\mathbb Q}
\newcommand{\Gm}{\mathbb G_m}
\newcommand{\Spin}{\operatorname{Spin}}
\newcommand{\SO}{\operatorname{SO}}
\newcommand{\Or}{\operatorname{O}}
\newcommand{\GL}{\operatorname{GL}}
\newcommand{\SL}{\operatorname{SL}}
\newcommand{\Sp}{\operatorname{Sp}}
\newcommand{\Sym}{\operatorname{Sym}}
\newcommand{\Hom}{\operatorname{Hom}}
\newcommand{\End}{\operatorname{End}}
\newcommand{\Lie}{\operatorname{Lie}}
\newcommand{\Ad}{\operatorname{Ad}}
\newcommand{\ad}{\operatorname{ad}}
\newcommand{\Tr}{\operatorname{Tr}}
\newcommand{\tr}{\operatorname{tr}}
\newcommand{\rk}{\operatorname{rk}}
\newcommand{\supp}{\operatorname{supp}}
\newcommand{\Supp}{\operatorname{Supp}}
\newcommand{\codim}{\operatorname{codim}}
\newcommand{\pr}{\operatorname{pr}}
\newcommand{\Res}{\operatorname{Res}}
\newcommand{\diag}{\operatorname{diag}}
\newcommand{\Id}{\operatorname{Id}}
\newcommand{\fg}{\mathfrak g}
\newcommand{\fz}{\mathfrak z}
\newcommand{\fh}{\mathfrak h}
\newcommand{\fp}{\mathfrak p}
\newcommand{\cN}{\mathcal N}
\newcommand{\cO}{\mathcal O}
\newcommand{\cS}{\mathcal S}
\newcommand{\cG}{\mathcal G}
\newcommand{\one}{\mathbf 1}
\newcommand{\eps}{\varepsilon}
\newcommand{\muT}{\mu_2}
\newcommand{\absC}[1]{\lvert #1\rvert_{\C}}
\newcommand{\dd}{\,d}
\newcommand{\defect}{\operatorname{def}}
\newcommand{\wt}{\operatorname{wt}}
\newcommand{\sslash}{/\!\!/}
\newcommand{\wedgeTwo}{\bigwedge^{2}}

\hypersetup{
  pdftitle={Regularity and the Gelfand Property for Complex Symmetric Pairs},
  pdfsubject={The Aizenbud--Gourevitch and van Dijk conjectures for complex symmetric pairs},
  pdfkeywords={complex symmetric pair, regularity, Gelfand pair, Schwartz distribution, nilpotent orbit}
}

\title[Regularity of complex symmetric pairs]{Regularity and the Gelfand Property for Complex Symmetric Pairs}
% Insert author information before submission.
\author[Y.F.~Li]{Yufeng Li}
\address{School of Mathematical Sciences, Peking University, No. 5 Yiheyuan Road, Beijing 100871, P.R. China}
\email{2401110016@stu.pku.edu.cn}

\author[J.Y.~Xiao]{Junyan Xiao}
\address{School of Mathematical Sciences, Peking University, No. 5 Yiheyuan Road, Beijing 100871, P.R. China}
\email{jyxiao25@stu.pku.edu.cn}
 
\author[J.~Yu]{Jun Yu}
\address{School of Mathematical Sciences and Beijing International Center for Mathematical Research, Peking University, 
No. 5 Yiheyuan Road, Beijing 100871, P.R. China}
\email{junyu@bicmr.pku.edu.cn}

\date{}
\subjclass[2020]{Primary 20G05; Secondary 20G41, 22E46, 46F10}
\keywords{complex symmetric pair, Gelfand pair, regularity, Schwartz distribution, nilpotent orbit}

\begin{document}
\raggedbottom

\begin{abstract}
We prove that every symmetric pair of a connected complex reductive group is
regular in the sense of Aizenbud--Gourevitch.  This settles the
Aizenbud--Gourevitch regularity conjecture over the complex numbers.  Generalized
Harish--Chandra descent then makes the canonical central cover of every complex
symmetric pair a Gelfand--Kazhdan pair.  An anti-automorphism arising from a
compatible Chevalley involution upgrades the resulting GP2 bound to GP1 on the
cover, and finite central descent transfers GP1 to the original pair.
In particular, van Dijk's conjecture on complex symmetric pairs follows.

Rubio reduced the unresolved irreducible regularity problem to four families: the
DIII family $(D_r,A_{r-1}+\C)$, the balanced CII family
$(C_{2r},C_r+C_r)$, some remaining Spin block pairs, and the EVII pair
$(E_7,E_6+\C)$.  We treat these cases by four different mechanisms.  For DIII
we construct a sign-equivariant Schwartz distribution on the regular set and
extend it across a common orbit boundary by the Chen--Sun theorem.  For balanced
CII we combine homogeneity, distinguished nilpotent orbits, and a stable-density
theorem for the centralizer representation.  For Spin blocks we prove pleasantness for unequal odd--odd blocks, use
Przebinda's orthogonal-distribution theorem in odd smaller rank, and construct
a finite orbit closure with automatic extension in even smaller rank.  For
EVII we compute the graded-$\mathfrak{sl}_2$ data for all twenty-two nilpotent
orbits and use central-torus characters to eliminate the remaining resonances,
including the two residual triple-centralizer cases.  A finite-component
assembly theorem then handles arbitrary connected central quotients and diagonal
couplings among simple factors.
\end{abstract}

\maketitle
\setcounter{tocdepth}{2}
\tableofcontents

\section{Introduction}\label{sec:introduction}

\subsection{The conjectures and the main theorems}

Let $G$ be a connected complex reductive algebraic group and $\theta$ be an
algebraic involution. Put $H=G^\theta$.  Write
\[
  \fg=\fh\oplus\fp
\]
for the $\pm1$ eigenspace decomposition of $d\theta$.  Aizenbud and Gourevitch
introduced regularity of $(G,H,\theta)$ as the local distributional condition
needed in their generalized Harish--Chandra descent.  They conjectured that
every symmetric pair is regular \cite[Conjecture~4]{AG09}.  Aizenbud and
Gourevitch proved that every connected complex symmetric pair is good in their
sense \cite[Corollary~7.1.7]{AG09}.  Rubio classified the cases already covered
by pleasantness or niceness and reduced the remaining complex problem to four
families \cite[Proposition~6.3]{Rubio}.

The first main result closes that reduction for arbitrary connected complex reductive group 
$G$ and the full fixed subgroup $G^\theta$.

\begin{theorem}[Aizenbud--Gourevitch conjecture over $\C$]\label{thm:regularity-all}
Let $G$ be a connected complex reductive group and let $\theta$ be an
involution of $G$.  Then the symmetric pair $(G,G^\theta,\theta)$ is regular.
\end{theorem}

The second result is the global multiplicity-one consequence.  Throughout,
$E$ denotes an irreducible Casselman--Wallach representation of the underlying
real Lie group of $G$, and $E^\vee$ denotes its smooth contragredient.  We use
GP1 for
\[
  \dim\Hom_H(E,\C)\leq1
\]
and GP2 for
\[
  \dim\Hom_H(E,\C)\,\dim\Hom_H(E^\vee,\C)\leq1.
\]
For a symmetric pair, ``Gelfand--Kazhdan'' means that every
$H\times H$-invariant Schwartz distribution on $G$ is invariant under
$g\mapsto\theta(g^{-1})$; the Gelfand--Kazhdan criterion implies GP2.

\begin{theorem}[Gelfand property and van Dijk's conjecture]\label{thm:gelfand-all}
Let $(G,H,\theta)$ be as in Theorem~\ref{thm:regularity-all}, and let
\[
 m:\widetilde G=Z(G)^\circ\times G_{\mathrm{der}}^{\mathrm{sc}}
   \longrightarrow G
\]
be the canonical central cover with lifted involution $\widetilde\theta$.
Then $(\widetilde G,\widetilde G^{\widetilde\theta},\widetilde\theta)$
is a Gelfand--Kazhdan pair.  Moreover,
\[
  \dim\Hom_H(E,\C)\leq 1
\]
for every irreducible Casselman--Wallach representation $E$ of $G$.  Hence every
complex symmetric pair is a unitary Gelfand pair.
\end{theorem}

At the irreducible level, Rubio's reduction leaves the following families.

\begin{theorem}[The residual irreducible families]\label{thm:four-inputs}
The following complex symmetric pairs are regular.
\begin{enumerate}[label=\textup{(\roman*)}]
\item every DIII pair $(D_r,A_{r-1}+\C)$, $r\geq4$, in every connected
isogeny form;
\item every balanced CII pair $(C_{2r},C_r+C_r)$, $r\geq1$;
\item every Spin block pair
\[
 (\Spin_{r+s},\Spin_r\times_{\mu_2}\Spin_s),\qquad r,s>0,\quad r+s\geq5;
\]
\item the EVII pair $(E_7,E_6+\C)$, which is in fact special.
\end{enumerate}
\end{theorem}

Two further group-theoretic ingredients are needed for the passage from
regularity to the Gelfand property.  Proposition~\ref{spin:prop:pleasant} shows that every
unequal odd--odd Spin block is pleasant.  Proposition~\ref{prop:compatible-chevalley}
constructs a Chevalley involution compatible with the symmetric involution and
yields the GP1--GP2 equivalence required on the canonical central cover.

\subsection{Proof architecture}

The proof separates the infinitesimal regularity problem from the passage
between isogeny forms and full fixed subgroups.

For a simple simply connected factor, Rubio's classification leaves precisely
the four families of Theorem~\ref{thm:four-inputs}.  Sections
\ref{sec:DIII-family}--\ref{sec:EVII-family} prove them.  The four families require different arguments, reflecting the geometry of
their isotropy representations.

A finite central quotient can enlarge the full fixed subgroup and can couple
different simple factors diagonally.  Section~\ref{sec:component-assembly}
handles this phenomenon through the finite abelian group of effective central
multipliers, a character-sector decomposition of invariant distributions, and
product and intermediate-subgroup assembly theorems.  Together they extend the four simply connected calculations to arbitrary
connected central quotients and full fixed subgroups.

Finally, Section~\ref{sec:completion} applies generalized Harish--Chandra
descent.  We work first on the canonical central cover
$Z(G)^\circ\times G_{\mathrm{der}}^{\mathrm{sc}}$, whose semisimple
centralizers are connected.  Theorem~\ref{thm:regularity-all} therefore applies
to every descendant.  The Aizenbud--Gourevitch criterion gives GP2; the
compatible Chevalley anti-automorphism gives GP1; and GP1 descends through the
finite central isogeny.

\subsection{Conventions}

All algebraic groups, Lie algebras, and finite-dimensional representations are
over $\C$.  When Schwartz spaces, distributions, normal bundles, or Frobenius
reciprocity occur, the complex algebraic variety is regarded as its underlying
real Nash manifold.  Thus complexification of a real normal space contains
both holomorphic and antiholomorphic summands.  We use
$|z|_{\C}=z\bar z$.

\subsection{Acknowledgements and disclosure} The proof was developed with assistance from ChatGPT 5.6 Sol. 
The system was used for proof exploration and drafting. The third named author (Jun Yu) would like to thank 
Professors Binyong Sun and Dmitry Gourevitch for helpful communication.  


\section{Regularity, homogeneity, and finite components}\label{sec:framework}

\subsection{The regularity condition and descent}

Let a reductive group $L$ act linearly on a finite-dimensional complex vector
space $V$.  Following \cite[Notation~2.3.10]{AG09}, put
\[
  Q_L(V)=(V/V^L)(\C),
\]
identify this quotient with its canonical $L$-stable linear model, let
$\pi:Q_L(V)\to Q_L(V)\sslash L$ be the categorical quotient, and define
\[
  \Gamma_L(V)=\pi^{-1}(\pi(0)),\qquad
  R_L(V)=Q_L(V)\setminus\Gamma_L(V).
\]
For an isotropy representation, $\Gamma_H(\fp)$ is the nilpotent cone
\cite{KR,AG09}.  We omit the subscript when the acting group is clear.

\begin{definition}[Aizenbud--Gourevitch]\label{def:regularity}
An element $g\in G$ is \emph{admissible} for $(G,H,\theta)$ if
\begin{enumerate}[label=\textup{(\roman*)}]
\item $\Ad(g)$ commutes with $\theta$, equivalently
      $g^{-1}\theta(g)\in Z(G)$;
\item $\Ad(g)|_{\fp}$ normalizes the $H$-action, its square belongs to the
      $H$-action, and it preserves every closed $H$-orbit in $\fp$.
\end{enumerate}
The pair is \emph{regular} if, for every admissible $g$,
\[
 \cS^*(R_H(\fp))^H\subset\cS^*(R_H(\fp))^{\Ad(g)}
 \quad\Longrightarrow\quad
 \cS^*(Q_H(\fp))^H\subset\cS^*(Q_H(\fp))^{\Ad(g)}.
\]
\end{definition}

A descendant of $(G,H,\theta)$ is the symmetric pair
$(G_x,H_x,\theta|_{G_x})$ attached to a semisimple element
$x\in\{g\theta(g)^{-1}:g\in G\}$, with $G_x$ and $H_x$ its centralizers.
A connected complex symmetric pair is good in the sense of
Aizenbud--Gourevitch \cite[Corollary~7.1.7]{AG09}, and a good symmetric pair
whose descendants, including itself, are regular is a Gelfand--Kazhdan pair
\cite[Theorem~7.4.5]{AG09}.  Products of regular pairs are regular
\cite[Proposition~7.4.4]{AG09}.

A pair is \emph{pleasant} when every inner automorphism commuting with $\theta$
is induced, on the adjoint level, by the fixed subgroup.  Pleasant pairs are
regular \cite[Lemma~4.2]{Rubio}.  A pair is \emph{nice} in Sekiguchi's sense
when its distinguished defect is negative; nice pairs are special and hence
regular \cite{Sekiguchi,Aiz13}.  Rubio's classification of pleasant and nice
complex pairs, together with his descendant computation, leaves the four
families stated in Theorem~\ref{thm:four-inputs} \cite[Proposition~6.3]{Rubio}.

\subsection{Schwartz distributions and null cones}
For a Nash manifold $X$, let $\cS(X)$ be its Schwartz space,
$\cS^*(X)$ its continuous dual, and $\cG(X)$ the space of generalized
Schwartz functions in the sense of \cite[Appendix~B]{AG09}.

Let $L$ be a reductive algebraic group and $V$ an $L$-module.  Put
\[
 Q(V)=V/V^L,
\]
and let $\Gamma(V)\subset Q(V)$ be the null cone, namely the set of vectors
whose $L$-orbit closure contains $0$.  Write
\[
 R(V)=Q(V)\setminus\Gamma(V).
\]
For isotropy representations of complex symmetric pairs, the null cone is the
nilpotent cone by Kostant--Rallis \cite{KR}; see also
\cite[Lemma~7.3.8]{AG09}.  If $\chi$ is a character of $L$, then
$\cS^*(X)^{L,\chi}$ denotes the distributions satisfying
$l\cdot T=\chi(l)T$.

\subsection{Adapted homogeneity}

Let $B$ be an $L$-invariant nondegenerate symmetric bilinear form on a complex
vector space $V$ of dimension $n$.  We use the normalized absolute value
\[
 \absC{z}=z\bar z.
\]
Let $\rho(\lambda)$ act on $V$ by $v\mapsto\lambda^{-1}v$.  If $B$ and
$\lambda B$ are isometric for every $\lambda\in\C^\times$, the Weil-factor
character $\delta_B$ of \cite[Notation~5.1.3]{AG09} is trivial.  A
$B$-adapted distribution has one of the homogeneity characters
\begin{align}
 \eta_0(\lambda)&=\absC{\lambda}^{n/2},\label{core:eq:eta0-general}\\
 \eta_1(\lambda)&=\lambda\absC{\lambda}^{n/2},\label{core:eq:eta1-general}
\end{align}
with the Fourier condition in the first case as in
\cite[Definition~5.1.6]{AG09}.  Only the homogeneity law is needed in the
vanishing arguments below.

\begin{lemma}[Homogeneity reduction]\label{core:lem:homogeneity-reduction}
Let $L$ be reductive, let $V^L=0$, let $B$ be an $L$-invariant nondegenerate
symmetric form, and let $\chi$ be a character of $L$.  Assume
\[
 \cS^*(R(V))^{L,\chi}=0.
\]
If $\cS^*(V)^{L,\chi}\ne0$, then there exists a nonzero
$(L,\chi)$-equivariant $B$-adapted distribution supported on $\Gamma(V)$.
The same statement holds for a finite extension of $L$.
\end{lemma}

\begin{proof}
This is the contrapositive of \cite[Theorem~5.2.2]{AG09}, applied to the
one-summand decomposition $Q(V)=V$ and the form $B$.  In the conventions of
\cite{AG09}, reductive groups need not be connected and the equivariance
character is arbitrary.  The theorem is stated for a $G$-invariant
decomposition
\[
 Q_G(V)=W\oplus\bigoplus_i V_i
\]
with each quadratic form $B_i$ invariant under the full acting group.  There is
no separate hypothesis $V^G=0$; our assumption $V^L=0$ is used only to identify
$Q_L(V)$ with $V$.  Thus the one-summand contrapositive says that vanishing on
$R(V)$ together with vanishing of all $\chi$-equivariant adapted distributions
on the null cone implies vanishing on $Q(V)$, and it applies unchanged to the
finite extensions used below.
\end{proof}

\begin{lemma}[The origin]
\label{core:lem:origin-adapted}
Let $V$ be a nonzero complex vector space, let $B$ be a nondegenerate
symmetric bilinear form for which $\delta_B=1$, and let
$\rho(\lambda)v=\lambda^{-1}v$.  No nonzero $B$-adapted Schwartz
distribution on the underlying real vector space of $V$ is supported at the
origin.
\end{lemma}

\begin{proof}
A distribution supported at $0$ is a finite linear combination of jets
\[
  \partial_z^\alpha\partial_{\bar z}^\beta\delta_0.
\]
The chain rule gives
\[
 \rho(\lambda)
 \bigl(\partial_z^\alpha\partial_{\bar z}^\beta\delta_0\bigr)
 =\lambda^{-|\alpha|}\bar\lambda^{-|\beta|}
 \partial_z^\alpha\partial_{\bar z}^\beta\delta_0.
\]
Restrict the homogeneity action to $r\in\R_{>0}$.  Every origin jet then has
weight $r^{-q}$ for an integer $q\geq0$.  If $n=\dim_\C V>0$, the two adapted
characters restrict respectively to
\[
  |r|_\C^{n/2}=r^n,
  \qquad
  r|r|_\C^{n/2}=r^{n+1}.
\]
Both have strictly positive exponent.  Since the characters $r^{-q}$ are
linearly independent on $\R_{>0}$ from $r^n$ and $r^{n+1}$, no nonzero finite
sum of origin jets is adapted.
\end{proof}

\subsection{The orbitwise resonance calculation}

For the remainder of this subsection assume $\fp^H=0$.  This is the situation
for every semisimple isotropy block to which the resonance calculation is
applied below; for a reductive pair, the central fixed summand is removed by
the definition of $Q(\fp)$.

Let $0\ne e\in\fp$ be nilpotent and choose a normal $\mathfrak{sl}_2$-triple
\begin{equation}\label{core:eq:normal-triple-general}
 [d,e]=2e,\qquad [d,f]=-2f,\qquad [e,f]=d,
 \qquad d\in\fh,\quad e,f\in\fp.
\end{equation}
By the algebraic Jacobson--Morozov theorem, choose a homomorphism
$\phi:\SL_2\to G$ whose differential sends the standard triple to
$(e,d,f)$, and put
\begin{equation}\label{core:eq:Dt-algebraic}
 D_t=\phi\!\begin{pmatrix}t&0\\0&t^{-1}\end{pmatrix},
 \qquad t\in\C^\times.
\end{equation}
The cocharacter $t\mapsto D_t$ has differential $d$.  Since
$d\in\fh$, the cocharacters $\theta(D_t)$ and $D_t$ have the same
differential; in characteristic zero an algebraic cocharacter is determined
by its differential.  Hence $D_t\in H$ and $\Ad(D_t)e=t^2e$.  For the
action of $H\times\C^\times$ in which the second factor acts by
$\rho(\lambda)x=\lambda^{-1}x$, the element $(D_t,t^2)$ fixes $e$; see
\cite[\S3.2]{CM} for the algebraic $\mathfrak{sl}_2$ construction.  Put
\begin{equation}\label{core:eq:T-of-e}
 T(e)=\Tr\bigl(\ad(d)|_{\fh^e}\bigr).
\end{equation}
The stabilizer modular character on this one-parameter subgroup is
\begin{equation}\label{core:eq:modular-general}
 \Delta_{(H\times\C^\times)_e}(D_t,t^2)=\absC{t}^{T(e)}.
\end{equation}
With the modular-character convention in the Frobenius formula of
\cite[Theorem~2.5.7]{AG09}, this is the absolute Jacobian of the adjoint
action on the stabilizer Lie algebra.  That Lie algebra is
$\fh^e\oplus\C(d,2)$; $D_t$ fixes the second summand and has determinant
$t^{T(e)}$ on $\fh^e$, giving \eqref{core:eq:modular-general}.

Choose an invariant nondegenerate symmetric form on $\fg$ for which $\fh$ and
$\fp$ are orthogonal.  Invariance gives
\begin{equation}\label{core:eq:normal-duality-general}
 [e,\fh]^\perp=\fp^e,
 \qquad
 \bigl(\fp/[e,\fh]\bigr)^*\simeq\fp^e.
\end{equation}
Put
\[
 N_e:=\fp/[e,\fh].
\]
Equation~\eqref{core:eq:normal-duality-general} identifies $N_e^*$, not
$N_e$, with $\fp^e$.  More precisely, on every highest-weight isotypic component the invariant
form induces an equivariant perfect pairing between the highest-weight
multiplicity space in $\fp^e$ and the surviving lowest-weight multiplicity
space in $N_e$.  If the highest weight is even, the paired highest and lowest
lines have the same $\theta$-parity.  If it is odd, the form pairs the two
oppositely graded copies.  Concretely, write one copy as
$v_+^{\fh},\ldots,v_-^{\fp}$ and the opposite copy as
$w_+^{\fp},\ldots,w_-^{\fh}$, where $+$ and $-$ denote the highest and
lowest lines.  Since $B(\fh,\fp)=0$ and the invariant form pairs opposite
$d$-weights, its pairing between the two highest and two lowest lines has
the pattern
\[
\begin{array}{c|cc}
 & v_-^{\fp} & w_-^{\fh}\\ \hline
 v_+^{\fh} & 0 & *\\
 w_+^{\fp} & * & 0
\end{array}
\]
with the two indicated entries nonzero after choosing dual multiplicity
bases.  Thus a surviving lowest line from an $\fh$-highest copy is paired
with the highest line of the corresponding $\fp$-highest copy.
Consequently the $d$-weight and every commuting-torus
character are inverted
on passage from $\fp^e$ to $N_e$.  The Bruhat filtration is formulated with the conormal
bundle $CN_{H\cdot e}^{\fp}$; after Frobenius reciprocity and dualization,
its point fiber is the symmetric algebra of $N_e$.  All degrees below are
therefore degrees on the surviving lowest lines in $N_e$.

\begin{lemma}[Resonance formula]\label{core:lem:resonance-general}
Let an irreducible graded $\mathfrak{sl}_2$-summand have highest weight
$\ell$.  Its contribution to $N_e$ is one-dimensional exactly when its lowest
line lies in $\fp$.  In that case the quotient is represented by the lowest
line, has intrinsic $d$-weight $-\ell$, and $(D_t,t^2)$ acts on it by
\begin{equation}\label{core:eq:elementary-degree-general}
 t^{-(\ell+2)}.
\end{equation}

Assume $\delta_B=1$ and let $n=\dim_{\C}\fp$.  Let $T$ be an
$H$-invariant $B$-adapted distribution, and suppose that a nonzero term in its
Bruhat--Frobenius normal symbol along $H\cdot e$ has total holomorphic and
antiholomorphic weighted degrees $(A,B')$, where a surviving quotient line of
highest weight $\ell$ is assigned weighted degree $\ell+2$.  Then necessarily
\begin{align}
 \textup{Type I:}&\qquad A=B'=T(e)-n,\label{core:eq:resonance-I-general}\\
 \textup{Type II:}&\qquad (A,B')=(T(e)-n-2,T(e)-n).
 \label{core:eq:resonance-II-general}
\end{align}
The same necessary equations apply to a $(K,\chi)$-equivariant distribution
for a finite algebraic extension $K$ of the effective image of $H$ whenever
$\chi$ is trivial on that image, since such a distribution is $H$-invariant.
\end{lemma}

\begin{proof}
On an irreducible highest-weight module, $e$ maps every weight line except the
highest line isomorphically to the next higher line.  Since $e\in\fp$, the map
reverses parity.  Hence the quotient of the $\fp$-part by $e$ applied to the
$\fh$-part is zero unless the lowest line lies in $\fp$; in that case exactly
that line survives.  Its $D_t$-weight is $t^{-\ell}$, and the external
homothety contributes $t^{-2}$.

For Type I,
$\eta_0(t^2)=(t\bar t)^n$.  A normal tensor of bidegree $(A,B')$ contributes
$t^{-A}\bar t^{-B'}$, while \eqref{core:eq:modular-general} contributes
$(t\bar t)^{T(e)}$.  Triviality of the stabilizer character gives
$A=B'=T(e)-n$.  For Type II, the additional factor $t^2$ in
$\eta_1(t^2)$ subtracts two from the holomorphic target, giving
\eqref{core:eq:resonance-II-general}.
\end{proof}

The convention here is the same one written out later in
\eqref{cii:eq:exact-frobenius-fiber}.  The raw conormal Frobenius coefficient
carries $\Delta_{L_e}^{-1}$; dualizing a distribution supported at the point
changes this to $\Delta_{L_e}$ and changes $\Sym(N_e^*)$ to $\Sym(N_e)$.
Together with $N_e^*\simeq\fp^e$ from
\eqref{core:eq:normal-duality-general}, this is exactly why both the
$\mathfrak{sl}_2$ weights and commuting-torus characters are inverted in the
normal variables used in the resonance equations.

\begin{corollary}[Strict trace inequality]\label{core:cor:strict-trace}
Assume $\fp^H=0$ and $e\ne0$.  If $T(e)<\dim_{\C}\fp$, the orbit $H\cdot e$ cannot be
maximal in the support of a nonzero $H$-invariant $B$-adapted distribution
supported on the nilpotent cone.  The same conclusion holds for a
$(K,\chi)$-equivariant adapted distribution under a finite algebraic extension
$K$ of the effective image of $H$, provided $\chi$ is trivial on that image.
\end{corollary}

\begin{proof}
Every normal weighted degree is nonnegative, whereas the right-hand side of
\eqref{core:eq:resonance-I-general} is negative.  Type II is more restrictive.
Additional equivariance can only reduce the Frobenius fiber.
\end{proof}

The passage from orbitwise vanishing to distributional vanishing uses
\cite[Theorem~2.5.6]{AG09}.  In the form needed here, let $K$ be a Nash group
acting Nashly on a Nash manifold $M$, and let a locally closed subset
$N\subset M$ admit a finite Nash $K$-invariant stratification
$N=\bigcup_iN_i$.  If
\[
 \cS^*\!\left(N_i,\Sym^k CN_{N_i}^{M}\right)^{K,\chi}=0
 \qquad\text{for every }i\text{ and every }k\ge0,
\]
then $\cS_M^*(N)^{K,\chi}=0$.  Frobenius reciprocity
\cite[Theorem~2.5.7]{AG09} identifies each equivariant conormal-distribution
space with the corresponding stabilizer-equivariant symmetric tensors in the
normal quotient, with the standard modular-character and equivariance twists.
This is why the calculations below are made on $N_e$ and why the modular
character is retained explicitly.  In the one place where
we use a boundary-residue obstruction rather than vanishing of the conormal
space itself, the argument is local on a slice: an ordinary distribution has
finite order on a relatively compact coordinate neighborhood, so a nonzero
local restriction has a highest nonzero transverse order and the resulting
finite descent in transverse order is legitimate.

\subsection{Speciality and regularity}

\begin{lemma}[Aizenbud--Gourevitch speciality criterion]
\label{core:lem:speciality-criterion}
Let $L$ be a reductive algebraic group over $\C$ and let $V$ be a
finite-dimensional algebraic $L$-module.  Suppose there exists an
$L$-invariant decomposition
\[
 Q(V)=\bigoplus_i W_i
\]
and invariant nondegenerate symmetric forms $B_i$ on $W_i$ such that every
$L$-invariant distribution supported on $\Gamma(V)$ and adapted to all $B_i$
is zero.  Then $V$ is special.  A special symmetric pair is weakly linearly
tame and hence regular.
\end{lemma}

\begin{proof}
The first assertion is \cite[Lemma~6.0.10]{AG09}.  The remaining implications
are \cite[Proposition~7.3.5 and Remark~7.4.3]{AG09}.
\end{proof}

Products of special pairs are special
\cite[Proposition~7.3.6]{AG09}, and products of regular pairs are regular
\cite[Proposition~7.4.4]{AG09}.


\section{Central covers and finite-component assembly}\label{sec:component-assembly}

\subsection{Canonical covers, fixed vectors, and multiplier groups}

\begin{lemma}[Canonical central cover]\label{lem:canonical-cover}
Let $(G,G^\theta,\theta)$ be a connected reductive complex symmetric pair.
Multiplication gives a finite central isogeny
\[
 m:Z(G)^\circ\times G_{\mathrm{der}}^{\mathrm{sc}}\longrightarrow G.
\]
The involution $\theta$ lifts to an involution $\widetilde\theta$ of the source,
$m$ is equivariant, and $\ker m$ is $\widetilde\theta$-stable.
\end{lemma}

\begin{proof}
The characteristic subgroups $Z(G)^\circ$ and $G_{\mathrm{der}}$ are
$\theta$-stable.  The restriction to the derived group lifts uniquely to its
simply connected cover, while the restriction to the central torus is already
an automorphism there.  The two lifts commute with multiplication.  The usual
decomposition $G=Z(G)^\circ G_{\mathrm{der}}$ makes $m$ surjective with finite
central kernel, and equivariance makes the kernel stable.
\end{proof}

\begin{lemma}[Fixed vectors]\label{lem:fixed-vectors}
Let $\fg=\fz\oplus\fg_{\mathrm{der}}$ and
$\fp=\fp_{\fz}\oplus\fp_{\mathrm{der}}$.  Then
\[
  \fp^{G^\theta}=\fp^{(G^\theta)^\circ}=\fp_{\fz}.
\]
Consequently $Q_{G^\theta}(\fp)$ is canonically $\fp_{\mathrm{der}}$, and every
inner automorphism acts trivially on the discarded central summand.
\end{lemma}

\begin{proof}
The inclusion $\fp_{\fz}\subset\fp^{G^\theta}$ is immediate.  Let
$x\in\fp_{\mathrm{der}}$ be fixed by $(G^\theta)^\circ$, so
$[\fh\cap\fg_{\mathrm{der}},x]=0$.  Choose a nondegenerate invariant
$\theta$-invariant form $B$ on the semisimple algebra $\fg_{\mathrm{der}}$.
Its restrictions to the two eigenspaces are nondegenerate.  For
$y\in\fp_{\mathrm{der}}$ and $z\in\fh\cap\fg_{\mathrm{der}}$,
\[
  B([x,y],z)=-B([x,z],y)=0.
\]
Because $[x,y]\in\fh\cap\fg_{\mathrm{der}}$, it follows that $[x,y]=0$.
Thus $x$ commutes with both eigenspaces of $\fg_{\mathrm{der}}$ and is central
in a semisimple Lie algebra; hence $x=0$.
\end{proof}

Work now on the canonical cover
\[
 \widetilde G=Z\times G_{\mathrm{ss}}
\]
from Lemma~\ref{lem:canonical-cover}, where $G_{\mathrm{ss}}$ is simply
connected semisimple.  Put $\widetilde H=\widetilde G^{\widetilde\theta}$ and
$V=Q_{\widetilde H}(\fp)$.  The involution preserves the two characteristic
factors.  The fixed subgroup $G_{\mathrm{ss}}^{\widetilde\theta}$ is connected
by the Steinberg--Borel connectedness theorem \cite{Borel,Steinberg}, while all
components of $Z^{\widetilde\theta}$ act trivially on $V$.  Hence the effective
image
\[
 L=\Ad(\widetilde H)|_V
\]
is connected.

Set
\[
 \widetilde A_{\widetilde\theta}
 =\{g\in\widetilde G:g^{-1}\widetilde\theta(g)\in Z(\widetilde G)\},
 \qquad
 J=\Ad(\widetilde A_{\widetilde\theta})|_V,
 \qquad A=J/L.
\]
The group $J$ is a reductive algebraic group: it is the effective image on
$V$ of the fixed subgroup of the induced involution on the connected adjoint
group.  Its identity component is $L$.  Consequently
$A=J/L=\pi_0(J)$ is exactly the finite component group of this effective
algebraic action.
It is also abelian.  Indeed, on $\widetilde A_{\widetilde\theta}$ the multiplier
\[
 c(g)=g^{-1}\widetilde\theta(g)
\]
is a homomorphism because its values are central.  Thus
$c([g_1,g_2])=1$, so $[g_1,g_2]\in\widetilde H$ and the images of $g_1,g_2$
commute modulo $L$.  We write $B_0=\{1\}\leq A$ for the effective component
image of $\widetilde H$.

The following elementary lifting point is needed when an admissible element is
given only after a central quotient.

\begin{lemma}[Lifting central multipliers]\label{lem:lifting-multiplier}
Let $m:\widetilde G\to G$ be a central isogeny of connected complex algebraic
groups, equivariant for involutions $\widetilde\theta$ and $\theta$.  If
$\bar g\in G$ satisfies
$\bar g^{-1}\theta(\bar g)\in Z(G)$, then every lift $g\in\widetilde G$ of
$\bar g$ satisfies
\[
 g^{-1}\widetilde\theta(g)\in Z(\widetilde G).
\]
\end{lemma}

\begin{proof}
Put $y=g^{-1}\widetilde\theta(g)$.  Since $m(y)$ is central, for every
$x\in\widetilde G$ the commutator $[y,x]$ lies in the finite group $\ker m$.
The algebraic map $x\mapsto[y,x]$ from the connected variety
$\widetilde G$ to this finite set is constant.  Its value at the identity is
$1$, so $y$ is central.
\end{proof}

Let $F\subset Z(\widetilde G)$ be finite and $\widetilde\theta$-stable, put
$G=\widetilde G/F$, and let $H=G^\theta$.  The preimage of the full fixed
subgroup is
\begin{equation}\label{eq:preimage-fixed}
 \widehat H_F
 =\{g\in\widetilde G:g^{-1}\widetilde\theta(g)\in F\}.
\end{equation}
Its effective component image is a subgroup
\[
 B_0\leq B_F\leq A.
\]
Indeed, every nonempty fiber of $c:\widehat H_F\to F$ is a left coset of
$\widetilde H$.  By Lemma~\ref{lem:lifting-multiplier}, every inner
transformation commuting with the quotient involution is represented by an
element of $A$.  If it is admissible, the square condition in admissibility
says that its class $a\in A$ satisfies $a^2\in B_F$.

\begin{lemma}[Exact component identification]
\label{lem:exact-component-identification}
Group the simple factors of $G_{\mathrm{ss}}$ by their
$\widetilde\theta$-orbits.  For such an orbit $\Omega$, let
$(L_\Omega,J_\Omega,A_\Omega,V_\Omega)$ denote the corresponding effective
datum.  Then
\[
 V=\bigoplus_\Omega V_\Omega,\qquad
 L=\prod_\Omega L_\Omega,\qquad
 J=\prod_\Omega J_\Omega,\qquad
 A=J/L\simeq\prod_\Omega A_\Omega.
\]
The central torus contributes no effective factor.  Moreover, for every
finite $\widetilde\theta$-stable central subgroup $F$, the effective image of
$\widehat H_F$ in $J$ is exactly the inverse image $J_{B_F}$ of $B_F\leq A$.
Thus a diagonal central coupling is represented precisely by a possibly
non-product subgroup $B_F\leq\prod_\Omega A_\Omega$; there is no further
component group hidden in the quotient.
\end{lemma}

\begin{proof}
On the canonical cover the factor-orbit groups act on pairwise distinct direct
summands of $V$.  The multiplier condition and the fixed-point condition are
factorwise, while the central torus acts trivially after passage to $Q(\fp)$.
Taking effective images therefore gives the displayed direct products and the
product description of $A$.

Since $\widehat H_F\subset\widetilde A_{\widetilde\theta}$, its effective
image is contained in $J$ and has image $B_F$ in $A$.  Conversely, let
$j\in J$ have class in $B_F$.  Choose $h\in\widehat H_F$ with the same class.
Then $j\,\Ad(h)^{-1}\in L$.  The group $L$ is the effective image of
$\widetilde H\subset\widehat H_F$, so $j$ itself lies in the effective image
of $\widehat H_F$.  Hence that image is exactly $J_{B_F}$.
\end{proof}

Already at the level of this finite datum, $B_F$ need not split factorwise.
For example, if two factor orbits have $A_1=A_2=\mu_2$ and a diagonal
central quotient allows precisely equal multiplier classes, then
\[
 B_F=\Delta\mu_2=\{(1,1),(-1,-1)\}
       \subsetneq\mu_2\times\mu_2.
\]
Neither $(-1,1)$ nor $(1,-1)$ belongs to the full fixed-component group.  This
is the elementary finite-component mechanism behind the diagonal couplings
that make a factorwise treatment of the full fixed subgroup insufficient.

The finite component group does not alter the effective null cone.  This point
is needed because regularity is defined using the full fixed subgroup, not only
its identity component.

\begin{lemma}[Finite components and the null cone]
\label{lem:finite-components-nullcone}
Let $B\leq A$, and let $J_B$ be the inverse image of $B$ in $J$.  Then
\[
 V^{J_B}=V^L=0,\qquad Q_{J_B}(V)=Q_L(V)=V,
\]
and
\[
 \Gamma_{J_B}(V)=\Gamma_L(V),\qquad
 R_{J_B}(V)=R_L(V).
\]
Thus the notation $Q_L(V)$ and $R_L(V)$ used below is also the one entering the
regularity condition for every intermediate full fixed subgroup.
\end{lemma}

\begin{proof}
The first assertion follows from $V^L=0$ and $L\subset J_B$.  Moreover
\[
 \C[V]^{J_B}=(\C[V]^L)^B,
\]
so the natural morphism
\[
 V\sslash L\longrightarrow V\sslash J_B=(V\sslash L)\sslash B
\]
is the quotient by the finite group $B$.  Fibers of a finite-group quotient are
precisely finite group orbits; in particular the inverse image of the image of
$0$ is $B\cdot0=\{0\}$.  Pulling this equality back to $V$ proves equality of
the two null cones and hence of their complements.
\end{proof}

\subsection{Characterwise regularity}

For a finite-component datum $(L,A,B_0,V)$ as above and a character
$\lambda\in\widehat A$, write
\[
 \cS^*(X)^{L,\lambda}
 =\{T\in\cS^*(X)^L:j_*T=\lambda(jL)T\text{ for every }j\in J\}.
\]
This is well defined because $L$ acts trivially on $\cS^*(X)^L$.

\begin{definition}\label{def:componentwise-regular}
The datum is \emph{componentwise regular} if, for every nontrivial
$\lambda\in\widehat A$ with $\lambda|_{B_0}=1$,
\begin{equation}\label{eq:component-implication}
 \cS^*(R_L(V))^{L,\lambda}=0
 \quad\Longrightarrow\quad
 \cS^*(Q_L(V))^{L,\lambda}=0.
\end{equation}
\end{definition}

If $A=B_0$, this condition is vacuous.  In particular, pleasant factors create
no characterwise obstruction.  The following observation covers special
factors.

\begin{lemma}[Speciality implies the component implication]
\label{lem:special-component}
Assume that $A$ preserves every summand and every form in a witnessing
orthogonal decomposition for speciality of the connected $L$-module $V$.  Then
$(L,A,B_0,V)$ is componentwise regular.
\end{lemma}

\begin{proof}
Fix a nontrivial character $\lambda\in\widehat A$ that is trivial on $B_0$,
and suppose
\[
 \cS^*(R_L(V))^{L,\lambda}=0.
\]
Regard $\lambda$ as a character of $J$, trivial on $L$.  Assume for
contradiction that
\[
 \cS^*(Q_L(V))^{L,\lambda}
 =\cS^*(Q_L(V))^{J,\lambda}\ne0.
\]
Apply the full multi-summand form of
\cite[Theorem~5.2.2]{AG09} to the finite reductive extension $J$, the
character $\lambda$, and the complete witnessing family
$Q_L(V)=\bigoplus_iW_i$ with forms $B_i$.  The theorem allows the disconnected
reductive group $J$ and the arbitrary character $\lambda$; by hypothesis each
$B_i$ is invariant under all of $J$, not merely under $L$.  Its contrapositive
yields a single
nonzero $(J,\lambda)$-equivariant distribution supported on the null cone and
adapted simultaneously to every $B_i$.  After forgetting the finite
$J/L$-equivariance, this is a nonzero $L$-invariant distribution of precisely
the type excluded by speciality.  Hence
$\cS^*(Q_L(V))^{L,\lambda}=0$.
\end{proof}

For nice symmetric factors the forms may be taken from the restrictions of the
invariant form on the simple ideals; every inner component preserves them.
Thus Lemma~\ref{lem:special-component} applies to all nice factors.

We need one localization fact in order to tensor nonzero regular-set
distributions without introducing support on the product null cone.

\begin{lemma}[Invariant localization away from the null cone]
\label{lem:localize-away}
Let $T\in\cS^*(R_L(V))^{L,\lambda}$ be nonzero.  There exists a nonzero
$T'\in\cS^*(R_L(V))^{L,\lambda}$ whose support is closed in $Q_L(V)$ and is
disjoint from $\Gamma_L(V)$.
\end{lemma}

\begin{proof}
Compactly supported smooth functions are dense in the Schwartz space of a Nash
manifold \cite{AGSchwartz}; hence choose $f\in C_c^\infty(R_L(V))$ with
$T(f)\ne0$.  Let $\pi:Q_L(V)\to Q_L(V)\sslash L$ be the categorical quotient.
Choose an $A$-stable finite-dimensional generating subspace of
$\C[Q_L(V)]^L$; after subtracting the values at $0$, it gives an
$A$-equivariant closed embedding of $Q_L(V)\sslash L$ into a
finite-dimensional $A$-module with $\pi(0)$ represented by the origin.  The compact set $\pi(\supp f)$ and its finite $A$-orbit are
disjoint from the origin.  Choose a compactly supported smooth function on the
ambient affine space which equals one on that orbit and whose support misses
the origin, and average it over $A$.  Its pullback $\varphi$ is a
$J$-invariant Schwartz multiplier \cite{AGSchwartz}.  Then
$T'=\varphi T$ is nonzero, has the same equivariance, and is supported in
$\pi^{-1}(K)$ for a compact set $K$ not containing the origin.  The multiplier
vanishes on the inverse image of a neighborhood of the origin, so the support
of $T'$ is closed in $Q_L(V)$ and is separated from the null cone by an open
neighborhood.  Extension by zero therefore regards $T'$ as a Schwartz
distribution on $Q_L(V)$ with the same equivariance.
\end{proof}

\begin{proposition}[Products]\label{prop:component-product}
Let $(L_i,A_i,B_i,V_i)$, $1\leq i\leq m$, be componentwise regular.  Put
\[
 L=\prod_iL_i,\qquad A=\prod_iA_i,\qquad B_0=\prod_iB_i,
 \qquad V=\bigoplus_iV_i.
\]
Then $(L,A,B_0,V)$ is componentwise regular.
\end{proposition}

\begin{proof}
One has
\[
 Q_L(V)=\bigoplus_iQ_{L_i}(V_i),
 \qquad
 \Gamma_L(V)=\prod_i\Gamma_{L_i}(V_i).
\]
Let $\lambda=\boxtimes_i\lambda_i$ be nontrivial and trivial on $B_0$, and
suppose $\cS^*(Q_L(V))^{L,\lambda}\ne0$.  The Schwartz kernel theorem and the
density of algebraic tensor products in the completed Schwartz tensor product
\cite{AGSchwartz} allow one to choose a decomposable test function
$f_1\otimes\cdots\otimes f_m$ on which a nonzero distribution $T$ does not
vanish.  For fixed $i$, define the slice
\[
 T_i(\varphi)=T(f_1\otimes\cdots\otimes f_{i-1}\otimes\varphi
                  \otimes f_{i+1}\otimes\cdots\otimes f_m).
\]
It is nonzero because $T_i(f_i)=T(f_1\otimes\cdots\otimes f_m)$.  With the push-forward convention
$(j_*T)(f)=T(f\circ j)$, acting in the $i$th factor gives
$(j_i)_*T_i=\lambda_i(j_iL_i)T_i$; hence $T_i$ is $L_i$-invariant and has
character $\lambda_i$.  Hence
\[
 \cS^*(Q_{L_i}(V_i))^{L_i,\lambda_i}\ne0
\]
for every $i$.  If $\lambda_i$ is nontrivial, componentwise regularity gives a
nonzero distribution on $R_{L_i}(V_i)$; replace it using
Lemma~\ref{lem:localize-away} so that its support is closed in
$Q_{L_i}(V_i)$.  If $\lambda_i$ is trivial, use $\delta_0$ on
$Q_{L_i}(V_i)$.

The external tensor product of these distributions is nonzero, has character
$\lambda$, and has closed support in $Q_L(V)$.  At least one factor lies in a
regular set, so its support is contained in
\[
 R_L(V)=Q_L(V)\setminus\prod_i\Gamma_{L_i}(V_i).
\]
It therefore restricts to a nonzero element of
$\cS^*(R_L(V))^{L,\lambda}$.  This proves the contrapositive of
\eqref{eq:component-implication}.
\end{proof}

\subsection{Intermediate fixed subgroups}

For a subgroup $B\leq A$, the notation $\cS^*(X)^{L,B}$ means invariance under
$L$ and under the inverse image of $B$ in $J$.

\begin{proposition}[Finite-component assembly]\label{prop:finite-assembly}
Assume $(L,A,B_0,V)$ is componentwise regular and let
$B_0\leq B\leq A$.  Let $a\in A$ satisfy $a^2\in B$.  If
\[
 \cS^*(R_L(V))^{L,B}\subset\cS^*(R_L(V))^a,
\]
then
\[
 \cS^*(Q_L(V))^{L,B}\subset\cS^*(Q_L(V))^a.
\]
Consequently every symmetric central quotient represented by the intermediate
fixed-component group $B$ is regular.
\end{proposition}

\begin{proof}
If $a\in B$ there is nothing to prove.  Otherwise put $C=\langle B,a\rangle$.
Since $A$ is abelian and $a^2\in B$, the quotient $C/B$ has order two.  Let
$\chi:C\to\{\pm1\}$ be the character trivial on $B$ and satisfying
$\chi(a)=-1$.  The premise says precisely that
\[
 \cS^*(R_L(V))^{L,C,\chi}=0.
\]
Every character of the subgroup $C$ extends to the finite abelian group $A$.
The $\chi$-isotypic space for $C$ is the direct sum of the $A$-character spaces
for the extensions $\lambda$ of $\chi$.  Each such $\lambda$ is nontrivial and
is trivial on $B_0$.  Componentwise regularity therefore makes all corresponding
spaces on $Q_L(V)$ vanish.  Hence
$\cS^*(Q_L(V))^{L,C,\chi}=0$, which is exactly the asserted $a$-invariance.
For a quotient pair, the square condition $a^2\in B$ is part of admissibility,
and \eqref{eq:preimage-fixed} identifies its full fixed-component group with
such an intermediate $B$.
\end{proof}

\begin{corollary}[Central products]\label{cor:central-products}
Suppose the canonical-cover factors of a connected complex symmetric pair are
pleasant, special, or one of the four factors proved componentwise regular in
Sections~\ref{sec:DIII-family}--\ref{sec:EVII-family}.  Then every finite
$\theta$-stable central quotient, including every diagonal coupling of the
factors and the central torus, is regular.
\end{corollary}

\begin{proof}
Group the simple factors by the orbits of $\theta$.  A two-element orbit gives
a diagonal pair and is pleasant.  For the remaining factors, pleasantness makes
the character condition vacuous, special factors satisfy it by
Lemma~\ref{lem:special-component}, and the four residual factors satisfy it by
the results cited in the statement.  Proposition~\ref{prop:component-product}
handles the direct product on the canonical cover.  The central anti-invariant
summand disappears by Lemma~\ref{lem:fixed-vectors}; a finite central quotient
only replaces $B_0$ by the exact intermediate subgroup $B_F$ described in
Lemma~\ref{lem:exact-component-identification}.  Apply
Proposition~\ref{prop:finite-assembly}.
\end{proof}

\section{The DIII family}\label{sec:DIII-family}

\begin{theorem}\label{thm:DIII-main}
Let $r\ge4$.  Let $G$ be a connected complex semisimple group with Lie algebra
$\mathfrak{so}_{2r}(\C)$, let $\theta$ be an involution of DIII type, and put
$H=G^\theta$.  Then $(G,H,\theta)$ is regular.
\end{theorem}

The proof is carried out in the standard orthogonal model and then read on the
effective adjoint action.  In even rank the only nontrivial component interchanges
$\bigwedge^2V$ and $\bigwedge^2V^\vee$.  We construct two exchanged orbits whose
closures meet in one closed boundary orbit, take the difference of their orbital
measures, and cross the boundary by the automatic-extension theorem of
Chen--Sun.

\subsection{Effective model and component reduction}
\subsubsection{The orthogonal model}

By DIII type we mean that the symmetric Lie algebra pair is isomorphic to
$(\mathfrak{so}_{2r},\mathfrak{gl}_r)$.  Since any automorphism of a complex semisimple Lie algebra 
integrate uniquely to an automorphism of its connected adjoint group, the induced adjoint symmetric 
pair is, up to isomorphism, the standard DIII pair below.  We therefore fix this model for the 
effective-action calculations.

Put $V=\C^r$ and $E=V\oplus V^\vee$, equipped with the hyperbolic symmetric form
\[
  q(v+\lambda,w+\mu)=\lambda(w)+\mu(v).
\]
Relative to dual bases its matrix is
\[
  S=\begin{pmatrix}0&I_r\\ I_r&0\end{pmatrix}.
\]
Set
\[
  G_0=\SO(E,q),\qquad
  s=\begin{pmatrix}iI_r&0\\0&-iI_r\end{pmatrix},\qquad
  \theta_0=\Ad(s).
\]
Since $s^2=-I_E$ is central, $\theta_0$ is an involution.  Its fixed subgroup is
\[
 H_0=G_0^{\theta_0}
 =\left\{\begin{pmatrix}a&0\\0&a^{-t}\end{pmatrix}:a\in\GL(V)\right\}
 \simeq\GL(V).
\]
The $(-1)$-eigenspace in $\mathfrak{so}(E,q)$ is
\begin{equation}\label{diii:eq:pmodel}
 \fp=
 \left\{
   X(B,C)=\begin{pmatrix}0&B\\ C&0\end{pmatrix}:
   B^t=-B,\ C^t=-C
 \right\}
 \simeq\wedgeTwo V\oplus\wedgeTwo V^\vee,
\end{equation}
and the action of $H_0$ is
\begin{equation}\label{diii:eq:Haction}
  a\cdot(B,C)=\bigl(aBa^t,a^{-t}Ca^{-1}\bigr).
\end{equation}
The one-parameter subgroup $tI_V\subset H_0$ acts on the two summands by
$t^2$ and $t^{-2}$, respectively.  Hence
\begin{equation}\label{diii:eq:no-fixed}
  \fp^{H_0}=0,
  \qquad Q_{H_0}(\fp)=\fp.
\end{equation}

\subsubsection{The effective component group}

\begin{lemma}\label{diii:lem:components-SO}
Let $a\in G_0$ and suppose that $\Ad(a)$ commutes with $\theta_0$.  Then precisely one of the following holds:
\begin{enumerate}[label=\textup{(\roman*)}]
  \item $a\in H_0$;
  \item $r$ is even and
  \[
    a\in H_0\tau,
    \qquad
    \tau=\begin{pmatrix}0&I_r\\ I_r&0\end{pmatrix}\in G_0.
  \]
\end{enumerate}
For even $r$, the action of $\tau$ on $\fp$ is
\begin{equation}\label{diii:eq:tauaction}
  \tau(B,C)=(C,B).
\end{equation}
\end{lemma}

\begin{proof}
Commutation of the two inner automorphisms is equivalent to
\[
  asa^{-1}s^{-1}\in Z(G_0)=\{\pm I_E\}.
\]
If $asa^{-1}=s$, then $a$ preserves the $i$- and $-i$-eigenspaces $V$ and
$V^\vee$ of $s$.  The orthogonality relation therefore gives
\[
  a=\begin{pmatrix}g&0\\0&g^{-t}\end{pmatrix}\in H_0.
\]
If $asa^{-1}=-s$, then $a$ exchanges $V$ and $V^\vee$, so
\[
  a=\begin{pmatrix}0&g\\h&0\end{pmatrix}.
\]
The equation $a^tSa=S$ yields $h=g^{-t}$, and hence
\[
  \det a=(-1)^r\det(g)\det(g^{-t})=(-1)^r.
\]
Thus the second case can occur only for even $r$.  When $r$ is even, $\tau\in G_0$ and every such $a$ belongs to $H_0\tau$.  Formula \eqref{diii:eq:tauaction} follows by direct multiplication.
\end{proof}

Let $G_{\mathrm{ad}}$ be the adjoint group of type $D_r$, and let
$q:G_0\to G_{\mathrm{ad}}$ be the central quotient.  Denote by
$\theta_{\mathrm{ad}}$ the induced DIII involution and put
\[
  C=(G_{\mathrm{ad}})^{\theta_{\mathrm{ad}}}.
\]
Lemma~\ref{diii:lem:components-SO} immediately gives the component group of $C$.

\begin{corollary}\label{diii:cor:Ccomponents}
The identity component of $C$ is $q(H_0)$.  If $r$ is odd, then
$C=q(H_0)$.  If $r$ is even, then
\[
  C=q(H_0)\rtimes\langle q(\tau)\rangle,
  \qquad C/C^\circ\simeq\mathbb Z/2\mathbb Z.
\]
\end{corollary}

\begin{proof}
An element $q(a)$ is fixed by $\theta_{\mathrm{ad}}$ exactly when
$\Ad(a)$ commutes with $\theta_0$.  The assertion follows from
Lemma~\ref{diii:lem:components-SO}.  The two cosets remain distinct after quotienting by $Z(G_0)=\{\pm I_E\}$ because $H_0$ is block diagonal whereas $H_0\tau$ is block anti-diagonal.
\end{proof}

We now pass to an arbitrary connected isogeny form.  Let $(G,H,\theta)$ be as in Theorem~\ref{thm:DIII-main}, let
\[
  \rho:G\longrightarrow G_{\mathrm{ad}}
\]
be the adjoint central isogeny, and put $H_{\mathrm{eff}}=\rho(H)$.  The action on $\fp$ factors through $G_{\mathrm{ad}}$.

\begin{lemma}\label{diii:lem:Heff}
One has
\[
  C^\circ\subset H_{\mathrm{eff}}\subset C.
\]
Consequently, if $r$ is odd then $H_{\mathrm{eff}}=C$, whereas if $r$ is even then either $H_{\mathrm{eff}}=C^\circ$ or $H_{\mathrm{eff}}=C$.
\end{lemma}

\begin{proof}
If $h\in H$, then $\theta(h)=h$, so $\rho(h)\in C$; hence
$H_{\mathrm{eff}}\subset C$.  The differential of $\rho$ identifies
$\operatorname{Lie}(H)$ with the fixed Lie algebra of $d\theta_{\mathrm{ad}}$, which is
$\operatorname{Lie}(C)$.  Therefore the connected algebraic subgroup
$\rho(H^\circ)$ has the same Lie algebra as $C^\circ$, and hence
$\rho(H^\circ)=C^\circ$.  The final assertion follows from
Corollary~\ref{diii:cor:Ccomponents}.
\end{proof}

\begin{proposition}\label{diii:prop:effective-reduction}
Theorem~\ref{thm:DIII-main} is automatic if $r$ is odd or if
$H_{\mathrm{eff}}=C$.  In the remaining case, necessarily $r$ is even and
$H_{\mathrm{eff}}=C^\circ$, and it is enough to prove
\begin{equation}\label{diii:eq:sign-goal}
  \cS^*(R_{H_0}(\fp))^{K,\chi}\ne0,
  \qquad
  K:=H_0\rtimes\langle\tau\rangle,
\end{equation}
where $\chi|_{H_0}=1$ and $\chi(\tau)=-1$.
\end{proposition}

\begin{proof}
Let $g\in G$ be admissible.  The first admissibility condition implies
$\rho(g)\in C$.  If $\rho(g)\in H_{\mathrm{eff}}$, choose $h\in H$ with
$\rho(h)=\rho(g)$.  Then $gh^{-1}\in Z(G)$, so
$\Ad(g)|_{\fp}=\Ad(h)|_{\fp}$.  Thus every $H$-invariant distribution on both
$R_H(\fp)$ and $Q_H(\fp)$ is fixed by $\Ad(g)$, and the regularity implication is automatic.  This proves the first assertion by Lemma~\ref{diii:lem:Heff}.

Suppose now that $r$ is even and $H_{\mathrm{eff}}=C^\circ=q(H_0)$.  The actions of $H$ and $H_0$ on $\fp$ have the same effective image $C^\circ$; hence their invariant distribution spaces, invariant polynomial rings, and nullcones coincide.  In particular, by \eqref{diii:eq:no-fixed}, $Q_H(\fp)=\fp$.  If an admissible $g$ satisfies $\rho(g)\notin C^\circ$, write its effective action as $h_0\tau$ with $h_0\in H_0$.  A nonzero
$T\in\cS^*(R_{H_0}(\fp))^{K,\chi}$ is $H$-invariant and satisfies
\[
  (h_0\tau)_*T=-T\ne T.
\]
Hence the antecedent in Definition~\ref{def:regularity} is false for this $g$.  The elements with effective action in $C^\circ$ were already treated above, so \eqref{diii:eq:sign-goal} is sufficient.
\end{proof}

For the rest of the proof we assume that $r\ge4$ is even and construct the distribution in \eqref{diii:eq:sign-goal}.

\subsection{Orbit geometry and a sign distribution away from the boundary}

Choose a decomposition
\[
  V=V_0\oplus P,
  \qquad \dim V_0=r-2,\quad \dim P=2.
\]
Choose nondegenerate alternating matrices $J_0$ and $J_2$, with real entries, such that
\[
  J_0^2=-I_{V_0},\qquad J_2^2=-I_P.
\]
Define
\[
 B_0=C_0=\begin{pmatrix}J_0&0\\0&0\end{pmatrix},
\]
\[
 B_+=\begin{pmatrix}J_0&0\\0&J_2\end{pmatrix},
 \qquad
 C_+=\begin{pmatrix}J_0&0\\0&0\end{pmatrix},
\]
and
\[
  x_0=X(B_0,C_0),\qquad
  x_+=X(B_+,C_+),\qquad
  x_-:=\tau x_+.
\]
Then $\tau x_0=x_0$ and $\tau x_+=x_-$.  Put
\[
  A=\begin{pmatrix}-I_{V_0}&0\\0&0\end{pmatrix}.
\]
For all three points, $BC=CB=A$.  The polynomial
\[
  F(B,C)=\tr(BC)
\]
is $H_0$-invariant, $F(0)=0$, and
\[
  F(x_0)=F(x_+)=F(x_-)=-(r-2)\ne0.
\]
Every point of $\Gamma_{H_0}(\fp)=\pi^{-1}(\pi(0))$ has the same value as $0$ under every invariant polynomial.  Therefore
\begin{equation}\label{diii:eq:pointsR}
  x_0,x_+,x_-\in R_{H_0}(\fp).
\end{equation}
We shall also use
\begin{equation}\label{diii:eq:square}
  X(B,C)^2=\begin{pmatrix}BC&0\\0&CB\end{pmatrix}.
\end{equation}

Consider the $H_0$-equivariant morphism
\[
  \Phi:\fp\to\End(V),\qquad \Phi(B,C)=BC,
\]
where $H_0\simeq\GL(V)$ acts on $\End(V)$ by conjugation.  Let
\[
  L=Z_{H_0}(A)=\GL(V_0)\times\GL(P).
\]

\begin{lemma}\label{diii:lem:fiber}
The fibre over $A$ is the disjoint union of three $L$-orbits:
\[
  \Phi^{-1}(A)=Lx_+\sqcup Lx_-\sqcup Lx_0.
\]
Consequently,
\begin{equation}\label{diii:eq:preimage-conjclass}
  \Phi^{-1}(H_0\cdot A)
  =H_0x_+\sqcup H_0x_-\sqcup H_0x_0.
\end{equation}
\end{lemma}

\begin{proof}
Let $(B,C)\in\Phi^{-1}(A)$.  Since $B$ and $C$ are skew-symmetric,
\[
  CB=(BC)^t=A^t=A.
\]
Thus $B$ and $C$ commute with $A$ and preserve the eigenspace decomposition
$V=V_0\oplus P$; they are block diagonal.

On $V_0$,
\[
  B_{00}C_{00}=C_{00}B_{00}=-I_{V_0}.
\]
Hence $B_{00}$ is nondegenerate and $C_{00}=-B_{00}^{-1}$.  By a
$\GL(V_0)$-congruence we may arrange $B_{00}=J_0$, after which
$C_{00}=J_0$.  On the two-dimensional space $P$, write
\[
  B_{22}=bJ_2,\qquad C_{22}=cJ_2.
\]
The equation $B_{22}C_{22}=0$ gives $bc=0$.  For $g_P\in\GL(P)$,
\[
  b\longmapsto\det(g_P)b,
  \qquad
  c\longmapsto\det(g_P)^{-1}c.
\]
Thus the cases $(b\ne0,c=0)$, $(b=0,c\ne0)$, and $(b,c)=(0,0)$ are exactly the three $L$-orbits displayed above.

For the second assertion, let $y\in\Phi^{-1}(H_0\cdot A)$.  Choose
$h\in H_0$ such that $\Phi(y)=hAh^{-1}$.  Then
$h^{-1}y\in\Phi^{-1}(A)$, so the first assertion places $y$ in exactly one of the three stated $H_0$-orbits.
\end{proof}

\begin{proposition}\label{diii:prop:closures}
The orbit closures in $\fp$ satisfy
\begin{align}
 \overline{H_0x_+}&=H_0x_+\sqcup H_0x_0,\label{diii:eq:clplus}\\
 \overline{H_0x_-}&=H_0x_-\sqcup H_0x_0,\label{diii:eq:clminus}
\end{align}
and $H_0x_0$ is closed.
\end{proposition}

\begin{proof}
The matrix $A$ is semisimple, so its $H_0$-conjugacy class in $\End(V)$ is Zariski closed and hence closed in the ordinary topology of the underlying real Nash manifold.  By equivariance of $\Phi$, every point in
$\overline{H_0x_+}$ therefore belongs to $\Phi^{-1}(H_0\cdot A)$.
Lemma~\ref{diii:lem:fiber} leaves only the three orbits in
\eqref{diii:eq:preimage-conjclass} as possible limits.

Let $g_t=I_{V_0}\oplus tI_P$.  For $t\ne0$, $g_t\in H_0$ and
$g_t\cdot x_+\to x_0$ as $t\to0$.  Hence
$H_0x_0\subset\overline{H_0x_+}$.  On $H_0x_+$ one has
$\rk C=r-2$, while on $H_0x_-$ one has $\rk C=r$.  Since matrix rank cannot increase under specialization, $H_0x_-$ is not contained in
$\overline{H_0x_+}$.  This proves \eqref{diii:eq:clplus}; applying $\tau$ gives \eqref{diii:eq:clminus}.

Finally, let $y\in\overline{H_0x_0}$.  Again $y$ lies in the set
\eqref{diii:eq:preimage-conjclass}.  On $H_0x_0$ both $B$ and $C$ have rank
$r-2$, whereas $H_0x_+$ has $\rk B=r$ and $H_0x_-$ has $\rk C=r$.
Equivalently, the loci defined by $\rk B\le r-2$ and $\rk C\le r-2$ are closed; hence neither latter orbit can occur in the closure of $H_0x_0$.  It follows that $H_0x_0$ is closed.  These arguments determine the Zariski closures.  Since ordinary closure is
contained in Zariski closure and the displayed one-parameter degeneration
realizes the only added orbit, the same formulas hold in the ordinary topology
of the underlying real Nash manifold.  Each algebraic orbit is a locally closed
smooth subvariety and hence a Nash submanifold; in particular, $H_0x_0$ is Nash
closed and $H_0x_\pm$ are closed Nash submanifolds after removing $H_0x_0$.
\end{proof}

Put
\[
  \mathcal B=H_0x_0=Kx_0,
  \qquad
  M=R_{H_0}(\fp),
  \qquad
  \Omega=M\setminus\mathcal B.
\]
Since $\tau$ normalizes $H_0$, it preserves $\Gamma_{H_0}(\fp)$ and hence $M$; therefore $\mathcal B$, $M$, and $\Omega$ are $K$-stable.  By Proposition~\ref{diii:prop:closures}, $H_0x_+$ and $H_0x_-$ are closed Nash submanifolds of $\Omega$.

\begin{lemma}\label{diii:lem:stabilizers}
There are natural isomorphisms
\begin{align}
 (H_0)_{x_+}&\simeq \Sp(V_0,J_0)\times\Sp(P,J_2),\label{diii:eq:stabplus}\\
 (H_0)_{x_0}&\simeq \Sp(V_0,J_0)\times\GL(P).\label{diii:eq:stabzero}
\end{align}
Moreover $K_{x_+}=(H_0)_{x_+}$.
\end{lemma}

\begin{proof}
For $x_+$, the form $B_+$ is nondegenerate and $\ker C_+=P$.  An element preserving both tensors preserves $P$, and hence its $B_+$-orthogonal complement $V_0$.  The restrictions are precisely the two symplectic groups in \eqref{diii:eq:stabplus}.

For $x_0$, write $g\in\GL(V)$ in blocks relative to $V_0\oplus P$.
The equations
\[
  gB_0g^t=B_0,
  \qquad
  g^{-t}C_0g^{-1}=C_0
\]
force both off-diagonal blocks to vanish.  The $V_0$ block lies in
$\Sp(V_0,J_0)$ and the $P$ block is arbitrary, proving \eqref{diii:eq:stabzero}.

Finally, every element of $H_0\tau$ swaps the two ranks.  At $x_+$ they are
$r$ and $r-2$, respectively.  Hence no element of $H_0\tau$ fixes $x_+$.
\end{proof}

Both $H_0$ and $(H_0)_{x_+}$ are complex reductive groups and are therefore unimodular as real Nash groups.  Hence
$H_0x_+\simeq H_0/(H_0)_{x_+}$ carries a nonzero $H_0$-invariant smooth density.  Pushing it forward through the closed embedding
$H_0x_+\hookrightarrow\Omega$ gives an $H_0$-invariant distribution $\mu_+$ on $\Omega$.  Set
\[
  \mu_-:=\tau_*\mu_+,
  \qquad
  T^\circ:=\mu_+-\mu_-.
\]
The two supports are disjoint, so $T^\circ\ne0$.  Since $\tau$ normalizes $H_0$ and $\tau^2=1$,
\begin{equation}\label{diii:eq:Topen-equiv}
  h_*T^\circ=T^\circ\quad(h\in H_0),
  \qquad
  \tau_*T^\circ=-T^\circ.
\end{equation}

\begin{lemma}\label{diii:lem:orbital-schwartz}
The distributions $\mu_+$, $\mu_-$, and $T^\circ$ are Schwartz distributions on $\Omega$.
\end{lemma}

\begin{proof}
The support of $\mu_+$ is the closed semialgebraic orbit $H_0x_+$.  It is Nashly compact modulo $H_0$ in the sense of \cite[Definition~4.0.3]{AG09}: take the compact set $C=\{x_+\}$ and the closed semialgebraic set
$Z=H_0x_+$, so
\[
  \Supp(\mu_+)\subset Z=H_0C.
\]
The compact-modulo-group theorem \cite[Theorem~B.4.1]{AG09} therefore gives
$\mu_+\in\cS^*(\Omega)^{H_0}$.  The same argument applies to $\mu_-$, and hence to their difference.
\end{proof}

We have proved
\begin{equation}\label{diii:eq:Topen}
  0\ne T^\circ\in\cS^*(\Omega)^{K,\chi}.
\end{equation}
It remains to extend $T^\circ$ across the single boundary orbit $\mathcal B$.

\subsection{The boundary normal representation}

Let $z=x_0$.  Since $M$ is open in the underlying real vector space of $\fp$ and $K/H_0$ is finite, the complexification of the real normal space to $Kz$ in $M$ is obtained from the complex normal space to $H_0z$ in $\fp$.

\subsubsection{The holomorphic normal space}

By \eqref{diii:eq:square}, the minimal polynomial of $z$ divides
$t(t^2+1)$, so $z$ is semisimple in $\fg$.  Hence $\operatorname{ad}(z)$ is semisimple and
\begin{equation}\label{diii:eq:normal-centralizer}
  \fp=[\fh_0,z]\oplus\fp^z.
\end{equation}
Thus the holomorphic normal space is naturally
$N_z^{\mathrm{hol}}\simeq\fp^z$.

Write
\[
 B=\begin{pmatrix}B_{00}&M\\-M^t&B_{22}\end{pmatrix},
 \qquad
 C=\begin{pmatrix}C_{00}&N\\-N^t&C_{22}\end{pmatrix}.
\]
A direct calculation of $[X(B,C),z]=0$ gives
\[
  M=N=0,
  \qquad
  C_{00}=-J_0B_{00}J_0,
\]
with $B_{22}$ and $C_{22}$ arbitrary skew-symmetric $2\times2$ matrices.  Therefore
\begin{equation}\label{diii:eq:Nhol}
  N_z^{\mathrm{hol}}
  \simeq
  \wedgeTwo V_0\oplus\wedgeTwo P\oplus\wedgeTwo P^\vee.
\end{equation}
By Lemma~\ref{diii:lem:stabilizers},
\[
  (H_0)_z=\Sp(V_0,J_0)\times\GL(P).
\]

\subsubsection{The component involution}

On the $\wedgeTwo V_0$ summand, $\tau$ sends the parameter $B_{00}$ to
\begin{equation}\label{diii:eq:tauV0}
  B_{00}\longmapsto -J_0B_{00}J_0.
\end{equation}
The matrix $J_0$ belongs to $\Sp(V_0,J_0)$, and its action on
$\wedgeTwo V_0$ is
\[
  B_{00}\longmapsto J_0B_{00}J_0^t=-J_0B_{00}J_0.
\]
Thus $\tau$ acts on this summand exactly as the stabilizer element $J_0$.
On the one-dimensional summands $\wedgeTwo P$ and $\wedgeTwo P^\vee$, the group $\GL(P)$ acts through $\det$ and $\det^{-1}$, respectively, while $\tau$ interchanges the two lines.

Chen--Sun use the complexification of the underlying real normal space.  Hence we put
\begin{equation}\label{diii:eq:Nrealcomplex}
  N_z:=N_{z,\mathbb R}\otimes_{\mathbb R}\C
  \simeq N_z^{\mathrm{hol}}\oplus\overline{N_z^{\mathrm{hol}}}.
\end{equation}
Because $J_0$ has real entries, the comparison between $\tau$ and the stabilizer element $J_0$ holds on both summands in \eqref{diii:eq:Nrealcomplex}.

\begin{proposition}\label{diii:prop:normal-parity}
For every $m\ge0$,
\begin{equation}\label{diii:eq:normal-even}
  \bigl(\Sym^mN_z^*\bigr)^{(H_0)_z}
  \subset
  \bigl(\Sym^mN_z^*\bigr)^{\tau=1}.
\end{equation}
\end{proposition}

\begin{proof}
Write $N_z^*=A\oplus B$, where $A$ is the dual of the $V_0$-part in the two summands of \eqref{diii:eq:Nrealcomplex}, and $B$ is the dual of the $P$-part.  The factor $\Sp(V_0)$ acts only on $A$, whereas $\GL(P)$ acts only on $B$.  Therefore
\begin{equation}\label{diii:eq:sym-factor}
  \Sym(N_z^*)^{\Sp(V_0)\times\GL(P)}
  =\Sym(A)^{\Sp(V_0)}\otimes\Sym(B)^{\GL(P)}.
\end{equation}

On $A$, the action of $\tau$ is the action of the group element
$J_0\in\Sp(V_0)$ on both the holomorphic and antiholomorphic copies.  Hence $\tau$ acts trivially on $\Sym(A)^{\Sp(V_0)}$.

For $B$, choose dual coordinates $u,v,\bar u,\bar v$ with characters
\[
  \det^{-1},\quad\det,\quad\overline{\det}^{\,-1},\quad\overline{\det}
\]
under the underlying real group $\GL(P)=\GL_2(\C)$.  A monomial
$u^av^b\bar u^c\bar v^d$ is invariant precisely when
\[
  a=b,\qquad c=d.
\]
Indeed, the determinant map is surjective onto $\C^\times$, and varying independently the modulus and argument of the determinant forces both equalities.  Thus
\[
  \Sym(B)^{\GL(P)}=\C[uv,\bar u\bar v].
\]
The involution $\tau$ interchanges $u$ with $v$ and $\bar u$ with $\bar v$, so it fixes both generators.  By \eqref{diii:eq:sym-factor}, it acts trivially on the full invariant algebra, and hence on every homogeneous degree.
\end{proof}

\subsection{Automatic extension and proof of the main theorem}

In the case under consideration,
\[
  K_z=(H_0)_z\rtimes\langle\tau\rangle.
\]

\begin{lemma}\label{diii:lem:modular}
The Nash groups $K$ and $K_z$ are unimodular.  Consequently
$\delta_{K/K_z}=1$.
\end{lemma}

\begin{proof}
The identity components are complex reductive groups and hence unimodular.  The modular character is a continuous homomorphism to $\mathbb R_{>0}$; after restricting trivially to the identity component it factors through a finite component group, and is therefore trivial.
\end{proof}

\begin{lemma}\label{diii:lem:CScondition}
For every $m\ge0$, the trivial representation of $K_z$ does not occur as a subquotient of
\begin{equation}\label{diii:eq:CSfiber}
  \Sym^m(N_z^*)\otimes\delta_{K/K_z}\otimes\chi|_{K_z}.
\end{equation}
\end{lemma}

\begin{proof}
The restriction of \eqref{diii:eq:CSfiber} to the reductive real algebraic group
$(H_0)_z$ is a finite-dimensional algebraic representation and is completely reducible.  Since $K_z/(H_0)_z$ is finite, an $(H_0)_z$-equivariant projection onto any $K_z$-submodule can be averaged over the finite component group; hence \eqref{diii:eq:CSfiber} is completely reducible as a $K_z$-representation.  It is therefore enough to exclude a fixed vector.  By Lemma~\ref{diii:lem:modular}, such a vector would be an element
\[
  w\in\bigl(\Sym^mN_z^*\bigr)^{(H_0)_z}
\]
satisfying $\tau w=-w$.  Proposition~\ref{diii:prop:normal-parity} gives instead
$\tau w=w$, hence $w=0$.
\end{proof}

\begin{proposition}\label{diii:prop:automatic-extension}
Restriction induces an isomorphism
\begin{equation}\label{diii:eq:restriction-iso}
  \cS^*(M)^{K,\chi}
  \xrightarrow{\ \sim\ }
  \cS^*(\Omega)^{K,\chi}.
\end{equation}
In particular, $T^\circ$ extends uniquely to a nonzero element of
$\cS^*(M)^{K,\chi}$.
\end{proposition}

\begin{proof}
We apply \cite[Theorem~1.13]{ChenSun} to the almost linear Nash group $K$, the $K$-stable open Nash manifold $M$, its open submanifold
$\Omega=M\setminus Kz$, and the trivial rank-one tempered $K$-vector bundle.
The boundary consists of the single orbit $Kz$.  The sign character $\chi$ has finite image and therefore moderate growth.  Lemmas~\ref{diii:lem:modular} and
\ref{diii:lem:CScondition} give exactly the relative modular and boundary-subquotient hypotheses of Chen--Sun for every symmetric degree.

Theorem~1.13 gives the corresponding isomorphism on Schwartz homology; in degree zero, the invariant-distribution consequence explained immediately after
\cite[Theorem~1.12]{ChenSun} identifies the dual restriction map with
\eqref{diii:eq:restriction-iso}.  Since $\chi^{-1}=\chi$, there is no character-convention change.  Applying \eqref{diii:eq:restriction-iso} to the nonzero distribution
$T^\circ$ from \eqref{diii:eq:Topen} gives the asserted extension.
\end{proof}

\begin{proof}[Proof of Theorem~\ref{thm:DIII-main}]
By Proposition~\ref{diii:prop:automatic-extension}, the sign-equivariant space in
\eqref{diii:eq:sign-goal} is nonzero for every even $r\ge4$.  Proposition~\ref{diii:prop:effective-reduction} then proves regularity in the only nonautomatic case.  The same proposition already proves regularity in odd rank and in every isogeny form for which the image of the fixed subgroup contains both effective components.  This exhausts all connected complex DIII isogeny forms.
\end{proof}

\begin{corollary}\label{diii:cor:standard-forms}
For every $r\ge4$, both the simply connected DIII pair with ambient group
$\Spin_{2r}(\C)$ and the standard orthogonal pair
$(\SO_{2r}(\C),\GL_r(\C))$ are regular.
\end{corollary}

\begin{remark}
The proof uses no classification of distinguished nilpotent DIII orbits.  In the even-rank case the only distribution-theoretic input beyond the definition of regularity is the compact-modulo-group theorem of Aizenbud--Gourevitch and the automatic-extension theorem of Chen--Sun.
\end{remark}


\begin{corollary}[DIII component implication]\label{cor:DIII-component}
For the canonical DIII factor, the finite-component datum of
Definition~\ref{def:componentwise-regular} is componentwise regular.
\end{corollary}

\begin{proof}
In odd rank the obstruction quotient $A/B_0$ is trivial.  In even rank it has
order two and its nontrivial character is the sign character in
\eqref{diii:eq:sign-goal}.  Proposition~\ref{diii:prop:automatic-extension}
constructs a nonzero sign-equivariant distribution on the regular set.  Hence
the premise of the unique nontrivial component implication is false.
\end{proof}


\section{The balanced CII family}\label{sec:CII-family}

Let $U$ and $W$ be complex symplectic spaces of dimension $2r$, put
$E=U\oplus W$, and set
\[
 G=\Sp(E),\qquad H=\Sp(U)\times\Sp(W).
\]
For $s=\Id_U\oplus(-\Id_W)$ and $\theta=\Ad(s)$, the isotropy representation is
\[
 \fp\simeq U\otimes W,
 \qquad \fp^H=0.
\]
The purpose of this section is to prove regularity of
$(\Sp_{4r},\Sp_{2r}\times\Sp_{2r})$ for every $r\ge1$ and to verify all of its
descendants.  The only effective outer component is the block swap.  The proof
reduces its sign sector by homogeneity to distinguished nilpotent orbits, proves
a stable-density theorem in the conormal centralizer, and then annihilates the
Frobenius fiber including its modular character.

\subsection{The block-swap reduction}
Fix a symplectic isomorphism \(j:U\to W\) and define
\[
\tau=
\begin{pmatrix}
0&j^{-1}\\
j&0
\end{pmatrix}.
\]
Then
\[
\tau^2=1,\qquad \tau s=-s\tau,
\]
and \(\tau\) normalizes \(H\).

\begin{lemma}[Reduction of admissible elements]\label{cii:lem:admissible-components}
Let \(g\in G\) be admissible in the sense of
\cite[Definition~7.4.1]{AG09}.  Then
\[
g\in H\cup H\tau.
\]
Moreover, on the space of \(H\)-invariant distributions on \(\mathfrak p\),
every admissible element in \(H\tau\) has the same action as \(\tau\).
Consequently the only possibly nontrivial regularity implication is
equivalent to
\[
\cS^*(R(\mathfrak p))^{K,\chi}=0
\quad\Longrightarrow\quad
\cS^*(\mathfrak p)^{K,\chi}=0,
\]
where
\[
K:=H\rtimes\langle\tau\rangle,
\qquad
\chi|_H=1,\qquad
\chi(\tau)=-1.
\]
\end{lemma}

\begin{proof}
We use only the first condition in the definition of admissibility.
Since \(\Ad(g)\) commutes with \(\theta\),
\[
\Ad(\theta(g))=\Ad(g),
\]
and hence
\[
g^{-1}\theta(g)\in Z(G).
\]
For \(G=\Sp_{4r}(\C)\),
\[
Z(G)=\{\pm\Id_E\},
\]
so
\[
\theta(g)=g
\quad\text{or}\quad
\theta(g)=-g.
\]
This is a necessary consequence of admissibility; we do not claim that
every element satisfying one of these two equations is automatically
admissible.

If \(\theta(g)=g\), then \(g\in H\).  If \(\theta(g)=-g\), then
\(sg=-gs\), so \(g\) interchanges \(U\) and \(W\).  Therefore
\(h:=g\tau^{-1}\) preserves both symplectic summands, hence \(h\in H\), and
\(g=h\tau\in H\tau\).

Finally, if \(T\) is \(H\)-invariant and \(g=h\tau\), then
\[
\Ad(g)T=\Ad(h)\Ad(\tau)T=\Ad(\tau)T.
\]
Thus, for every admissible element whose action is not already trivial on
\(H\)-invariants, the regularity implication is precisely the implication
that the \(\tau\)-anti-invariant part vanishes.  This is equivalent to the
displayed \((K,\chi)\)-equivariant implication.
\end{proof}

Equivalently, the only implication that remains to be proved is
\begin{equation}\label{cii:eq:regularity-implication}
\cS^*(R(\mathfrak p))^H
\subset
\cS^*(R(\mathfrak p))^\tau
\quad\Longrightarrow\quad
\cS^*(\mathfrak p)^H
\subset
\cS^*(\mathfrak p)^\tau.
\end{equation}
When convenient, a distribution that is \(H\)-invariant and
\(\tau\)-odd will be viewed as a \((K,\chi)\)-equivariant distribution,
where
\[
K=H\rtimes\langle\tau\rangle,
\qquad
\chi|_H=1,\quad \chi(\tau)=-1.
\]

\subsection{Distributional reduction and resonance}
\label{cii:sec:distributional}

The symmetric bilinear form
\begin{equation}\label{cii:eq:B}
B(u\otimes w,u'\otimes w')
=
\omega_U(u,u')\,\omega_W(w,w')
\end{equation}
is nondegenerate and \(H\)-invariant.  If \(j:U\to W\) is the symplectic
identification used to define the block swap, then \(\tau\) exchanges the
two symplectic factors; hence \(B\) is also \(\tau\)-invariant.  Moreover,
\begin{equation}\label{cii:eq:nilcone-isotropic}
\cN(\mathfrak p)\subset
Z(B):=\{x\in\mathfrak p:B(x,x)=0\}.
\end{equation}
Indeed \(x\mapsto B(x,x)\) is \(H\)-invariant, and its value at a point
whose orbit closure contains \(0\) must equal its value at \(0\).

\begin{lemma}[Invariant form and orbit conormals]\label{cii:lem:B-centralizer}
Let
\[
  \kappa(X,Y)=\Tr_E(XY),\qquad X,Y\in\mathfrak{sp}(E).
\]
Then \(B\) is a nonzero scalar multiple of
\(\kappa|_{\mathfrak p\times\mathfrak p}\).  Consequently, for every
\(e\in\mathfrak p\),
\begin{equation}\label{cii:eq:orthogonal-centralizer}
  [\mathfrak h,e]^{\perp_B}=\mathfrak p^e.
\end{equation}
After passage to the underlying real vector space, the same statement holds
for \(B_{\R}=\operatorname{Re}B\).
\end{lemma}

\begin{proof}
The trace form \(\kappa\) is a nondegenerate invariant symmetric form on the
simple Lie algebra \(\mathfrak{sp}(E)\).  Its \(\pm1\)-eigenspaces for
\(d\theta\) are orthogonal, so its restrictions to \(\mathfrak h\) and
\(\mathfrak p\) are nondegenerate.  Both \(B\) and
\(\kappa|_{\mathfrak p}\) are nonzero \(H\)-invariant symmetric bilinear
forms on the irreducible \(H\)-module
\(\mathfrak p\simeq U\otimes W\).  Since
\(\dim\Hom_H(\mathfrak p,\mathfrak p^*)=1\), they differ by a nonzero
scalar.

For \(y\in\mathfrak p\), invariance of \(\kappa\) gives
\[
  \kappa([X,e],y)=\kappa(X,[e,y])\qquad(X\in\mathfrak h).
\]
Here \([e,y]\in\mathfrak h\).  Nondegeneracy of \(\kappa|_{\mathfrak h}\)
therefore shows that \(y\) annihilates \([\mathfrak h,e]\) if and only if
\([e,y]=0\), proving \eqref{cii:eq:orthogonal-centralizer}.  Taking real parts
and viewing the complex spaces as real spaces preserves nondegeneracy and the
same annihilator identity.
\end{proof}

\subsubsection{Homogeneity}

The following form of the reduction avoids any need to regard
\(H\rtimes\langle\tau\rangle\) as the group of complex points of a
reductive algebraic group.

\begin{lemma}[Homogeneity reduction]\label{cii:lem:homogeneity}
Assume the left-hand side of \eqref{cii:eq:regularity-implication}.  If the
right-hand side fails, then there exists a nonzero \(H\)-invariant,
\(\tau\)-odd, \(B\)-adapted Schwartz distribution \(\xi\) on
\(\mathfrak p\) such that
\[
\supp\xi\subset\cN(\mathfrak p).
\]
Equivalently, \(\xi\) is \((K,\chi)\)-equivariant.
\end{lemma}

\begin{proof}
Let
\[
\mathcal L
=
\{T\in\cS^*(\mathfrak p)^H:\tau T=-T\}.
\]
The failure of the right-hand side of
\eqref{cii:eq:regularity-implication} implies \(\mathcal L\neq0\).
By the left-hand side, every element of \(\mathcal L\) restricts to zero
on \(R(\mathfrak p)\); therefore
\[
\supp T\subset\cN(\mathfrak p)\subset Z(B)
\qquad(T\in\mathcal L).
\]

The Fourier transform \(\mathcal F_B\) preserves \(H\)-invariance by
\cite[Lemma~5.1.8]{AG09}.  Since \(\tau\) is a \(B\)-isometry,
\(\mathcal F_B\) commutes with \(\tau\), and consequently
\(\mathcal F_B(\mathcal L)\subset\mathcal L\).
Multiplication by the quadratic function \(x\mapsto B(x,x)\) also
preserves \(\mathcal L\).  Thus the nonzero subspace \(\mathcal L\)
satisfies the hypotheses of the Archimedean homogeneity theorem
\cite[Theorem~5.1.7]{AG09}.  It therefore contains a nonzero
\(B\)-adapted distribution.  The support assertion follows from the
definition of \(\mathcal L\).
\end{proof}

\begin{lemma}[Weil character]\label{cii:lem:weil-trivial}
For the form \(B\), the Weil character
\[
\delta_B(t)=\frac{\gamma(B)}{\gamma(tB)}
\]
is identically \(1\) on \(\C^\times\).
\end{lemma}

\begin{proof}
Choose \(a\in\C^\times\) with \(a^2=t\).  Multiplication by \(a\) identifies
the quadratic spaces defined by \(B\) and \(tB\), so their Weil constants
are equal.
\end{proof}

\begin{lemma}[The zero orbit]\label{cii:lem:zero-orbit}
The zero orbit cannot be maximal in the support of a nonzero
$B$-adapted distribution supported on $\cN(\mathfrak p)$.
\end{lemma}

\begin{proof}
If the zero orbit were maximal, the distribution would be supported at the
origin, since $0$ lies in the closure of every nilpotent orbit.  This is
impossible by Lemma~\ref{core:lem:origin-adapted}; the Weil character is
trivial by Lemma~\ref{cii:lem:weil-trivial}.
\end{proof}

\subsubsection{Resonance}

Let \(0\neq e\in\cN(\mathfrak p)\).  By the normal
Jacobson--Morozov theorem for symmetric pairs
\cite[Proposition~4]{KR}, choose a normal
\(\mathfrak{sl}_2\)-triple
\[
[d,e]=2e,\qquad
[d,f]=-2f,\qquad
[e,f]=d,
\qquad
d\in\mathfrak h,\quad e,f\in\mathfrak p.
\]
Let \(D_t\) be the associated cocharacter, so
\(\Ad(D_t)e=t^2e\), and put
\begin{equation}\label{cii:eq:trace}
T(e)=\Tr\!\left(\ad(d)|_{\mathfrak h^e}\right),
\qquad
n=\dim_\C\mathfrak p=4r^2.
\end{equation}
If an irreducible graded \(\mathfrak{sl}_2\)-summand has highest weight
\(\ell\), its surviving normal line has weighted degree
\[
q=\ell+2.
\]
For a real normal tensor let \((A,B')\) be its total holomorphic and
antiholomorphic weighted degrees.

\begin{lemma}[Resonance equation]\label{cii:lem:resonance}
A nonzero Frobenius normal symbol of a \(B\)-adapted distribution must
satisfy one of
\begin{align}
\textnormal{Type I:}\qquad
&A=B'=T(e)-n, \label{cii:eq:typeI}\\
\textnormal{Type II:}\qquad
&(A,B')=(T(e)-n-2,T(e)-n), \label{cii:eq:typeII}
\end{align}
\end{lemma}

\begin{proof}
The element \((D_t,t^2)\in H\times\C^\times\), where the second factor is
the external homothety, fixes \(e\).  Its stabilizer modular character is
\[
|t|_\C^{T(e)}.
\]
A normal line arising from highest weight \(\ell\) contributes
\(t^{-(\ell+2)}\).  Since \(\delta_B=1\) by
Lemma~\ref{cii:lem:weil-trivial}, comparison with the two adapted homogeneity
characters gives \eqref{cii:eq:typeI} and \eqref{cii:eq:typeII}.
\end{proof}

\subsubsection{Frobenius reciprocity and the conormal filtration}

We record the precise convention used below.  Let a unimodular Nash group
\(L\) act transitively on a Nash manifold \(\mathcal O=L\cdot e\), let
\(L_e\) be the stabilizer, and let
\[
  N_e=T_eM/T_e\mathcal O
\]
be the real normal space in an ambient Nash manifold \(M\).  For a character
\(\psi\) of \(L\), Frobenius reciprocity
\cite[Theorem~2.5.7]{AG09}, followed by dualization of the point coefficient,
gives a canonical isomorphism
\begin{equation}\label{cii:eq:exact-frobenius-fiber}
\cS^*\!\left(
  \mathcal O,\Sym^k(CN_{\mathcal O}^{M})
\right)^{L,\psi}
\simeq
\left(
  \Sym^k(N_e)\otimes_{\R}\C
  \otimes\Delta_{L_e}\otimes\psi^{-1}|_{L_e}
\right)^{L_e}.
\end{equation}
Since \(L\) is unimodular, the Frobenius factor on the conormal
coefficient is
\(\Delta_L|_{L_e}\Delta_{L_e}^{-1}=\Delta_{L_e}^{-1}\).  Passing to a
distribution supported at \(e\) dualizes the coefficient space.  Thus
\[
  \Sym^k(CN_{\mathcal O,e}^{M})\otimes
  \Delta_{L_e}^{-1}\otimes\psi|_{L_e}
\]
is replaced by
\[
  \Sym^k(N_e)\otimes\Delta_{L_e}\otimes\psi^{-1}|_{L_e},
\]
which gives \eqref{cii:eq:exact-frobenius-fiber}.

We also need the following standard consequence of the filtration used in
\cite[Theorem~2.5.6]{AG09} (see the cited proof and
\cite[Section~B.2]{AGS08}).

\begin{lemma}[Maximal-stratum symbol]\label{cii:lem:maximal-symbol}
Let a Nash group \(L\) act on a Nash manifold \(M\), and let
\(Z\subset M\) be a finite union of locally closed \(L\)-orbits.  If a
nonzero \((L,\psi)\)-equivariant distribution supported in \(Z\) has
nonzero germ along an orbit \(\mathcal O\) that is open in its support,
then for some \(k\geq0\),
\[
  \cS^*\!\left(
    \mathcal O,\Sym^k(CN_{\mathcal O}^{M})
  \right)^{L,\psi}\neq0.
\]
\end{lemma}

\begin{proof}
Because \(\mathcal O\) is open in the support,
the complement of \(\mathcal O\) in the support is closed in \(M\).  Since an orbit is
locally closed, \(\overline{\mathcal O}\setminus\mathcal O\) is also
closed.  Removing these two invariant closed sets gives an
\(L\)-invariant open neighborhood \(M'\) of \(\mathcal O\) in which
\(\mathcal O\) is closed and the restricted distribution is supported on
\(\mathcal O\).  The standard filtration of distributions supported on the
closed Nash submanifold \(\mathcal O\subset M'\) by transverse order has
associated graded contained in
\[
  \bigoplus_{k\geq0}
  \cS^*\!\left(\mathcal O,
  \Sym^k(CN_{\mathcal O}^{M})\right).
\]
Choose one point $x\in\mathcal O$ and a relatively compact Nash coordinate
neighborhood $U$ of $x$.  Continuity of the restricted distribution against
finitely many Schwartz seminorms gives a finite transverse-order bound, say
$k$, on $U$.  For $y=l\cdot x$, equivariance identifies the restriction on
$lU$ with the translate of the restriction on $U$, multiplied by the scalar
$\psi(l)$.  The same bound $k$ therefore holds near every point of the
transitive orbit.  Hence the distribution belongs globally to the $k$th term
of the transverse-order filtration on distributions supported on
$\mathcal O$.  Choose the least such $k$.  Its image in the $k$th associated
graded quotient is nonzero and, by functoriality of the filtration, is
$(L,\psi)$-equivariant.  This gives a nonzero element of
\[
 \cS^*\!\left(\mathcal O,\Sym^k(CN_{\mathcal O}^{M})\right)^{L,\psi}.
\]
This is precisely the conormal filtration used in
\cite[Theorem~2.5.6]{AG09}; compare also the Nash submersion formalism in
\cite[Appendix~B]{AG09}.
\end{proof}

The groups used below are unimodular:
\[
  K=H\rtimes\langle\tau\rangle,
  \qquad L=H\times\C^\times.
\]
Indeed, \(H\) is complex reductive and hence unimodular, \(K/H\) is finite,
and \(L\) is a product of unimodular groups.

\begin{lemma}[Naturality of the modular character]
\label{cii:lem:modular-naturality}
Let \(J\) be a Lie group and \(\alpha:J\to J\) a continuous
automorphism.  Then
\[
\Delta_J(\alpha(j))=\Delta_J(j)
\qquad(j\in J).
\]
In particular, if \(e\) is fixed by \(\tau\), then
\[
\Delta_{K_e}(\tau h\tau^{-1})=\Delta_{K_e}(h)
\qquad(h\in K_e),
\]
and
\[
\Delta_{K_e}(\tau)=1.
\]
\end{lemma}

\begin{proof}
If \(m\) is a left Haar measure on \(J\), then \(\alpha_*m\) is another
left Haar measure, hence differs from \(m\) by a positive scalar.
Comparing right translations before and after applying \(\alpha\) proves
the first identity.  For the last assertion, \(\tau\) has finite order
and \(\Delta_{K_e}\) is positive-valued, so
\[
\Delta_{K_e}(\tau)^2=\Delta_{K_e}(\tau^2)=1
\]
forces \(\Delta_{K_e}(\tau)=1\).
\end{proof}


\subsubsection{Elimination of Type II}

The second adapted type can be ruled out on every nonzero nilpotent orbit
without using the orbit classification.

Let
\[
L:=H\times\C^\times
\]
act on \(\mathfrak p\) by the adjoint action of \(H\) together with the
external homothety used in adapted homogeneity.  Let \(\eta\) denote the
corresponding Type I or Type II homogeneity character of the
\(\C^\times\)-factor.  For a nilpotent \(e\), the corresponding Frobenius point fiber is,
with the convention of
\eqref{cii:eq:exact-frobenius-fiber},
\begin{equation}\label{cii:eq:typeII-fiber}
\left(
\Sym^k(N_e)\otimes_{\R}\C
\otimes\Delta_{L_e}
\otimes(1,\eta)^{-1}
\right)^{L_e}.
\end{equation}

Let
\[
z=(-\Id_U,\Id_W)\in H.
\]
Since \(\mathfrak p\simeq U\otimes W\), the element \(z\) acts on
\(\mathfrak p\) by \(-1\).  If \(D_t\) is the cocharacter of a normal
\(\mathfrak{sl}_2\)-triple, then
\[
\Ad(D_i)e=-e.
\]
Hence
\[
c:=zD_i\in H_e
\]
and \((c,1)\in L_e\).

\begin{proposition}\label{cii:prop:typeII}
No Type II Frobenius normal symbol can occur for a
\((K,\chi)\)-equivariant adapted distribution.
\end{proposition}

\begin{proof}
A surviving normal line of weighted degree
\[
q=\ell+2
\]
has \(D_i\)-eigenvalue \(i^{-\ell}\).  Multiplication by \(z\) contributes
\(-1\), so \(c=zD_i\) acts on that line by
\[
-i^{-\ell}=i^{-(\ell+2)}=i^{-q}.
\]
Thus on a real normal tensor of weighted bidegree \((A,B')\), the
eigenvalue is
\[
i^{B'-A}.
\]
For Type II, Lemma~\ref{cii:lem:resonance} gives
\[
B'-A=2,
\]
so \(c\) acts by \(-1\).

On the other hand, \(c\) is a finite-order element of \(H_e\).
Consequently
\[
\Delta_{L_e}(c,1)=1
\]
because the modular character is positive-valued, and
\[
(1,\eta)(c,1)=1.
\]
The point fiber \eqref{cii:eq:typeII-fiber} therefore requires the normal tensor itself to be
fixed by \(c\), contradicting the eigenvalue \(-1\).
\end{proof}

\begin{corollary}[Global exclusion of Type II]
\label{cii:cor:no-typeII-distribution}
There is no nonzero Type II $(K,\chi)$-equivariant adapted distribution
supported on $\cN(\mathfrak p)$.
\end{corollary}

\begin{proof}
Assume that such a distribution exists.  Choose an orbit $\mathcal O$ that is
maximal, for the closure order, among the finitely many orbits in its support.
The zero orbit is excluded by Lemma~\ref{cii:lem:zero-orbit}.  The germ along
$\mathcal O$ is nonzero, so Lemma~\ref{cii:lem:maximal-symbol} produces a
nonzero Type II Frobenius normal symbol on $\mathcal O$.  This contradicts
Proposition~\ref{cii:prop:typeII}.
\end{proof}

Henceforth only Type I needs to be considered.

\subsection{Distinguished nilpotent orbits}\label{cii:sec:distinguished}

For Type I adapted distributions, the Fourier transform is proportional to
the original distribution.  Therefore
\[
\supp\xi\subset\cN(\mathfrak p)
\quad\Longrightarrow\quad
\supp\widehat{\xi}\subset\cN(\mathfrak p^*).
\]

Aizenbud works throughout over a local field of characteristic zero
\cite[\S2]{Aiz13}; hence Corollary~4.3.5 of \cite{Aiz13} applies in
particular to \(F=\C\).  It states that
if \(V\) is a finite-dimensional algebraic representation of a reductive group, if
\(\Gamma(V)\) has finitely many group orbits, and if a distribution on
\(Q(V)\) and its Fourier transform are supported in
\(\Gamma(V)\) and \(\Gamma(V^*)\), respectively, then the distinguished
elements are dense in its support.  We apply the theorem with
\[
F=\C,\qquad V=\mathfrak p,\qquad G=H.
\]

In the present isotropy representation one has \(\mathfrak p^H=0\), so
\(Q(\mathfrak p)=\mathfrak p\).  By Kostant--Rallis \cite{KR} the null cone
\(\Gamma(\mathfrak p)\) is the nilpotent cone and contains only finitely
many \(H\)-orbits.

We next compare Aizenbud's distinguished condition with the usual
\(p\)-distinguished condition.  Since
\(\mathfrak g=\mathfrak{sp}(E)\) is semisimple,
\cite[Lemma~5.1.40]{Aiz13} identifies, via an invariant
nondegenerate bilinear form, the conormal fiber
\(CN^{\mathfrak p}_{H\cdot e,e}\) with \(\mathfrak p^e\) and gives
\(Q(\mathfrak p)=\mathfrak p\).  Moreover,
\cite[Lemma~5.2.2]{Aiz13} proves that Aizenbud's notion of
distinguished element coincides with Sekiguchi's standard symmetric-pair
notion.  In our normalization the conormal identification also follows
directly from Lemma~\ref{cii:lem:B-centralizer}:
\[
  N^*_{H\cdot e,e}\simeq\mathfrak p^e.
\]
Since \(B\) identifies the dual null cone with \(\mathcal N(\mathfrak p)\),
Aizenbud's condition is therefore equivalent to
\[
  \mathfrak p^e\subset\mathcal N(\mathfrak p),
\]
i.e. to \(p\)-distinguishedness.

\begin{proposition}\label{cii:prop:dense-distinguished}
Let \(\xi\) be a nonzero Type I adapted distribution satisfying
\[
\supp\xi\subset\cN(\mathfrak p),
\qquad
\supp\widehat{\xi}\subset\cN(\mathfrak p^*).
\]
Then the \(p\)-distinguished elements are dense in \(\supp\xi\).
In particular, among the maximal nilpotent orbits in the support one may
choose a \(p\)-distinguished orbit.
\end{proposition}

\begin{proof}
All hypotheses of Aizenbud's Corollary~4.3.5 are satisfied by the preceding
discussion.  The final assertion follows because the nilpotent cone is a
finite union of \(H\)-orbits.  If every maximal orbit contained in the
support were non-\(p\)-distinguished, their union would be a nonempty
relatively open subset of the support containing no \(p\)-distinguished
points, contradicting density.
\end{proof}

\subsubsection{Balanced CII diagrams and a normal form}

Nilpotent \(H\)-orbits in \(\mathfrak p\) are parametrized by CII
\(ab\)-diagrams.  We only need Bulois' classification of
\(p\)-distinguished diagrams.

\begin{proposition}[Bulois]\label{cii:prop:bulois}
For type CII, a \(p\)-distinguished \(ab\)-diagram has the following form:
\begin{enumerate}[label=(\roman*)]
\item for every even row length, there is either no row or exactly one
pair of rows, one beginning with \(a\) and one beginning with \(b\);
\item all odd-length rows, if any, begin with the same letter.
\end{enumerate}
In type CII, almost \(p\)-distinguished elements are \(p\)-distinguished as well.
\end{proposition}

This is Proposition~3.8 of \cite{Bulois}.

\begin{corollary}\label{cii:cor:stable-distinguished}
Every \(p\)-distinguished orbit in the balanced CII pair
\[
(C_{2r},C_r+C_r)
\]
is invariant under the block swap \(\tau\).  More precisely its diagram is
\[
(2a_1)^+(2a_1)^-\cdots
(2a_s)^+(2a_s)^-,
\qquad
a_1>\cdots>a_s>0,
\qquad
\sum_{i=1}^s a_i=r.
\]
\end{corollary}

\begin{proof}
Because the two symplectic blocks have equal dimension, the total number
of \(a\)'s and \(b\)'s in the diagram is the same.

Every even row contains the same number of \(a\)'s and \(b\)'s.
An odd row beginning with \(a\) contributes \(+1\) to
\(\#a-\#b\), whereas one beginning with \(b\) contributes \(-1\).
By Proposition~\ref{cii:prop:bulois}, all odd rows have the same initial
letter.  Hence the balance condition forces the number of odd rows to be
zero.  The remaining even rows occur in \(a/b\)-pairs, so the diagram is
fixed by exchanging \(a\) and \(b\), i.e.\ by \(\tau\).
\end{proof}

We now construct such a representative explicitly.

\begin{proposition}[Normal form for balanced \(p\)-distinguished orbits]
\label{cii:prop:normal-form}
Let \(\mathcal O\) be a \(p\)-distinguished nilpotent orbit in balanced
CII.  Then \(\mathcal O\) contains a \(\tau\)-fixed element
\[
e=x_A=
\begin{pmatrix}
0&A\\
A&0
\end{pmatrix},
\qquad
A^*=-A,
\]
where \(A\in\mathfrak{sp}(V)\), \(\dim V=2r\), has pairwise distinct even
Jordan block sizes
\[
n_1>\cdots>n_s>0,
\qquad
\sum_i n_i=2r.
\]
\end{proposition}

\begin{proof}
By Corollary~\ref{cii:cor:stable-distinguished}, the signed diagram of
\(\mathcal O\) is
\[
n_1^+n_1^-\cdots n_s^+n_s^-,
\qquad
n_i=2a_i,
\qquad
n_1>\cdots>n_s.
\]
For each \(i\), let \(V_i\) have basis
\[
v_{i,0},\ldots,v_{i,n_i-1}
\]
and alternating form
\[
\langle v_{i,k},v_{i,\ell}\rangle
=
(-1)^k\delta_{k+\ell,n_i-1}.
\]
Because \(n_i\) is even, this form is skew-symmetric and nondegenerate.
Define
\[
Av_{i,k}=v_{i,k+1},
\qquad
Av_{i,n_i-1}=0.
\]
Then
\[
\langle Av,w\rangle+\langle v,Aw\rangle=0,
\]
so \(A^*=-A\).  On
\[
V=\bigoplus_iV_i
\]
with the orthogonal direct-sum symplectic form, \(A\) has precisely the
Jordan blocks \(n_1,\ldots,n_s\).

Identify \(U\) and \(W\) symplectically with \(V\), choose
\[
\tau(u,w)=(w,u),
\]
and define
\[
e=x_A=
\begin{pmatrix}
0&A\\
A&0
\end{pmatrix}.
\]
Since \(A^*=-A\), this element belongs to \(\mathfrak p\), and directly
\[
\tau e\tau^{-1}=e.
\]

For every block \(V_i\), the vectors
\[
(v_{i,0},0),\ (0,v_{i,1}),\ (v_{i,2},0),\ldots
\]
form an \(e\)-Jordan chain of length \(n_i\) beginning in \(U\), whereas
\[
(0,v_{i,0}),\ (v_{i,1},0),\ (0,v_{i,2}),\ldots
\]
form a second chain of the same length beginning in \(W\).  The two chains
are interchanged term-by-term by \(\tau\).  Hence the signed CII diagram
of \(e\) is exactly
\[
n_1^+n_1^-\cdots n_s^+n_s^-.
\]
Nilpotent \(H\)-orbits in CII are classified by their signed
\(ab\)-diagrams, so \(e\in\mathcal O\).
\end{proof}

\subsection{The stable density lemma}

We now prove the central geometric statement.

\begin{theorem}[Stable density lemma]\label{cii:thm:stable-density}
Let \(\mathcal O\) be a \(p\)-distinguished nilpotent orbit in balanced
CII, and let \(e=x_A\in\mathcal O\) be the \(\tau\)-fixed representative
supplied by Proposition~\ref{cii:prop:normal-form}.  Then
\[
\overline{
H_e^\circ\cdot(\mathfrak p^e)^\tau
}^{\,\mathrm{Zar}}
=
\mathfrak p^e.
\]
Moreover, the image of
\[
H_e^\circ\times(\mathfrak p^e)^\tau\longrightarrow\mathfrak p^e,
\qquad
(h,x)\longmapsto h\cdot x,
\]
contains a nonempty Euclidean open subset.
\end{theorem}

\subsubsection{Matrix model}

Identify
\[
U\simeq W\simeq V,
\qquad
\dim V=N=2r.
\]
For \(C\in\End(V)\), write \(C^*\) for the symplectic adjoint.
Then
\[
\mathfrak p
\simeq
\left\{
x_C=
\begin{pmatrix}
0&C\\
-C^*&0
\end{pmatrix}
:C\in\End(V)
\right\},
\]
and the block swap acts by
\[
\tau(C)=-C^*.
\]
Thus every \(C\) decomposes uniquely as
\[
C=R+S,
\qquad
R^*=-R,\quad S^*=S,
\]
with
\[
\tau R=R,\qquad
\tau S=-S.
\]

For this representative,
\[
e=x_A=
\begin{pmatrix}
0&A\\
A&0
\end{pmatrix},
\qquad
A^*=-A,
\]
where the Jordan block sizes of \(A\) are distinct even integers
\[
n_1>n_2>\cdots>n_s,
\qquad
\sum_i n_i=N.
\]

Let
\[
h=
\begin{pmatrix}
X&0\\
0&Y
\end{pmatrix}\in\mathfrak h,
\qquad
P=\frac{X+Y}{2},
\qquad
Q=\frac{X-Y}{2}.
\]
The equation \([h,e]=0\) becomes
\[
[P,A]=0,
\qquad
QA+AQ=0.
\]
Similarly,
\[
(\mathfrak p^e)^\tau
=
\{R\in\mathfrak{sp}(V):[R,A]=0\},
\]
and
\[
(\mathfrak p^e)^{-\tau}
=
\{S\in\End(V):S^*=S,\ SA+AS=0\}.
\]
At a point \(R\in(\mathfrak p^e)^\tau\), the infinitesimal action of
\((\mathfrak h^e)^{-\tau}\) on the odd tangent space is
\[
\Phi_R:
(\mathfrak h^e)^{-\tau}
\longrightarrow
(\mathfrak p^e)^{-\tau},
\qquad
Q\longmapsto QR+RQ.
\]

Thus it suffices to find \(R\) for which \(\Phi_R\) is surjective.

\subsubsection{Construction of a cyclic commuting partner}

Decompose
\[
V=V_1\oplus\cdots\oplus V_s
\]
into the Jordan blocks of \(A\).  Choose bases
\[
v_{i,0},v_{i,1},\ldots,v_{i,n_i-1}
\]
such that
\[
Av_{i,k}=v_{i,k+1},
\qquad
Av_{i,n_i-1}=0,
\]
and
\[
\langle v_{i,k},v_{i,\ell}\rangle
=
(-1)^k\delta_{k+\ell,n_i-1}.
\]
Because \(n_i\) is even, this realizes \(A^*=-A\).

Define
\[
Jv_{i,k}=(-1)^{i-1+k}v_{i,k}.
\]
Then
\[
J^2=1,\qquad
J^*=-J,\qquad
JAJ=-A.
\]

For \(1\leq i<s\), define
\[
F_i:V_i\to V_{i+1},
\qquad
F_i(v_{i,k})=
\begin{cases}
v_{i+1,k},&0\leq k<n_{i+1},\\
0,&k\geq n_{i+1}.
\end{cases}
\]
Then
\[
F_iA=AF_i.
\]
If \(d_i=n_i-n_{i+1}\), direct computation gives
\[
F_i^*v_{i+1,k}=v_{i,k+d_i}.
\]
Define \(R\in\End(V)\) by the nonzero blocks
\[
R_{i+1,i}=F_i,
\qquad
R_{i,i+1}=-F_i^*.
\]
Then
\[
[R,A]=0,\qquad
R^*=-R,\qquad
JRJ=-R.
\]
Therefore
\[
R\in(\mathfrak p^e)^\tau.
\]

\begin{lemma}\label{cii:lem:cyclic}
The commuting pair \((A,R)\) is cyclic: there exists \(v\in V\) such that
\[
\C[A,R]v=V.
\]
Consequently
\[
Z_{\End(V)}(A,R)=\C[A,R],
\qquad
\dim Z_{\End(V)}(A,R)=N.
\]
\end{lemma}

\begin{proof}
Consider \(V/AV\).  Let
\[
\bar v_i=v_{i,0}\pmod{AV}.
\]
The reverse block \(-F_{i-1}^*\) sends \(v_{i,0}\) into \(AV_{i-1}\):
for a $p$-distinguished diagram the row lengths are distinct even integers, so
\(d_{i-1}=n_{i-1}-n_i\ge2\), and explicitly
\[
 F_{i-1}^*v_{i,0}=v_{i-1,d_{i-1}}
   =A^{d_{i-1}}v_{i-1,0}\in AV_{i-1}.
\]
By contrast, \(F_i\) sends \(v_{i,0}\) to \(v_{i+1,0}\).  Hence
\[
\bar R\bar v_i=\bar v_{i+1}\quad(i<s),
\qquad
\bar R\bar v_s=0.
\]
Thus \(\bar v_1\) is cyclic for \(\bar R\) on \(V/AV\).

Put
\[
M=\C[A,R]v_{1,0}.
\]
Then \(M+AV=V\).  The algebra \(\C[A]\) is a local Artinian algebra
isomorphic to \(\C[t]/(t^{n_1})\), and \(A\) belongs to its Jacobson
radical.  Nakayama's lemma applied to the \(\C[A]\)-module \(V/M\)
gives \(V/M=0\).  Hence
\[
\C[A,R]v_{1,0}=V.
\]

Now every endomorphism commuting with both \(A\) and \(R\) is uniquely
determined by the image of the cyclic vector \(v_{1,0}\).  Given any such
image, cyclicity produces a polynomial in \(A,R\) with the same value on
\(v_{1,0}\).  Therefore
\[
Z_{\End(V)}(A,R)=\C[A,R].
\]
Since evaluation on the cyclic vector identifies this algebra with \(V\)
as a vector space, its dimension is \(N\).
\end{proof}

\subsubsection{The centralizer involution}

Put
\[
\mathfrak c=Z_{\End(V)}(A)
\]
and define
\[
\Theta(Z)=-JZ^*J.
\]
Then \(\Theta^2=1\) on \(\mathfrak c\); write
\[
\mathfrak c=\mathfrak c_+\oplus\mathfrak c_-
\]
for its eigenspaces.

Multiplication by \(J\) gives identifications
\[
(\mathfrak h^e)^{-\tau}\xrightarrow{\sim}\mathfrak c_-,
\qquad
Q\longmapsto QJ,
\]
and
\[
(\mathfrak p^e)^{-\tau}\xrightarrow{\sim}\mathfrak c_+,
\qquad
S\longmapsto SJ.
\]
Since \(JR=-RJ\),
\[
(QR+RQ)J=[R,QJ].
\]
Thus the surjectivity of \(\Phi_R\) is equivalent to the surjectivity of
\[
\ad R:\mathfrak c_-\longrightarrow\mathfrak c_+.
\]

\begin{lemma}\label{cii:lem:dimension-difference}
One has
\[
\dim\mathfrak c_- -\dim\mathfrak c_+=N.
\]
Moreover
\[
Z_{\End(V)}(A,R)\subset\mathfrak c_-.
\]
\end{lemma}

\begin{proof}
Decompose
\[
\mathfrak c
=
\bigoplus_{i,j}\Hom_A(V_j,V_i).
\]
For \(i\neq j\), the involution \(\Theta\) exchanges the \((i,j)\) and
\((j,i)\) summands, so their combined trace is zero.

For a diagonal block,
\[
\End_A(V_i)=\C[A_i].
\]
For every \(m\geq0\),
\[
\Theta(A_i^m)=-A_i^m.
\]
Therefore
\[
\Tr(\Theta|\mathfrak c)
=
-\sum_i n_i=-N.
\]
Since
\[
\Tr(\Theta|\mathfrak c)
=
\dim\mathfrak c_+ -\dim\mathfrak c_-,
\]
the first assertion follows.

For the second, note that \(A\) and \(R\) commute and satisfy
\[
A^*=-A,\qquad R^*=-R,\qquad
JAJ=-A,\qquad JRJ=-R.
\]
Hence for every monomial \(A^aR^b\),
\[
\Theta(A^aR^b)=-A^aR^b.
\]
Lemma~\ref{cii:lem:cyclic} then gives
\[
Z_{\End(V)}(A,R)=\C[A,R]\subset\mathfrak c_-.
\]
\end{proof}

\begin{proposition}\label{cii:prop:surjectivity}
The map
\[
\Phi_R:
(\mathfrak h^e)^{-\tau}
\longrightarrow
(\mathfrak p^e)^{-\tau}
\]
is surjective.
\end{proposition}

\begin{proof}
By Lemma~\ref{cii:lem:cyclic},
\[
\ker(\ad R|\mathfrak c_-)
=
Z_{\End(V)}(A,R)
\]
has dimension \(N\).  Hence
\[
\operatorname{rank}(\ad R|\mathfrak c_-)
=
\dim\mathfrak c_- -N.
\]
By Lemma~\ref{cii:lem:dimension-difference}, this is
\[
\dim\mathfrak c_+.
\]
Therefore
\[
\ad R:\mathfrak c_-\twoheadrightarrow\mathfrak c_+,
\]
and the equivalent map \(\Phi_R\) is surjective.
\end{proof}

\begin{proof}[Proof of Theorem~\ref{cii:thm:stable-density}]
Consider
\[
\mu:
H_e^\circ\times(\mathfrak p^e)^\tau
\longrightarrow
\mathfrak p^e,
\qquad
(h,x)\longmapsto h\cdot x.
\]
At \((1,R)\), the tangent directions from the second factor give all of
\((\mathfrak p^e)^\tau\), while
Proposition~\ref{cii:prop:surjectivity} shows that the orbit tangent directions
cover all of \((\mathfrak p^e)^{-\tau}\).  Hence the complex differential
\[
d\mu_{(1,R)}
\]
is surjective.  The complex-analytic submersion theorem therefore implies
that the image of \(\mu\) contains a nonempty Euclidean open subset of
\(\mathfrak p^e\).  In particular \(\mu\) is dominant and
\[
\overline{
H_e^\circ\cdot(\mathfrak p^e)^\tau
}^{\,\mathrm{Zar}}
=
\mathfrak p^e.\qedhere
\]
\end{proof}


\subsection{Regularity}\label{cii:sec:regularity}

The distribution theory is carried out on the underlying real Nash
manifold.  Hence normal tensors include mixed
\(z^\alpha\bar z^\beta\) terms.  We use the real polynomial ring
\[
\C[(\mathfrak p^e)_{\R}].
\]

Let \(e\) be the \(\tau\)-fixed representative from
Proposition~\ref{cii:prop:normal-form}.  Since
\[
K=H\rtimes\langle\tau\rangle
\]
and \(\tau e=e\),
\[
K_e=H_e\rtimes\langle\tau\rangle.
\]
The Frobenius point fiber controlling the orbit term in
\eqref{cii:eq:exact-frobenius-fiber} is
\begin{equation}\label{cii:eq:K-fiber}
\mathcal F_{e,k}
=
\left(
\Sym^k(N_e)\otimes_{\R}\C
\otimes\Delta_{K_e}
\otimes\chi^{-1}
\right)^{K_e},
\end{equation}
because \(K\) is unimodular.

The invariant real bilinear form
\[
  B_{\R}=\operatorname{Re}B
\]
identifies
\[
  N_e\simeq((\mathfrak p^e)_{\R})^*.
\]
This is exactly the real form of the annihilator identity
\eqref{cii:eq:orthogonal-centralizer} from
Lemma~\ref{cii:lem:B-centralizer}.
Therefore
\begin{equation}\label{cii:eq:real-poly}
\Sym^k(N_e)\otimes_{\R}\C
\simeq
\C[(\mathfrak p^e)_{\R}]_k.
\end{equation}

Working with $H_e^\circ$ in the next two lemmas is sufficient for the
contradiction: a $K_e$-invariant Frobenius vector is automatically an
$H_e^\circ$-semi-invariant, the density argument then forces its polynomial
part to be $\tau$-fixed, while the remaining sign character makes
$K_e$-invariance impossible.  No density assertion for the full possibly
disconnected stabilizer $H_e$ is used.

\begin{lemma}[Frobenius character]
\label{cii:lem:frobenius-character}
On \(H_e^\circ\), a vector in the fiber
\eqref{cii:eq:K-fiber} corresponds, via \eqref{cii:eq:real-poly}, to a real polynomial \(f\) satisfying
\[
f(hx)=\nu_e(h)f(x),
\qquad
\nu_e:=\Delta_{K_e}|_{H_e^\circ}.
\]
Moreover
\[
\nu_e(\tau h\tau^{-1})=\nu_e(h)
\qquad(h\in H_e^\circ),
\]
and
\[
\Delta_{K_e}(\tau)=1.
\]
\end{lemma}

\begin{proof}
With the standard action on polynomial functions
\[
(h\cdot f)(x)=f(h^{-1}x),
\]
the \(H_e^\circ\)-invariance of
\[
f\otimes\Delta_{K_e}
\]
in \eqref{cii:eq:K-fiber} is equivalent to
\[
h\cdot f=\Delta_{K_e}(h)^{-1}f,
\]
or, equivalently,
\[
f(hx)=\Delta_{K_e}(h)f(x).
\]
The sign character $\chi$ is trivial on \(H_e^\circ\).  Here $\psi$ in
\eqref{cii:eq:exact-frobenius-fiber} denotes a generic Frobenius
equivariance character, whereas $\chi$ denotes the fixed sign character of
$K/H$ used in the regularity argument.
The conjugation invariance of \(\nu_e\) and the equality
\(\Delta_{K_e}(\tau)=1\) are exactly
Lemma~\ref{cii:lem:modular-naturality}.
\end{proof}

\begin{lemma}[Real-polynomial semi-invariants]
\label{cii:lem:real-semiinvariant}
Let
\[
f\in\C[(\mathfrak p^e)_{\R}]
\]
satisfy
\[
f(hx)=\nu_e(h)f(x)
\qquad(h\in H_e^\circ).
\]
Then
\[
f(\tau x)=f(x)
\qquad
\text{for every }x\in\mathfrak p^e.
\]
\end{lemma}

\begin{proof}
By the proof of Theorem~\ref{cii:thm:stable-density}, the action map
\[
\mu:
H_e^\circ\times(\mathfrak p^e)^\tau
\longrightarrow
\mathfrak p^e
\]
has surjective complex differential at \((1,R)\).  Hence its image contains
a nonempty Euclidean open subset
\(\Omega\subset\mathfrak p^e\).

For \(x\in\Omega\), write \(x=hy\) with
\[
h\in H_e^\circ,\qquad
y\in(\mathfrak p^e)^\tau.
\]
Using Lemma~\ref{cii:lem:frobenius-character},
\[
\begin{aligned}
f(\tau x)
&=f((\tau h\tau^{-1})y)\\
&=\nu_e(\tau h\tau^{-1})f(y)\\
&=\nu_e(h)f(y)\\
&=f(hy)=f(x).
\end{aligned}
\]
Thus the real polynomial \(f\circ\tau-f\) vanishes on the nonempty
Euclidean open set \(\Omega\), and hence vanishes identically.
\end{proof}

\begin{corollary}\label{cii:cor:no-sign-fiber}
For every \(p\)-distinguished nilpotent orbit in balanced CII and every
\(k\geq0\),
\[
\mathcal F_{e,k}=0.
\]
Equivalently, every real Frobenius normal fiber relevant to a Type I
\((K,\chi)\)-equivariant distribution has zero sign-isotypic component.
\end{corollary}

\begin{proof}
Suppose a nonzero element of \(\mathcal F_{e,k}\) exists and let \(f\) be
the corresponding polynomial.  By
Lemma~\ref{cii:lem:real-semiinvariant},
\[
\tau\cdot f=f.
\]
On the one-dimensional modular factor,
\[
\Delta_{K_e}(\tau)=1,
\]
whereas
\[
\chi^{-1}(\tau)=-1.
\]
Hence \(\tau\) acts by \(-1\) on the complete tensor appearing in
\eqref{cii:eq:K-fiber}, contradicting \(K_e\)-invariance.
\end{proof}


\subsubsection{Proof of regularity}

\begin{theorem}\label{cii:thm:regular}
For every \(r\geq1\), the symmetric pair
\[
(C_{2r},C_r+C_r)
\]
is regular.
\end{theorem}

\begin{proof}
Assume the left-hand side of
\eqref{cii:eq:regularity-implication} and suppose that the right-hand side
fails.  By Lemma~\ref{cii:lem:homogeneity}, there is a nonzero \(H\)-invariant,
\(\tau\)-odd, adapted distribution \(\xi\) supported on
\(\cN(\mathfrak p)\).  Equivalently, \(\xi\) is
\((K,\chi)\)-equivariant.

Lemma~\ref{cii:lem:zero-orbit} excludes the zero orbit from being a maximal
support stratum.  Corollary~\ref{cii:cor:no-typeII-distribution} excludes Type II, so
\(\xi\) is of Type I.  Its Fourier transform is proportional to itself,
hence
\[
\supp\widehat{\xi}\subset\cN(\mathfrak p^*).
\]
By Proposition~\ref{cii:prop:dense-distinguished},
the \(p\)-distinguished elements are dense in \(\supp\xi\).
Since \(\supp\xi\) is closed and the nilpotent cone has only finitely
many \(H\)-orbits, \(\supp\xi\) is a finite union of orbit closures.
Let \(\mathcal U\) be the union of the orbits that are maximal, for the
closure order, among the orbits contained in \(\supp\xi\).  The complement
\(\supp\xi\setminus\mathcal U\) is a finite union of closures of
nonmaximal orbits and is therefore closed in \(\supp\xi\); hence
\(\mathcal U\) is a nonempty relatively open subset of \(\supp\xi\).
The density of the \(p\)-distinguished elements implies that
\(\mathcal U\) meets the distinguished locus.  Since
\(p\)-distinguishedness is constant on \(H\)-orbits, one of the maximal
orbits in \(\mathcal U\) is \(p\)-distinguished.  Fix such an orbit
\[
  \mathcal O=H\cdot e.
\]

Choose the \(\tau\)-fixed representative
\(e\in\mathcal O\) supplied by Proposition~\ref{cii:prop:normal-form}.
Because \(\mathcal O\) is open in \(\supp\xi\), the germ of \(\xi\)
along \(\mathcal O\) is nonzero.
Lemma~\ref{cii:lem:maximal-symbol} therefore gives a \(k\geq0\) with
\[
  \cS^*\!\left(
    \mathcal O,
    \Sym^k(CN_{\mathcal O}^{\mathfrak p_{\R}})
  \right)^{K,\chi}\neq0.
\]
The exact Frobenius isomorphism \eqref{cii:eq:exact-frobenius-fiber} then gives
\[
  \mathcal F_{e,k}\neq0.
\]
This contradicts Corollary~\ref{cii:cor:no-sign-fiber}.

Hence no such \(\xi\) exists.  The sign implication holds, and
Lemma~\ref{cii:lem:admissible-components} proves regularity.
\end{proof}


\subsection{Descendants and completion of the proof}\label{cii:sec:descendants}

We also give a direct descendant analysis for balanced CII.  In particular,
every descendant is seen to be regular without invoking the finite-component
assembly of Section~\ref{sec:component-assembly}.

We make the descendant reduction explicit.  In the terminology of
\cite[Definition~7.2.2]{AG09}, a descendant is obtained from a semisimple
element \(x=s(g)=g\theta(g)^{-1}\); in particular \(x\) belongs to
\[
  P=\{g\theta(g)^{-1}:g\in G\}.
\]
Thus it suffices to analyze all semisimple \(x\in P\).  For such an
\(x\),
\[
\theta(x)=x^{-1}.
\]
Write the standard symplectic representation as
\[
E=\bigoplus_\lambda E_\lambda.
\]
Because \(x\) is symplectic, the form pairs
\(E_\lambda\) perfectly with \(E_{\lambda^{-1}}\), and
\(E_{\pm1}\) are nondegenerate symplectic subspaces.  Since
\(\theta=\Ad(s)\) and \(\theta(x)=x^{-1}\),
\[
sE_\lambda=E_{\lambda^{-1}}.
\]

Since $\Sp(E)$ is simply connected, the centralizer $G_x$ is connected by
Steinberg's connected-centralizer theorem.  The eigenspace decomposition is
orthogonal and gives the following actual direct-product decomposition of the
group centralizer, not merely a Lie-algebra decomposition up to isogeny.
Choose one representative from each inverse pair
\(\{\lambda,\lambda^{-1}\}\), \(\lambda\neq\pm1\).  Then
\[
G_x
\simeq
\Sp(E_1)\times\Sp(E_{-1})
\times
\prod_{\{\lambda,\lambda^{-1}\}}
\GL(E_\lambda),
\]
where the \(\GL(E_\lambda)\)-factor acts on
\(E_{\lambda^{-1}}\) through the contragredient representation determined
by the symplectic pairing.

\subsubsection{The factors with \texorpdfstring{\(\lambda=\pm1\)}{lambda = +/-1}}

For \(\varepsilon\in\{1,-1\}\), the space \(E_\varepsilon\) is
\(s\)-stable.  Put
\[
E_\varepsilon^+=E_\varepsilon\cap U,\qquad
E_\varepsilon^-=E_\varepsilon\cap W.
\]
Both are symplectic and
\[
E_\varepsilon=E_\varepsilon^+\oplus E_\varepsilon^-.
\]
If
\[
\dim E_\varepsilon^+=2a,\qquad
\dim E_\varepsilon^-=2b,
\]
the descendant factor is
\[
\left(
\Sp_{2(a+b)}(\C),
\Sp_{2a}(\C)\times\Sp_{2b}(\C)
\right).
\]
If \(a=0\) or \(b=0\), this is a trivial symmetric pair and hence regular.
If \(a,b>0\) and \(a\neq b\), it is pleasant by
\cite[Proposition~4.16]{Rubio}, hence regular.  If \(a=b\), it is a
balanced CII factor
\[
(C_{2a},C_a+C_a),
\]
regular by Theorem~\ref{cii:thm:regular}.

\subsubsection{The factors with \texorpdfstring{\(\lambda\neq\pm1\)}{lambda not equal to +/-1}}

Fix one inverse pair and set \(V=E_\lambda\).  The involution \(s\)
identifies \(V\) with \(E_{\lambda^{-1}}\).  Define
\[
\beta(v,w)=\omega(v,sw).
\]
It is nondegenerate.  Moreover
\[
\beta(w,v)
=
\omega(w,sv)
=
\omega(sw,v)
=
-\omega(v,sw)
=
-\beta(v,w),
\]
so \(\beta\) is alternating and \(\dim V=2a\).
We now compute the restricted involution rather than merely identify its
Lie type.  An element of the centralizer on
\(V\oplus E_{\lambda^{-1}}\) is determined by \(a\in\GL(V)\); its action
on \(E_{\lambda^{-1}}\) is the contragredient action determined by
\(\omega\).  Let \(a^{*_{\beta}}\) denote the adjoint with respect to
\(\beta\).  Using \(sV=E_{\lambda^{-1}}\), the equality
\(\omega(av,bw)=\omega(v,w)\), and the definition
\(\beta(v,w)=\omega(v,sw)\), one obtains
\[
  \theta(a)=(a^{*_{\beta}})^{-1}.
\]
Hence
\[
  \GL(V)^\theta
  =\{a\in\GL(V):a^{*_{\beta}}a=1\}
  =\Sp(V,\beta).
\]
Since \(\beta\) is nondegenerate alternating, \(\dim V=2a\), and the
factor is
\[
  (\GL_{2a}(\C),\Sp_{2a}(\C)).
\]

\begin{lemma}[Effective reduction of the \(\GL/\Sp\) factor]
\label{cii:lem:GLSL}
The regularity problem for
\[
(\GL_{2a}(\C),\Sp_{2a}(\C))
\]
is identical to the regularity problem for
\[
(\SL_{2a}(\C),\Sp_{2a}(\C)).
\]
\end{lemma}

\begin{proof}
Let \(V\) be the defining \(2a\)-dimensional symplectic space and let
\(X\mapsto X^*\) denote symplectic adjoint.  On Lie algebras,
\[
  d\theta(X)=-X^*,
\]
so
\[
  \mathfrak p_{\GL}=\{X:X^*=X\}
  =\C I\oplus\mathfrak p_{\SL},
  \qquad
  \mathfrak p_{\SL}=\{X:X^*=X,\ \Tr X=0\}.
\]
The \(\Sp(V)\)-fixed subspace of \(\End(V)\) is \(\C I\) by Schur's
lemma.  Therefore
\[
  Q(\mathfrak p_{\GL})
  =\mathfrak p_{\GL}/\C I
  \xrightarrow{\sim}\mathfrak p_{\SL}
  =Q(\mathfrak p_{\SL})
\]
\(\Sp(V)\)-equivariantly.  Under this isomorphism the null cones and their
complements agree, because both are defined by the same categorical quotient
of the effective \(\Sp(V)\)-module.

It remains to compare admissibility.  Any \(g\in\GL(V)\) has a central
scalar multiple \(cg\in\SL(V)\), and
\[
  \Ad(cg)=\Ad(g).
\]
Condition~(i) in the definition of an admissible element
\cite[Definition~7.4.1]{AG09} depends only on this adjoint transformation.
For condition~(ii), the action on
\(\mathfrak p_{\GL}=\C I\oplus\mathfrak p_{\SL}\) is the identity on
\(\C I\) and the given adjoint action on \(\mathfrak p_{\SL}\).  Closed
\(\Sp(V)\)-orbits in \(\mathfrak p_{\GL}\) are exactly translates
\(cI+\mathcal O\) with \(\mathcal O\) closed in
\(\mathfrak p_{\SL}\).  Hence the \(H\)-admissibility condition on the
full \(\GL\)-space is equivalent to the one on the effective
\(\SL\)-space.  Thus \(g\) is admissible for the \(\GL/\Sp\) problem if
and only if the same adjoint transformation, represented by \(cg\), is
admissible for the \(\SL/\Sp\) problem.

The two regularity implications are therefore identical under the
identification of effective quotients above.
\end{proof}

By \cite[Proposition~4.14]{Rubio}, the pair
\[
  (\SL_{2a},\Sp_{2a})
\]
is pleasant and hence regular.  Lemma~\ref{cii:lem:GLSL} therefore proves
regularity of each \((\GL_{2a},\Sp_{2a})\) factor.

\subsubsection{Products}

The centralizer decomposition above is an actual direct product: its
\(\GL\)-centres are already included in the \(\GL/\Sp\) factors treated by
Lemma~\ref{cii:lem:GLSL}.  Products of regular symmetric pairs are regular by
\cite[Proposition~7.4.4]{AG09}.  Therefore every descendant of the balanced
CII pair is regular.


\begin{corollary}[Balanced CII component implication]\label{cor:CII-component}
The finite-component datum for the balanced CII factor is componentwise
regular.
\end{corollary}

\begin{proof}
Its obstruction quotient has order two.  The unique nontrivial character is
$\chi$ in Lemma~\ref{cii:lem:admissible-components}, and the required implication
is exactly \eqref{cii:eq:regularity-implication}, proved in
Theorem~\ref{cii:thm:regular}.
\end{proof}

\section{The Spin block family}\label{sec:Spin-family}

Let $U$ and $W$ be nondegenerate quadratic spaces of dimensions $r$ and $s$,
put $E=U\oplus W$, and let $\theta$ be induced on $\Spin(E)$ by the isometry
$1_U\oplus(-1_W)$.  The fixed subgroup is
\[
 H=\Spin(U)\times_{\mu_2}\Spin(W),
\]
and the isotropy representation is
\begin{equation}\label{spin:eq:p-intro}
 \fp\simeq\Hom(W,U).
\end{equation}
Its connected effective image is
\begin{equation}\label{spin:eq:H0-action-intro}
 H_0=\SO(U)\times\SO(W),
 \qquad (a,b)X=aXb^{-1}.
\end{equation}
We use the Clifford-theoretic low-rank conventions
$\Spin_1=\mu_2$ and $\Spin_2\simeq\Gm$; their images are
$\SO_1=1$ and $\SO_2\simeq\Gm$.  The ambient group has dimension at
least five in the theorem below, so these conventions affect only a
fixed-subgroup factor and not the ambient Clifford calculations.
After interchanging the blocks, when $r>s$ put
\begin{equation}\label{spin:eq:Kchi-intro}
 K=S(\Or(U)\times\Or(W)),\qquad
 \chi(a,b)=\det a=\det b.
\end{equation}

\begin{theorem}\label{thm:Spin-residual}
If $r,s>0$, $|r-s|>2$, and $rs$ is even, then
\[
 (\Spin_{r+s},\Spin_r\times_{\mu_2}\Spin_s)
\]
is regular.
\end{theorem}

\begin{corollary}\label{cor:Spin-all}
If $r,s>0$ and $r+s\ge5$, then the Spin block pair above is regular.
\end{corollary}

We first dispose of the unequal odd--odd cases by pleasantness.  The
remaining nonpleasant, nonnice cases reduce to the sign implication for
$(K,\chi)$.  If the smaller
block is odd, Przebinda's theorem annihilates the entire sign sector.  If it is
even, a closed-support sign distribution is constructed on the regular set by
a finite orbit-closure calculation and Chen--Sun automatic extension.
\subsection{Conventions and analytic input}\label{spin:sec:conventions}

All algebraic varieties and algebraic groups in the body of the paper are
defined over $\C$.  When Schwartz functions or Schwartz distributions are
used, the complex variety is regarded as its underlying real Nash manifold.
We write $\cS(X)$ for the Schwartz space and $\cS^*(X)$ for its continuous
dual.  On a finite-dimensional complex vector space this agrees, after a
choice of nonzero translation-invariant density, with the tempered generalized
functions used in \cite{JSZ}.

If a real Nash group $L$ acts on $X$ and $\eta:L\to\C^\times$ is a character,
then $\cS^*(X)^{L,\eta}$ denotes the $\eta$-isotypic subspace for the standard
contragredient action on distributions.  In Chen--Sun's dual formulation the
same space is written as a Hom-space with target $\eta^{-1}$.  All characters
used in the proof are determinant characters and hence are self-inverse, so no
inverse-character ambiguity occurs.

For a complex quadratic space $V$, the bilinear form is denoted by
$(\,\cdot\,,\,\cdot\,)_V$ and $T^*$ denotes the adjoint.  Thus
\[
  (Tv,w)_V=(v,T^*w)_V.
\]
The notation $\Sym(V)$ means the vector space of self-adjoint endomorphisms of
$V$; it is canonically identified with the symmetric square of $V$ after the
quadratic form is used to identify $V$ and $V^*$.

We use the invariant-theoretic notation of Aizenbud--Gourevitch.  For a
finite-dimensional representation $V$ of a reductive group $L$, put
\[
  Q_L(V)=(V/V^L)(\C).
\]
Let $\bar\rho_L:Q_L(V)\to Q_L(V)/\!/L$ be the affine categorical quotient
map, and set
\[
  \Gamma_L(V)=\bar\rho_L^{-1}(\bar\rho_L(0)),
  \qquad
  R_L(V)=Q_L(V)\setminus\Gamma_L(V);
\]
see \cite[Notation~2.3.10]{AG09}.  We suppress the subscript $L$ when the
acting group is clear.  Thus $R(\fp)$ below is the complement of this
invariant-theoretic null cone; it is not defined as the complement of the
ambient Lie-algebra nilpotent cone.

For later comparison with tempered generalized functions, the
translation-invariant density on $\fp=\Hom(W,U)$ may be chosen
$\Or(U)\times\Or(W)$-invariant.  Indeed the complex determinant of
$X\mapsto aXb^{-1}$ is
\[
  (\det a)^s(\det b)^{-r}\in\{\pm1\},
\]
so the underlying real Jacobian has absolute value one.  Hence the
identification between Schwartz distributions defined using functions and
tempered generalized functions defined using densities is equivariant for
all orthogonal groups used below.

We use the following three analytic results.

\begin{theorem}[Przebinda, as stated by Jiang--Sun--Zhu]\label{spin:thm:Przebinda}
Let $F\in\{\R,\C\}$, let $V$ be a finite-dimensional nondegenerate quadratic
space over $F$, and let $\Or(V)$ act diagonally on $V^k$.  If
$k<\dim_F V$, then every $\SO(V)$-invariant tempered generalized function on
$V^k$ is $\Or(V)$-invariant.
\end{theorem}

This is \cite[Proposition~1.5]{JSZ}, which cites
\cite[Theorem~C.7]{Przebinda}.

\begin{theorem}[Aizenbud--Gourevitch]\label{spin:thm:AGcompact}
Let a Nash group $L$ act on a Nash manifold $M$, and let $\mathcal E$ be an
$L$-equivariant Nash bundle.  An $L$-invariant distribution with values in
$\mathcal E$ whose support is Nashly compact modulo $L$ is a Schwartz
distribution.
\end{theorem}

This is \cite[Theorem~B.4.1]{AG09}.  We shall apply it to a one-dimensional
character bundle; this is the precise way in which a finite-order equivariance
character is incorporated into the invariant formulation of the theorem.

The third input is the automatic-extension criterion of Chen--Sun.  We record
only the finite-rank consequence that will be used below.  For a $G$-orbit
$Gz$ in a real Nash manifold $M$, put
\[
  N_z=(T_zM/T_z(Gz))\otimes_{\R}\C
\]
and let $N_z^*$ be its complex dual.  Write
\[
  \delta_{G/G_z}=\delta_G|_{G_z}\,\delta_{G_z}^{-1}.
\]

\begin{theorem}[Chen--Sun]\label{spin:thm:ChenSun}
Let $G$ be an almost linear Nash group, let $M$ be a $G$-Nash manifold, and
let $U_0\subset M$ be a $G$-stable open Nash submanifold such that
$M\setminus U_0$ is a finite union of $G$-orbits.  Let $\eta$ be a character
of moderate growth, meaning that $|\eta|$ is bounded above by a positive Nash
function.  Suppose that for every $z\in M\setminus U_0$ and every
$k\ge0$, the trivial representation of $G_z$ does not occur as a subquotient
of
\begin{equation}\label{spin:eq:CS-general}
  \Sym^k(N_z^*)\otimes\delta_{G/G_z}\otimes\eta|_{G_z}.
\end{equation}
Then restriction induces an isomorphism between $\eta^{-1}$-equivariant
Schwartz distributions on $M$ and on $U_0$.
\end{theorem}

This is the invariant-distribution consequence of
\cite[Theorems~1.12--1.13]{ChenSun}; the definitions of $N_z^*$,
$\delta_{G/G_z}$, and moderate growth above are the conventions of that
source, and the vector bundle is taken to be the trivial line bundle.  The theorem in \cite{ChenSun} is formulated through
Schwartz homology and the extension-by-zero map on Schwartz sections.  The
dual statement is a restriction isomorphism on equivariant Schwartz
distributions, which is the form used here.

\subsection{Pleasantness for odd--odd Spin blocks}\label{spin:sec:indexing}

We first determine the range not already covered by pleasantness or niceness.  Let $\theta$ be the involution on $\Spin(E)$ induced by
$1_U\oplus(-1_W)$.  For $n=\dim E$, write $\omega=e_1\cdots e_n$ for a
Clifford volume element when $n$ is even.  With an orthogonal basis adapted to
$U\oplus W$ one has
\begin{equation}\label{spin:eq:theta-omega}
  \theta(\omega)=(-1)^s\omega.
\end{equation}
Recall Rubio's notation
\[
  A_\theta=\{g\in G:\theta(g)\in gZ(G)\}.
\]
The symmetric pair is pleasant when $A_\theta\subset HZ(G)$; pleasant pairs
are regular \cite[Lemma~4.2]{Rubio}.

\begin{proposition}[odd--odd pleasantness]\label{spin:prop:pleasant}
Assume $r$ and $s$ are odd and $r\ne s$.  Then the Spin block pair
\[
  \bigl(\Spin_{r+s},\Spin_r\times_{\mu_2}\Spin_s\bigr)
\]
is pleasant.
\end{proposition}

\begin{proof}
Let $g\in A_\theta$.  Since $r+s$ is even, the center of $\Spin(E)$ is
$\{\pm1,\pm\omega\}$, with $\pi(\pm\omega)=-I_E$ for the standard covering
$\pi:\Spin(E)\to\SO(E)$.  We first rule out the multipliers $\pm\omega$.
If $\theta(g)=\pm\omega g$, then
\[
  I\pi(g)I^{-1}=-\pi(g).
\]
Relative to $E=U\oplus W$, this forces the diagonal blocks of $\pi(g)$ to
vanish, so $\pi(g)(U)\subset W$ and $\pi(g)(W)\subset U$.  Invertibility
then forces $r=s$, contrary to the hypothesis.  Thus every element of
$A_\theta$ has multiplier $\pm1$.

If $\theta(g)=g$, then $g\in H$.  If $\theta(g)=-g$, the oddness of $s$
gives $\theta(\omega)=-\omega$ by \eqref{spin:eq:theta-omega}; hence
\[
  \theta(\omega g)=(-\omega)(-g)=\omega g.
\]
Thus $\omega g\in H$, while $\omega$ is central.  Therefore
$\Ad(g)=\Ad(\omega g)\in\Ad(H)$.  This proves pleasantness.
\end{proof}

Rubio records the balanced Spin block Lie algebra pairs with $|r-s|\le2$
among the nice pairs, and nice symmetric pairs are regular
\cite[Lemmas~4.21--4.22 and Table~3]{Rubio}; see also
\cite{Sekiguchi,Aiz13}.  Consequently, after the nice cases and the
unequal odd--odd pleasant cases have been disposed of, it remains only to
treat the following range.

\begin{corollary}\label{spin:cor:remaining}
To prove Corollary~\ref{cor:Spin-all}, it remains only to prove regularity for
\[
  |r-s|>2,
  \qquad
  rs\equiv0\pmod2.
\]
\end{corollary}

\subsection{Component reduction and determinant sectors}\label{spin:sec:component}

From now on we work in the range of Corollary~\ref{spin:cor:remaining}.  After
interchanging the blocks, assume
\begin{equation}\label{spin:eq:standing}
  r>s>0,
  \qquad
  d:=r-s>2,
  \qquad
  rs\equiv0\pmod2.
\end{equation}
Let $H_0=\SO(U)\times\SO(W)$ denote the orthogonal image in
\eqref{spin:eq:H0-action-intro}.  Since the action of $H$ factors onto $H_0$, the
two groups have the same invariant distributions and the same invariant
polynomial algebra on $\fp$.

Moreover $\fp^{H_0}=0$.  Indeed, here $r\ge4$; if
$X\in\Hom(W,U)$ is fixed by $H_0$, then every vector in $X(W)$ is fixed by
the standard $\SO(U)$-module, which has no nonzero fixed vector.  Therefore
$X=0$.  Hence
\[
  Q(\fp)=\fp.
\]
Let
\[
  \rho:\fp\longrightarrow \fp/\!/H_0
\]
be the affine categorical quotient map.  We write
\begin{equation}\label{spin:eq:Gamma-R}
  \Gamma(\fp)=\rho^{-1}(\rho(0)),
  \qquad
  R(\fp)=\fp\setminus\Gamma(\fp).
\end{equation}
The group $\widehat K=\Or(U)\times\Or(W)$ introduced below normalizes
$H_0$, and therefore preserves $\Gamma(\fp)$ and $R(\fp)$.  This is exactly
the $Q$, $\Gamma$, and $R$ entering
\cite[Notation~2.3.10 and Definition~7.4.2]{AG09}.

\subsubsection{Central multipliers}

\begin{lemma}\label{spin:lem:central-multipliers}
Suppose $g\in\Spin(E)$ satisfies
\[
  \theta(g)=gz,
  \qquad
  z\in Z(\Spin(E)).
\]
Under \eqref{spin:eq:standing}, one has $z=\pm1$.
\end{lemma}

\begin{proof}
If $r+s$ is odd, then $Z(\Spin(E))=\{\pm1\}$.  Suppose $r+s$ is even.
The remaining central elements are $\pm\omega$, and
$\pi(\pm\omega)=-I_E$.  If $z=\pm\omega$, applying
$\pi:\Spin(E)\to\SO(E)$ to $\theta(g)=gz$ gives
\[
  I\pi(g)I^{-1}=-\pi(g).
\]
As in Proposition~\ref{spin:prop:pleasant}, the diagonal blocks of $\pi(g)$
relative to $E=U\oplus W$ vanish.  Since $\pi(g)$ is invertible, this
forces $r=s$, contradicting \eqref{spin:eq:standing}.  Hence $z=\pm1$.
\end{proof}

Choose nonisotropic vectors $u\in U$ and $w\in W$ and normalize them in the
Clifford realization of $\Spin(E)$.  Then
\begin{equation}\label{spin:eq:g0-minus}
  g_0:=uw\in\Spin(E),
  \qquad
  \theta(g_0)=-g_0.
\end{equation}
Its image under $\pi$ is the product of the reflections in $u$ and $w$; in
particular its restrictions to $U$ and $W$ both have determinant $-1$.  Thus
the group generated on $\fp$ by $H_0$ and $\Ad(g_0)$ is exactly the group
$K$ in \eqref{spin:eq:Kchi-intro}.  If $g_1$ and $g_2$ both satisfy
$\theta(g_i)=-g_i$, then
\[
  \theta(g_1g_2^{-1})=(-g_1)(-g_2)^{-1}=g_1g_2^{-1},
\]
so $g_1g_2^{-1}\in H$.  Hence all multiplier-$-1$ elements induce the same
component action on $H_0$-invariant distributions.

\begin{proposition}[sign reduction]\label{spin:prop:sign-reduction}
If
\begin{equation}\label{spin:eq:SR}
  \cS^*(R(\fp))^{K,\chi}=0
  \quad\Longrightarrow\quad
  \cS^*(\fp)^{K,\chi}=0,
\end{equation}
then the Spin block pair is regular.
\end{proposition}

\begin{proof}
Let $g$ be an admissible element in the definition of regularity.  By
Lemma~\ref{spin:lem:central-multipliers}, its multiplier is $\pm1$.  If the
multiplier is $+1$, then $\theta(g)=g$, so $g\in H$ and the required
invariance is automatic.

Suppose the multiplier is $-1$.  On the space of $H_0$-invariant
distributions, $\Ad(g)$ acts as the nontrivial element of $K/H_0$.  For every
representation $V$ on which $H_0$ acts trivially,
\[
  V=V^K\oplus V^{K,\chi}.
\]
Consequently every $H_0$-invariant distribution is fixed by $\Ad(g)$ if and
only if the $(K,\chi)$-isotypic space vanishes.  Apply this on $R(\fp)$ and on
$Q(\fp)=\fp$.  The defining implication for regularity becomes exactly
\eqref{spin:eq:SR}.
\end{proof}

\begin{remark}
No assertion that the explicit element $g_0$ is admissible is required.
Equation \eqref{spin:eq:SR} is a sufficient statement that simultaneously covers
every admissible multiplier-$-1$ element.
\end{remark}

\subsubsection{Passage to the full orthogonal product}

Set
\begin{equation}\label{spin:eq:Khat}
  \widehat K=\Or(U)\times\Or(W).
\end{equation}
Then $K$ is a normal subgroup of index two.  The character $\chi$ has precisely
two extensions to $\widehat K$:
\begin{equation}\label{spin:eq:det-extensions}
  \det_U(a,b)=\det a,
  \qquad
  \det_W(a,b)=\det b.
\end{equation}

\begin{lemma}[index-two decomposition]\label{spin:lem:index-two}
Let $X$ be any $\widehat K$-stable Nash manifold.  Then
\begin{equation}\label{spin:eq:index-two-decomp}
  \cS^*(X)^{K,\chi}
  =
  \cS^*(X)^{\widehat K,\det_U}
  \oplus
  \cS^*(X)^{\widehat K,\det_W}.
\end{equation}
\end{lemma}

\begin{proof}
Choose a reflection $r_U\in\Or(U)$ and put
$\tau=(r_U,I_W)\in\widehat K\setminus K$.  Conjugation by $\tau$ preserves
$K$ and $\chi$, so $\tau$ acts on $\cS^*(X)^{K,\chi}$.  Since $\tau^2=1$,
this space is the direct sum of its $\pm1$ eigenspaces.  On $\tau$ one has
\[
  \det_U(\tau)=-1,
  \qquad
  \det_W(\tau)=1.
\]
The $-1$ eigenspace is therefore the $\det_U$ extension and the $+1$
eigenspace is the $\det_W$ extension.  This proves
\eqref{spin:eq:index-two-decomp}.
\end{proof}

\subsection{Odd smaller block: global vanishing}\label{spin:sec:odd}

Assume throughout this section that $s$ is odd.  The standing parity condition
$rs\equiv0\pmod2$ then implies that $r$ is even.

\begin{proposition}\label{spin:prop:detU-vanish}
For every $r>s>0$ one has
\[
  \cS^*(\fp)^{\widehat K,\det_U}=0.
\]
\end{proposition}

\begin{proof}
Choose a basis of $W$.  As an $\Or(U)$-representation,
\[
  \fp=\Hom(W,U)\simeq U^{\oplus s},
\]
with diagonal action.  Let
$T\in\cS^*(\fp)^{\widehat K,\det_U}$.  Since $\det_U$ is trivial on
$\SO(U)$, the distribution $T$ is $\SO(U)$-invariant.  It is a tempered
generalized function on $U^s$.  Because $s<r$, Theorem~\ref{spin:thm:Przebinda}
forces $T$ to be invariant under all of $\Or(U)$.  A reflection in $\Or(U)$,
however, acts on $T$ by $-1$ by the assumed $\det_U$-equivariance.  Hence
$T=0$.
\end{proof}

\begin{lemma}\label{spin:lem:detW-odd}
If $s$ is odd and $r$ is even, then
\[
  \cS^*(\fp)^{\widehat K,\det_W}=0.
\]
\end{lemma}

\begin{proof}
The central element
\[
  c=(-I_U,-I_W)\in\widehat K
\]
acts trivially on $\fp$, because
\[
  (-I_U)X(-I_W)^{-1}=X.
\]
Thus it acts trivially on every distribution on $\fp$.  On the other hand,
\[
  \det_W(c)=(-1)^s=-1.
\]
A $\det_W$-equivariant distribution would therefore be equal to its negative.
\end{proof}

\begin{corollary}\label{spin:cor:odd-sign-zero}
Under \eqref{spin:eq:standing}, if $s$ is odd, then
\[
  \cS^*(\fp)^{K,\chi}=0.
\]
In particular, the implication \eqref{spin:eq:SR} holds.
\end{corollary}

\begin{proof}
Combine Lemma~\ref{spin:lem:index-two}, Proposition~\ref{spin:prop:detU-vanish}, and
Lemma~\ref{spin:lem:detW-odd}.
\end{proof}

The odd smaller-block half of Theorem~\ref{thm:Spin-residual} is therefore complete.
No nilpotent-orbit analysis and no automatic extension are needed in this
parity.

\subsection{The two-dimensional self-adjoint block}\label{spin:sec:ranktwo}

For the remainder of the proof assume
\begin{equation}\label{spin:eq:s-even}
  s=2m.
\end{equation}
We shall construct a nonzero element of
$\cS^*(R(\fp))^{\widehat K,\det_W}$.

Let $P$ be a complex hyperbolic plane.  Choose a basis $e,f$ with
\begin{equation}\label{spin:eq:hyperbolic}
  (e,e)=(f,f)=0,
  \qquad
  (e,f)=1.
\end{equation}
Define
\begin{equation}\label{spin:eq:Nranktwo}
  Ne=0,
  \qquad
  Nf=e.
\end{equation}
Then
\[
  N=N^*,
  \qquad
  N^2=0,
  \qquad
  N\ne0.
\]

\begin{lemma}\label{spin:lem:ranktwo-centralizer}
The centralizer of $N$ in $\Or(P)$ is
\[
  C_{\Or(P)}(N)=\{\pm I_P\}.
\]
In particular, every element of this centralizer has determinant $+1$.
\end{lemma}

\begin{proof}
Since $N$ is a single Jordan block,
\[
  C_{\End(P)}(N)=\C[N]=\{aI+bN:a,b\in\C\}.
\]
If $g=aI+bN$ is orthogonal, then $N^*=N$ gives
\[
  I=g^*g=g^2=a^2I+2abN.
\]
Thus $a^2=1$ and $2ab=0$.  Since $a\ne0$, one has $b=0$ and
$g=\pm I_P$.  Both determinants equal $1$ because $\dim P=2$.
\end{proof}

\begin{lemma}\label{spin:lem:ranktwo-orbit}
The nonzero nilpotent elements of $\Sym(P)$ form a single $\Or(P)$-orbit, and
its closure is that orbit together with $0$.
\end{lemma}

\begin{proof}
Relative to the basis in \eqref{spin:eq:hyperbolic}, a self-adjoint endomorphism is
of the form
\[
  A(a,b,c)=
  \begin{pmatrix}
    a&b\\ c&a
  \end{pmatrix}.
\]
If $A$ is nilpotent, then $a=0$ and $bc=0$.  Hence the nonzero nilpotents form
the two punctured coordinate lines $b\ne0,c=0$ and $b=0,c\ne0$.  The
reflection exchanging $e$ and $f$ interchanges the two lines.  The torus
\[
  t_z=\diag(z,z^{-1})\in\SO(P)
\]
acts by $b\mapsto z^2b$ and $c\mapsto z^{-2}c$.  Thus all nonzero points form
one orbit, and scaling $z$ to $0$ or $\infty$ shows that $0$ is its unique
boundary orbit.
\end{proof}

\subsection{An explicit regular orbit for an even smaller block}\label{spin:sec:even-orbit}

Choose an orthogonal decomposition
\begin{equation}\label{spin:eq:Wdecomp}
  W=P_1\perp\cdots\perp P_m
\end{equation}
into hyperbolic planes.  On each $P_j$ choose the nilpotent $N_j$ of
Section~\ref{spin:sec:ranktwo}.  Fix pairwise distinct nonzero scalars
\begin{equation}\label{spin:eq:lambdas}
  \lambda_1,\ldots,\lambda_m\in\C^\times
\end{equation}
and set
\begin{equation}\label{spin:eq:Sdef}
  S=\bigoplus_{j=1}^m(\lambda_jI_{P_j}+N_j)\in\Sym(W).
\end{equation}
The operator $S$ is invertible.

Choose $\alpha_j\in\C^\times$ with $\alpha_j^2=\lambda_j$ and define
\begin{equation}\label{spin:eq:Aj}
  A_j=\alpha_j\left(I_{P_j}+\frac{1}{2\lambda_j}N_j\right).
\end{equation}
Because $N_j^2=0$ and $N_j=N_j^*$,
\[
  A_j^*=A_j,
  \qquad
  A_j^2=\lambda_jI_{P_j}+N_j.
\]
Let $A=\bigoplus_jA_j$.  Since $r>s$, there exists an isometric embedding
\[
  \iota:W\hookrightarrow U.
\]
Set
\begin{equation}\label{spin:eq:Xdef}
  X=\iota A\in\fp.
\end{equation}
Then
\begin{equation}\label{spin:eq:Xgram}
  X^*X=A^*A=S.
\end{equation}
In particular $X$ is injective.

We shall repeatedly use the Gram map
\begin{equation}\label{spin:eq:gram-map}
  q:\fp\longrightarrow\Sym(W),
  \qquad
  q(Y)=Y^*Y.
\end{equation}
The following invariant is sufficient to keep the entire orbit closure away
from the invariant-theoretic null cone; no identification with the ambient
Lie-algebra nilpotent cone is needed.

\begin{lemma}\label{spin:lem:det-detects-R}
The polynomial
\[
  D(Y):=\det(Y^*Y)
\]
is $H_0$-invariant and satisfies $D(0)=0$.  Consequently
\[
  D(Y)\ne0\quad\Longrightarrow\quad Y\in R(\fp).
\]
\end{lemma}

\begin{proof}
For $(a,b)\in H_0$,
\[
  q(aYb^{-1})=bq(Y)b^{-1},
\]
so $D$ is invariant.  If $Y\in\Gamma(\fp)=\rho^{-1}(\rho(0))$, then every
$H_0$-invariant polynomial has the same value at $Y$ and $0$ by the
definition of the categorical quotient.  In particular $D(Y)=D(0)=0$.
The contrapositive proves the assertion.
\end{proof}

Since
\[
  D(X)=\det S=\prod_{j=1}^m\lambda_j^2\ne0,
\]
Lemma~\ref{spin:lem:det-detects-R} gives $X\in R(\fp)$.

Let
\[
  \cO=\widehat K\cdot X.
\]
Put $L=X(W)$.  Since $X^*X$ is invertible, $L$ is nondegenerate and
\[
  U=L\perp L^\perp,
  \qquad
  \dim L^\perp=d=r-s.
\]

\begin{lemma}\label{spin:lem:start-stabilizer}
There is a natural isomorphism
\[
  \widehat K_X
  \simeq
  C_{\Or(W)}(S)\times\Or(L^\perp)
  \simeq
  \mu_2^m\times\Or_d(\C).
\]
Moreover,
\[
  \det_W|_{\widehat K_X}=1.
\]
\end{lemma}

\begin{proof}
If $(a,b)\in\widehat K_X$, then $aX=Xb$.  Taking Gram operators gives
$b^{-1}Sb=S$, so $b\in C_{\Or(W)}(S)$.  Conversely, if $b$ centralizes $S$,
the map
\[
  Xw\longmapsto Xbw
\]
is an isometry of $L$.  Extending it orthogonally and choosing an arbitrary
$c\in\Or(L^\perp)$ gives an element of the stabilizer.  This proves the first
isomorphism.

The eigenvalues $\lambda_j$ are pairwise distinct, so every orthogonal
centralizer of $S$ preserves each primary space $P_j$.  On $P_j$ it centralizes
$N_j$.  Lemma~\ref{spin:lem:ranktwo-centralizer} therefore gives
\[
  C_{\Or(W)}(S)
  =\prod_{j=1}^m\{\pm I_{P_j}\}
  \simeq\mu_2^m.
\]
Every $\pm I_{P_j}$ has determinant $+1$, proving the final assertion.
\end{proof}

Because $\widehat K$ and $\widehat K_X$ are reductive real Nash groups, both
are unimodular.  Lemma~\ref{spin:lem:start-stabilizer} therefore yields a nonzero
$\det_W$-equivariant distribution on the homogeneous space
$\cO\simeq\widehat K/\widehat K_X$.  For example, with an invariant quotient
density one may write
\begin{equation}\label{spin:eq:orbital-integral}
  T_{\cO}(f)
  =\int_{\widehat K/\widehat K_X}
       f(kX)\det_W(k)^{-1}\,d\dot k
\end{equation}
for compactly supported smooth test functions on $\cO$.  The integrand is
well-defined because $\det_W$ is trivial on the stabilizer.

The orbit $\cO$ will not be closed in $R(\fp)$, so formula
\eqref{spin:eq:orbital-integral} is not yet the desired global Schwartz
distribution.  We next compute its boundary exactly.

\subsection{The finite orbit closure}\label{spin:sec:closure}

For a subset $J\subset\{1,\ldots,m\}$ define
\begin{equation}\label{spin:eq:SJ}
  N_j^{(J)}=
  \begin{cases}
    0,&j\in J,\\
    N_j,&j\notin J,
  \end{cases}
  \qquad
  S_J=\bigoplus_{j=1}^m
  \bigl(\lambda_jI_{P_j}+N_j^{(J)}\bigr).
\end{equation}
Every $S_J$ is invertible.

\begin{proposition}\label{spin:prop:base-closure}
One has the disjoint orbit decomposition
\begin{equation}\label{spin:eq:base-closure}
  \overline{\Or(W)\cdot S}
  =
  \coprod_{J\subset\{1,\ldots,m\}}
  \Or(W)\cdot S_J.
\end{equation}
More generally,
\[
  \overline{\Or(W)\cdot S_J}
  =
  \coprod_{J\subset J'}\Or(W)\cdot S_{J'}.
\]
Here and below, closures are taken in the ordinary topology of the
underlying real Nash manifold.  No comparison with algebraic orbit closures
is used.
\end{proposition}

\begin{proof}
Let $T$ be a limit of conjugates of $S$.  Orthogonal conjugation preserves
self-adjointness and $\Sym(W)$ is closed, so $T=T^*$.  The coefficients of
the characteristic polynomial are conjugation-invariant and continuous, hence
\begin{equation}\label{spin:eq:charpoly}
  \det(tI-T)=\prod_{j=1}^m(t-\lambda_j)^2.
\end{equation}
For each $j$ put
\[
  W_j(T)=\ker(T-\lambda_jI)^2.
\]
The primary decomposition gives
$W=\bigoplus_jW_j(T)$ and $\dim W_j(T)=2$.

We claim that the primary spaces are mutually orthogonal.  If $i\ne j$, the
polynomials
\[
  p_i(t)=(t-\lambda_i)^2,
  \qquad
  p_j(t)=(t-\lambda_j)^2
\]
are coprime.  Choose polynomials $a,b$ with $ap_i+bp_j=1$.  If
$v\in W_i(T)$ and $w\in W_j(T)$, then, using $T=T^*$, polynomial expressions
in $T$ may be moved from one side of the quadratic pairing to the other.  The
Bezout identity and the relations $p_i(T)v=0$, $p_j(T)w=0$ then give
$(v,w)=0$.  Since the direct sum is the whole nondegenerate space $W$, each
$W_j(T)$ is itself nondegenerate.

On $W_j(T)$ write
\[
  T=\lambda_jI+M_j.
\]
By Cayley--Hamilton, $M_j^2=0$, and $M_j=M_j^*$.  Lemma
\ref{spin:lem:ranktwo-orbit} says that $M_j$ is either $0$ or belongs to the unique
nonzero nilpotent $\Or(W_j(T))$-orbit.  Labeled orthogonal decompositions into
nondegenerate two-planes are conjugate under $\Or(W)$.  Thus $T$ is conjugate
to some $S_J$.  The subset $J$ is determined by whether
$T|_{W_j(T)}-\lambda_jI$ has rank zero or one, so the union in
\eqref{spin:eq:base-closure} is disjoint.

Conversely, every $S_J$ is a limit of conjugates of $S$.  Indeed, on $P_j$ the
torus $t_z=\diag(z,z^{-1})$ satisfies
\[
  t_zN_jt_z^{-1}=z^2N_j.
\]
Letting $z\to0$ independently on the factors indexed by $J$ replaces precisely
those $N_j$ by zero.  The same argument beginning with $S_J$ gives the more
general closure formula.
\end{proof}

The next lemma lifts the base calculation to $\fp$.

\begin{lemma}[invertible Gram fibers]\label{spin:lem:gram-fiber}
If $Y_1,Y_2\in\fp$ satisfy
\[
  Y_1^*Y_1=Y_2^*Y_2=T
\]
with $T$ invertible, then there exists $a\in\Or(U)$ such that
\[
  aY_1=Y_2.
\]
\end{lemma}

\begin{proof}
Both maps are injective and their images are nondegenerate.  The map
$Y_1w\mapsto Y_2w$ preserves the quadratic form because the two Gram
operators coincide.  Witt's extension theorem extends it to an orthogonal
transformation of $U$.
\end{proof}

For each $J$, choose $X_J\in\fp$ with
\[
  X_J^*X_J=S_J.
\]
One may take the canonical choice obtained from \eqref{spin:eq:Aj} by deleting the
nilpotent term on the factors in $J$.

\begin{proposition}\label{spin:prop:p-closure}
One has
\begin{equation}\label{spin:eq:p-closure}
  \overline{\widehat K\cdot X}
  =
  \coprod_{J\subset\{1,\ldots,m\}}
  \widehat K\cdot X_J.
\end{equation}
More generally,
\[
  \overline{\widehat K\cdot X_J}
  =
  \coprod_{J\subset J'}\widehat K\cdot X_{J'}.
\]
The entire closure in \eqref{spin:eq:p-closure} is contained in $R(\fp)$.
\end{proposition}

\begin{proof}
The Gram map is continuous and $\widehat K$-equivariant in the sense
\[
  q(aYb^{-1})=bq(Y)b^{-1}.
\]
Thus every limit point of $\widehat KX$ has Gram operator in the base closure
of Proposition~\ref{spin:prop:base-closure}.  Since all $S_J$ are invertible,
Lemma~\ref{spin:lem:gram-fiber} shows that the inverse image of each base orbit is a
single $\widehat K$-orbit.  Hence the left side of \eqref{spin:eq:p-closure} is
contained in the right side.

For the reverse inclusion we give an explicit lift of the degeneration.  Put
\[
  A_j(t)=\alpha_j
  \left(I+\frac{t}{2\lambda_j}N_j\right),
  \qquad
  A_j(t)^2=\lambda_jI+tN_j.
\]
Let $t_j\to0$ for $j\in J$ and $t_j=1$ for $j\notin J$, and define
\[
  X(t)=\iota\bigoplus_jA_j(t_j).
\]
Then $X(t)\to X_J$.  When every $t_j\ne0$, choose $z_j$ with $z_j^2=t_j$;
the product of the torus elements $t_{z_j}$ conjugates $S$ to $q(X(t))$.
After applying that right orthogonal transformation, the Gram operator agrees
with $S$, and Lemma~\ref{spin:lem:gram-fiber} supplies the required left
orthogonal transformation.  Thus $X(t)\in\widehat KX$ for all nonzero
parameters, proving $X_J\in\overline{\widehat KX}$.  The same proof gives the
general closure formula.

Finally,
\[
  \det S_J=\prod_{j=1}^m\lambda_j^2\ne0
\]
for every $J$.  Since the determinant of the Gram operator is constant on
$\widehat K$-orbits, Lemma~\ref{spin:lem:det-detects-R} shows that every orbit in
the closure lies in $R(\fp)$.
\end{proof}

Put
\begin{equation}\label{spin:eq:Bboundary}
  B=\overline{\cO}\setminus\cO
   =\coprod_{\varnothing\ne J\subset\{1,\ldots,m\}}
     \widehat K\cdot X_J.
\end{equation}
The refined closure formula in Proposition~\ref{spin:prop:p-closure} shows
more strongly that the closure of every orbit indexed by
$\varnothing\ne J$ is a union of orbits still indexed by nonempty subsets.
Hence $B$ is closed in $\fp$ (and therefore also in $R(\fp)$).  Each orbit of the real algebraic group
$\widehat K$ is locally closed and semialgebraic, so the finite union $B$ is
semialgebraic.  Hence
\begin{equation}\label{spin:eq:Vopen}
  V=R(\fp)\setminus B
\end{equation}
is a $\widehat K$-stable open Nash submanifold, and $\cO$ is closed in $V$.
The ordinary orbital distribution \eqref{spin:eq:orbital-integral}, pushed forward
to $V$, has support $\cO$, which is Nashly compact modulo $\widehat K$.
Equivalently, one may regard the determinant character as a one-dimensional
$\widehat K$-equivariant Nash bundle and the distribution as invariant with
values in that bundle.  Theorem~\ref{spin:thm:AGcompact} therefore gives
\begin{equation}\label{spin:eq:TO-Schwartz}
  0\ne T_{\cO}\in\cS^*(V)^{\widehat K,\det_W}.
\end{equation}

\subsection{Boundary stabilizers and conormal representations}\label{spin:sec:conormal}

Fix a nonempty subset $J\subset\{1,\ldots,m\}$ and put $z=X_J$.  Since
$S_J=z^*z$ is invertible, the same argument as in
Lemma~\ref{spin:lem:start-stabilizer} gives
\begin{equation}\label{spin:eq:boundary-stabilizer}
  \widehat K_z
  \simeq
  C_{\Or(W)}(S_J)\times\Or_d(\C).
\end{equation}
The distinct eigenvalues separate the two-dimensional blocks, and
Lemma~\ref{spin:lem:ranktwo-centralizer} gives
\begin{equation}\label{spin:eq:boundary-centralizer}
  C_{\Or(W)}(S_J)
  =
  \prod_{j\in J}\Or(P_j)
  \times
  \prod_{j\notin J}\{\pm I_{P_j}\}.
\end{equation}
Consequently
\begin{equation}\label{spin:eq:boundary-stabilizer2}
  \widehat K_z
  \simeq
  \left(\prod_{j\in J}\Or(P_j)\right)
  \times\mu_2^{m-|J|}
  \times\Or_d(\C).
\end{equation}
Both $\widehat K$ and $\widehat K_z$ are reductive real Nash groups and hence
unimodular.  Thus
\begin{equation}\label{spin:eq:modular-trivial}
  \delta_{\widehat K/\widehat K_z}=1.
\end{equation}
Moreover,
\begin{equation}\label{spin:eq:detW-boundary}
  \det_W|_{\widehat K_z}
  =\prod_{j\in J}\det_{P_j}.
\end{equation}

We next compute the full conormal representation required by
Theorem~\ref{spin:thm:ChenSun}.

\begin{lemma}\label{spin:lem:dq}
Let $z\in\fp$ satisfy $z^*z$ invertible.  Then the differential
\[
  (dq)_z:\fp\longrightarrow\Sym(W),
  \qquad
  (dq)_z(Y)=z^*Y+Y^*z
\]
is surjective, and
\begin{equation}\label{spin:eq:kerdq}
  \ker(dq)_z=T_z(\Or(U)\cdot z)=\mathfrak{so}(U)z.
\end{equation}
\end{lemma}

\begin{proof}
Since $z^*z$ is invertible, $z$ is injective and $z^*:U\to W$ is surjective.
Given $C=C^*\in\Sym(W)$, choose $Y$ with $z^*Y=C/2$.  Then
$Y^*z=(z^*Y)^*=C/2$, proving surjectivity.

The left $\Or(U)$-action preserves the Gram map, so
$\mathfrak{so}(U)z\subset\ker(dq)_z$.  The kernel dimension is
\[
  rs-\frac{s(s+1)}2.
\]
The stabilizer of $z$ for the left $\Or(U)$-action is $\Or(L^\perp)$, where
$L=z(W)$ and $\dim L^\perp=d$.  Therefore
\[
  \dim \mathfrak{so}(U)z
  =\frac{r(r-1)-d(d-1)}2
  =rs-\frac{s(s+1)}2.
\]
The inclusion is consequently an equality.
\end{proof}

\begin{proposition}[normal and conormal identification]\label{spin:prop:normal-identification}
Let $z=X_J$ and $S_J=z^*z$.  The Gram differential induces a
$\widehat K_z$-equivariant isomorphism of complex algebraic normal spaces
\begin{equation}\label{spin:eq:normal-iso}
  \fp/T_z(\widehat K\cdot z)
  \simeq
  \Sym(W)/T_{S_J}(\Or(W)\cdot S_J).
\end{equation}
Under the trace pairing on $\Sym(W)$, the algebraic conormal is therefore
\begin{equation}\label{spin:eq:alg-conormal}
  C_{\Sym(W)}(S_J)
  :=\{C\in\Sym(W):[C,S_J]=0\}.
\end{equation}
For the underlying real Nash manifold, the complexified conormal appearing in
Chen--Sun is
\begin{equation}\label{spin:eq:real-conormal}
  N_z^*
  \simeq
  C_{\Sym(W)}(S_J)
  \oplus
  \overline{C_{\Sym(W)}(S_J)}.
\end{equation}
\end{proposition}

\begin{proof}
Lemma~\ref{spin:lem:dq} gives the exact sequence
\[
  0\longrightarrow T_z(\Or(U)\cdot z)
  \longrightarrow\fp
  \xrightarrow{(dq)_z}\Sym(W)
  \longrightarrow0.
\]
The image under $(dq)_z$ of the tangent contributed by the right $\Or(W)$
action is precisely the tangent space of the conjugacy orbit of $S_J$.
Quotienting the exact sequence therefore gives \eqref{spin:eq:normal-iso}.

The conjugacy tangent is
\[
  T_{S_J}(\Or(W)\cdot S_J)
  =[\mathfrak{so}(W),S_J].
\]
If $C=C^*$ and $A\in\mathfrak{so}(W)$, cyclicity of trace gives
\[
  \tr\bigl(C[A,S_J]\bigr)
  =\tr\bigl([S_J,C]A\bigr).
\]
The commutator $[S_J,C]$ is skew-adjoint.  Since the trace pairing is
nondegenerate on $\mathfrak{so}(W)$, the displayed expression vanishes for all
$A$ precisely when $[S_J,C]=0$.  This proves \eqref{spin:eq:alg-conormal}.

Finally, the orbit and the ambient vector space are complex submanifolds of
their underlying real Nash manifolds.  If $N^{\mathrm{alg}}$ denotes the
complex algebraic normal space, then the standard complexification identity
\[
  (N^{\mathrm{alg}})_{\R}\otimes_{\R}\C
  \simeq
  N^{\mathrm{alg}}\oplus\overline{N^{\mathrm{alg}}}
\]
gives \eqref{spin:eq:real-conormal} after dualizing.
\end{proof}

Because the $\lambda_j$ are pairwise distinct, the centralizer in
\eqref{spin:eq:alg-conormal} is block diagonal.  We obtain the explicit formula
\begin{equation}\label{spin:eq:conormal-blocks}
  C_{\Sym(W)}(S_J)
  =
  \bigoplus_{j\in J}\Sym(P_j)
  \oplus
  \bigoplus_{j\notin J}\C\{I_{P_j},N_j\}.
\end{equation}
Indeed, when $j\in J$ the block is the scalar $\lambda_jI_{P_j}$ and every
self-adjoint endomorphism centralizes it; when $j\notin J$, the self-adjoint
centralizer of $N_j$ is $\C\{I,N_j\}$.

As a check, if $k=|J|$, then
\[
  \dim_\C C_{\Sym(W)}(S_J)
  =3k+2(m-k)=s+k.
\]
On the other hand, using \eqref{spin:eq:boundary-stabilizer2}, a direct orbit
dimension calculation gives
\[
  \operatorname{codim}_\C(\widehat K\cdot z,\fp)=s+k.
\]
Thus the normal identification has the expected dimension; in particular,
there are no missing directions coming from $L^\perp$.

Fix $j_0\in J$.  Restricting \eqref{spin:eq:real-conormal} and
\eqref{spin:eq:conormal-blocks} to the factor $\Or(P_{j_0})$ gives
\begin{equation}\label{spin:eq:factor-conormal}
  N_z^*|_{\Or(P_{j_0})}
  \simeq
  \Sym(P_{j_0})
  \oplus\overline{\Sym(P_{j_0})}
  \oplus\mathbf1^{\oplus M_J}
\end{equation}
for some $M_J\ge0$.  Every other block and the factor $\Or_d(\C)$ are trivial
under this selected $\Or(P_{j_0})$.

\subsection{The rank-two determinant lemma}\label{spin:sec:detlemma}

We now prove the local character statement that eliminates every boundary
orbit simultaneously.

\begin{lemma}\label{spin:lem:det-ranktwo}
Let $P$ be a complex hyperbolic plane, and regard $\Or(P)=\Or_2(\C)$ as a real
Nash group.  For every $k\ge0$, the determinant character $\det_P$ does not
occur as a subquotient of
\begin{equation}\label{spin:eq:det-ranktwo-module}
  \Sym^k\bigl(\Sym(P)\oplus\overline{\Sym(P)}\bigr).
\end{equation}
The same conclusion holds after adding any number of trivial summands inside
the symmetric power.
\end{lemma}

\begin{proof}
Use the hyperbolic basis $e,f$.  A self-adjoint endomorphism is
\[
  A(a,b,c)=
  \begin{pmatrix}a&b\\c&a\end{pmatrix}.
\]
The identity component
\[
  \SO(P)\simeq\C^\times
\]
consists of $t_z=\diag(z,z^{-1})$.  Conjugation gives the weights
\begin{equation}\label{spin:eq:ranktwo-weights}
  a\mapsto a,
  \qquad
  b\mapsto z^2b,
  \qquad
  c\mapsto z^{-2}c.
\end{equation}
On the conjugate copy the corresponding weights are
$1,\bar z^2,\bar z^{-2}$.

Consider a monomial in the symmetric algebra.  Let $u$ be the exponent of the
holomorphic $+2$ variable minus the exponent of the holomorphic $-2$ variable,
and let $v$ be the analogous difference in the conjugate copy.  Its
$\SO(P)$-character is
\[
  z^{2u}\bar z^{2v}.
\]
If this character is trivial for every $z=re^{i\theta}\in\C^\times$, then the
radial and angular parts give
\[
  u+v=0,
  \qquad
  u-v=0.
\]
Hence $u=v=0$.  Thus every $\SO(P)$-invariant monomial contains equally many
$+2$ and $-2$ variables in each of the two copies separately.

Let $\tau\in\Or(P)$ be the reflection interchanging $e$ and $f$.  It exchanges
the $+2$ and $-2$ variables and fixes the weight-zero variable, in both the
holomorphic and the conjugate copies.  Every $\SO(P)$-invariant monomial is
therefore fixed by $\tau$.  Consequently $\tau$ acts trivially on the entire
$\SO(P)$-fixed subspace.

The determinant character is trivial on $\SO(P)$ and takes the value $-1$ on
$\tau$.  It therefore cannot occur.  The finite-dimensional modules here are
rational representations of the reductive real algebraic group obtained from
$\Or_2(\C)$ by restriction of scalars, hence are completely reducible; absence
as a subrepresentation is equivalent to absence as a subquotient.

Adding trivial summands merely tensors the weight calculation with a polynomial
algebra on which $\Or(P)$ acts trivially, and does not alter the conclusion.
\end{proof}

\subsection{Automatic extension in the even case}\label{spin:sec:automatic}

We now verify the hypotheses of Theorem~\ref{spin:thm:ChenSun} for
\begin{equation}\label{spin:eq:CS-data}
  G=\widehat K,
  \qquad
  M=R(\fp),
  \qquad
  U_0=V=R(\fp)\setminus B,
  \qquad
  \eta=\det_W.
\end{equation}
The group $\widehat K$ is an almost linear Nash group.  By
Proposition~\ref{spin:prop:p-closure}, the boundary $B=M\setminus V$ is the finite
union of the orbits $\widehat KX_J$ with $J\ne\varnothing$.

Fix such a boundary point $z=X_J$ and choose $j_0\in J$.  By
\eqref{spin:eq:modular-trivial},
\[
  \delta_{\widehat K/\widehat K_z}=1.
\]
By \eqref{spin:eq:detW-boundary}, the restriction of $\det_W$ to the selected
factor $\Or(P_{j_0})$ is $\det_{P_{j_0}}$.  Finally,
\eqref{spin:eq:factor-conormal} describes the restriction of the full complexified
real conormal.

\begin{proposition}\label{spin:prop:CS-vanishing}
For every boundary point $z=X_J$ and every $k\ge0$, the trivial
representation of $\widehat K_z$ does not occur as a subquotient of
\begin{equation}\label{spin:eq:CS-module-here}
  \Sym^k(N_z^*)\otimes\det_W|_{\widehat K_z}.
\end{equation}
\end{proposition}

\begin{proof}
Suppose the trivial representation occurred.  The module in
\eqref{spin:eq:CS-module-here} is a finite-dimensional rational representation of a
reductive real algebraic group, hence is completely reducible.  It would
therefore contain a nonzero $\widehat K_z$-fixed vector.  Restrict this vector
to the factor $\Or(P_{j_0})$.

Under that factor, $N_z^*$ has the form
\[
  \Sym(P_{j_0})\oplus\overline{\Sym(P_{j_0})}
  \oplus\mathbf1^{\oplus M_J},
\]
while the character twist is $\det_{P_{j_0}}$.  Decomposing the symmetric
power of a direct sum according to degree in the trivial variables reduces the
existence of a fixed vector to the occurrence of $\det_{P_{j_0}}$ in some
symmetric power of
$\Sym(P_{j_0})\oplus\overline{\Sym(P_{j_0})}$.  This is excluded by
Lemma~\ref{spin:lem:det-ranktwo}.
\end{proof}

We can now cross the entire finite boundary.

\begin{proposition}\label{spin:prop:even-regular-distribution}
If $s$ is even, then
\begin{equation}\label{spin:eq:regular-sign-nonzero}
  \cS^*(R(\fp))^{\widehat K,\det_W}\ne0.
\end{equation}
Consequently
\[
  \cS^*(R(\fp))^{K,\chi}\ne0.
\]
\end{proposition}

\begin{proof}
The distribution $T_{\cO}$ in \eqref{spin:eq:TO-Schwartz} is a nonzero
$\det_W$-equivariant Schwartz distribution on $V$.  Proposition
\ref{spin:prop:CS-vanishing}, together with
\eqref{spin:eq:modular-trivial}, verifies the hypothesis
\eqref{spin:eq:CS-general} of Theorem~\ref{spin:thm:ChenSun} at every boundary orbit and
for every conormal order.  The restriction map
\[
  \cS^*(R(\fp))^{\widehat K,\det_W}
  \longrightarrow
  \cS^*(V)^{\widehat K,\det_W}
\]
is therefore an isomorphism.  Hence $T_{\cO}$ has a
$\det_W$-equivariant Schwartz extension $\widetilde T$ to $R(\fp)$.  Since
its restriction is $T_{\cO}\ne0$, the extension is nonzero.  This proves
\eqref{spin:eq:regular-sign-nonzero}.

Finally $\det_W|_K=\chi$, so the same distribution belongs to
$\cS^*(R(\fp))^{K,\chi}$.
\end{proof}

\subsection{Completion of the Spin block family in ambient dimension at least five}\label{spin:sec:conclusion}

We now combine the two parities.

\begin{proof}[Proof of Theorem~\ref{thm:Spin-residual}]
Interchange the blocks if necessary so that $r>s$.  By
Proposition~\ref{spin:prop:sign-reduction}, it is enough to prove the implication
\eqref{spin:eq:SR}.

If $s$ is odd, Corollary~\ref{spin:cor:odd-sign-zero} gives
\[
  \cS^*(\fp)^{K,\chi}=0.
\]
Thus the conclusion of \eqref{spin:eq:SR} always holds.

If $s$ is even, Proposition~\ref{spin:prop:even-regular-distribution} gives
\[
  \cS^*(R(\fp))^{K,\chi}\ne0.
\]
Thus the premise of \eqref{spin:eq:SR} is false.

These two cases exhaust all possibilities for $s$, so \eqref{spin:eq:SR} always
holds.  The pair is regular.
\end{proof}

\begin{proof}[Proof of Corollary~\ref{cor:Spin-all}]
Assume $r+s\ge5$.  The Spin block cases with $|r-s|\le2$ are among the
nice symmetric Lie algebra pairs in Sekiguchi's classification and hence are
regular \cite[Lemmas~4.21--4.22 and Table~3]{Rubio}.  If $|r-s|>2$ and both
block dimensions are odd, Proposition~\ref{spin:prop:pleasant} makes the pair
pleasant, and pleasant pairs are regular.  In every remaining case $rs$ is
even, so Theorem~\ref{thm:Spin-residual} applies.  In particular, no
rank-one block is omitted: if, say, $s=1$ and $r+s\ge5$, then
$|r-s|>2$; odd $r$ is the unequal odd--odd pleasant case, whereas even $r$ is
in the residual $rs$-even case.
\end{proof}


\begin{corollary}[Spin component implication]\label{cor:Spin-component}
For every nonpleasant Spin block factor needed in Rubio's reduction, the
finite-component datum is componentwise regular.
\end{corollary}

\begin{proof}
In the residual range the obstruction quotient has order two and
Proposition~\ref{spin:prop:sign-reduction} identifies its nontrivial character
with $\chi$.  For odd smaller block,
Corollary~\ref{spin:cor:odd-sign-zero} makes the full $Q$-sector zero.  For even
smaller block, Proposition~\ref{spin:prop:even-regular-distribution} constructs
a nonzero regular-set distribution in that sector.  In the remaining Spin
ranges the pair is pleasant or nice, as proved in
Proposition~\ref{spin:prop:pleasant} and the proof of
Corollary~\ref{cor:Spin-all}.
\end{proof}

\section{The EVII pair}\label{sec:EVII-family}
Let $G$ be the simply connected complex group of type $E_7$ and let $\theta$
be the involution of type EVII.  Put
\[
 H=G^\theta,
 \qquad
 \fg=\fh\oplus\fp.
\]
The Lie algebra $\fh$ has type $E_6+\C$.  The isotropy representation has the
standard form
\begin{equation}\label{evii:eq:EVII-p}
 \fp\simeq J\oplus J^*,
\end{equation}
where $J$ is the $27$-dimensional minuscule $E_6$-module and the connected
central torus of $H$ acts by mutually inverse nontrivial characters on $J$ and
$J^*$.  Hence
\begin{equation}\label{evii:eq:EVII-dim}
 \dim_{\C}\fp=54,
 \qquad
 \fp^H=0.
\end{equation}
The invariant symmetric form is
\begin{equation}\label{evii:eq:EVII-B}
 B((x,\varphi),(y,\psi))=\varphi(y)+\psi(x).
\end{equation}
For $t\in\C^\times$, the linear map
$T_t(x,\varphi)=(tx,\varphi)$ satisfies
$B(T_tu,T_tv)=tB(u,v)$.  Hence $B$ and $tB$ are isometric for every $t$, so
$\delta_B=1$.

We prove
\begin{theorem}\label{evii:thm:EVII-special}
The complex symmetric pair $(E_7,E_6+\C)$ is special.  Hence it is weakly
linearly tame and regular.
\end{theorem}

\subsection{Nilpotent orbits and the trace reduction}

Djokovi\'c gives explicit Cayley triples for all nilpotent adjoint orbits in
the real form $E_{7(-25)}$ \cite[Table~6]{Djokovic}.  By the
Kostant--Sekiguchi correspondence \cite[Chapter~9]{CM}, these are in
bijection with nilpotent $H$-orbits in the complexified EVII isotropy
representation.  There are $22$ nonzero orbits.

Djokovi\'c uses triples $(E,H_0,F)$ satisfying
\[
 [H_0,E]=2E,
 \qquad
 [H_0,F]=-2F,
 \qquad
 [F,E]=H_0.
\]
We pass to the normal triple
\begin{equation}\label{evii:eq:EVII-Cayley-transform}
 e=\frac{E-F+iH_0}{2},
 \qquad
 f=\frac{E-F-iH_0}{2},
 \qquad
 d=i(E+F).
\end{equation}
A direct bracket calculation gives \eqref{core:eq:normal-triple-general}; the EVII
involution satisfies $\theta(d)=d$, $\theta(e)=-e$, and $\theta(f)=-f$.

For $j\geq0$ set
\[
 \fp_j(e)=\{x\in\fp:[d,x]=jx\},\qquad
 \fh_j(e)=\{x\in\fh:[d,x]=jx\},
\]
and write $p_j(e)=\dim\fp_j(e)$ and $h_j(e)=\dim\fh_j(e)$.
The following elementary observation turns the exact $d$-weight dimensions
into the highest-weight decompositions used below.

\begin{lemma}[Highest multiplicities from graded dimensions]
\label{evii:lem:graded-multiplicities}
Let $m_{j,\fp}$ and $m_{j,\fh}$ denote the multiplicities of the irreducible
graded $\mathfrak{sl}_2$-modules of highest weight $j$ whose highest lines lie
in $\fp$ and $\fh$, respectively.  Then, for every $j\geq0$,
\begin{equation}\label{evii:eq:graded-difference}
 m_{j,\fp}=p_j(e)-h_{j+2}(e),\qquad
 m_{j,\fh}=h_j(e)-p_{j+2}(e),
\end{equation}
where a missing weight space has dimension zero.  In particular,
\begin{align}
 \fp^e&\simeq\bigoplus_{j\geq0}m_{j,\fp}V_j^{\fp},
 \label{evii:eq:p-highest-from-weights}\\
 T(e)&=\sum_{j\geq0}j\,m_{j,\fh}.
 \label{evii:eq:T-from-weights}
\end{align}
\end{lemma}

\begin{proof}
On every finite-dimensional irreducible $\mathfrak{sl}_2$-module, the raising
operator $e$ maps the weight-$j$ line isomorphically onto the weight-$(j+2)$
line whenever $j\geq0$ and the latter line occurs.  Since $e\in\fp$, it
interchanges the two $\theta$-parities.  Hence
\[
 e:\fp_j(e)\longrightarrow\fh_{j+2}(e),\qquad
 e:\fh_j(e)\longrightarrow\fp_{j+2}(e)
\]
are surjective.  Their kernels are precisely the highest lines of weight $j$
in the indicated parity.  Taking dimensions proves
\eqref{evii:eq:graded-difference}; the remaining assertions follow from the
definitions of $\fp^e$ and $T(e)$.
\end{proof}

\subsubsection{The finite calculation}

We compute the rank and weight-space data for the twenty-two orbits in
Djokovi\'c's ordered
Chevalley basis
\[
 H_1,\ldots,H_7,\quad X_1,\ldots,X_{63},\quad
 X_{-1},\ldots,X_{-63}.
\]
We use the Bourbaki numbering of $E_7$, with edges
$(1,3),(3,4),(4,5),(5,6),(6,7),(2,4)$, and the corresponding Cartan matrix.
If the positive root numbered $r$ is
$\alpha_r=\sum_{k=1}^7a_{rk}\alpha_k$, our conventions are
\begin{align*}
 [H_k,X_{\pm r}]&=\pm\alpha_r(H_k)X_{\pm r},\\
 [X_r,X_{-r}]&=-\sum_{k=1}^7a_{rk}H_k,
 & [X_{-r},X_r]&=\sum_{k=1}^7a_{rk}H_k.
\end{align*}
For all other root brackets we use the signed constants $N(r,s)$ printed in
Djokovi\'c's Appendix.  That table lists the constants with positive first
index, including the mixed-sign cases.  We extend them by
\[
 N(-r,s)=-N(s,-r),\qquad N(-r,-s)=N(r,s),
\]
whenever the indicated root sum exists, and use skew-symmetry; all remaining
root-vector brackets are zero.  Thus the bracket table is determined without
any further choice.  An exact implementation of the finite row reductions in
this subsection is included as supplementary material.  Let $\sigma$ be the
conjugation in
\cite[Table~3]{Djokovic} and let $\sigma_u$ be the compact conjugation
$\sigma_u(H_k)=-H_k$, $\sigma_u(X_r)=X_{-r}$.  The EVII involution is
$\theta=\sigma_u\sigma$.  Direct substitution in the signed bracket table
gives $\theta^2=1$, $\theta([x,y])=[\theta x,\theta y]$, and
\[
 \dim\ker(\theta-1)=79,\qquad \dim\ker(\theta+1)=54.
\]

For each row below, if $E=\sum_rc_rX_r$, then
$F=\sum_rc_rX_{-r}$.  A sign $\pm$ gives the two numbered complex
$H$-orbits.  The entries are the Cayley triples in
\cite[Table~6]{Djokovic}; they also fix unambiguously the numbering used in
Proposition~\ref{evii:prop:EVII-orbit-invariants}.

\begingroup
\small
\renewcommand{\arraystretch}{1.08}
\begin{longtable}{@{}c >{\raggedright\arraybackslash}p{0.27\textwidth}
 >{\raggedright\arraybackslash}p{0.56\textwidth}@{}}
\caption{EVII Cayley data}\label{evii:tab:cayley-data}\\
\toprule
orbit(s)&$H_0=\sum_{k=1}^7b_kH_k$&$E$\\
\midrule
\endfirsthead
\toprule
orbit(s)&$H_0=\sum_{k=1}^7b_kH_k$&$E$\\
\midrule
\endhead
$1,2$ & $(2,2,3,4,3,2,1)$ & $\pm X_{63}$\\
$3,4$ & $(2,3,4,6,5,4,2)$ & $\pm(X_{49}-X_{63})$\\
$5$ & $(2,3,4,6,5,4,2)$ & $X_{49}+X_{63}$\\
$6,7$ & $(2,3,4,6,5,4,3)$ & $\pm(X_7-X_{49}+X_{63})$\\
$8,9$ & $(2,3,4,6,5,4,3)$ & $\pm(X_7+X_{49}+X_{63})$\\
$10$ & $(4,4,6,8,6,4,2)$ & $X_1+X_{37}+X_{55}+X_{61}$\\
$11,12$ & $(4,5,7,10,8,6,3)$ &
 $\pm(X_{27}+X_{39}+X_{49}+X_{53}-X_{54})$\\
$13,14$ & $(6,7,10,14,11,8,4)$ &
 $\pm(\sqrt3\,X_1+2X_{49}+\sqrt3\,X_{37})$\\
$15$ & $(4,6,8,12,10,8,4)$ &
 $\sqrt2\,(X_6+X_{34}+X_{40}+X_{56})$\\
$16,19$ & $(6,7,10,14,11,8,5)$ &
 $\pm(\sqrt3\,X_1-X_7+2X_{49}+\sqrt3\,X_{37})$\\
$17,18$ & $(6,7,10,14,11,8,5)$ &
 $\pm(\sqrt3\,X_1+X_7+2X_{49}+\sqrt3\,X_{37})$\\
$20$ & $(8,10,14,20,16,12,6)$ &
 $2X_1+\sqrt3\,X_6+\sqrt3\,X_{19}+2X_{37}+\sqrt3\,X_{40}+\sqrt3\,X_{41}$\\
$21,22$ & $(10,13,18,26,21,16,9)$ &
 $\pm(\sqrt5\,X_1+2\sqrt2\,X_6+3X_7+\sqrt5\,X_{37}+2\sqrt2\,X_{40})$\\
\bottomrule
\end{longtable}
\endgroup

\begin{remark}[The representative for orbits $6,7$]
\label{evii:rem:djokovic-sign}
The representative printed in \cite[Table~6]{Djokovic} is
\[
  \pm(X_7-X_{49}-X_{63}).
\]
With the structure constants and sign extension from Djokovi\'c's Appendix
used above, this printed vector satisfies the $\mathfrak{sl}_2$ relations
with $H_0=(2,3,4,6,5,4,3)$, but the exact weight calculation gives
\[
  (\operatorname{tr},\operatorname{inv})=(128,47),
\]
which is the invariant pair of the adjacent $8,9$ row.  Replacing only the
last sign gives
\[
  \pm(X_7-X_{49}+X_{63}),
\]
and again the $\mathfrak{sl}_2$ relations hold, while now
\[
  (\operatorname{tr},\operatorname{inv})=(0,79),
\]
which is the pair assigned to row $6,7$.  Both calculations use the same
structure constants and the same sign extension for $N(r,s)$, so the
discrepancy does not arise from that convention.
Accordingly, we use $\pm(X_7-X_{49}+X_{63})$ for orbits $6,7$ in
Table~\ref{evii:tab:cayley-data} and below.
\end{remark}

For a displayed triple define
\[
 \mathcal P_e(q)=\sum_{j\geq0}p_j(e)q^j,
 \qquad
 \mathcal H_e(q)=\sum_{j\geq0}h_j(e)q^j.
\]
The dimensions are obtained by the following explicitly displayed finite
systems on the $133$-dimensional Chevalley basis:
\begin{equation}\label{evii:eq:weight-kernels}
\begin{aligned}
 p_j(e)&=\dim\ker
 \begin{pmatrix}\theta+1\\ \ad d-j\end{pmatrix},\\
 h_j(e)&=\dim\ker
 \begin{pmatrix}\theta-1\\ \ad d-j\end{pmatrix}.
\end{aligned}
\end{equation}
All entries lie in $\Q(i,\sqrt2,\sqrt3,\sqrt5)$.  Exact row reduction of
\eqref{evii:eq:weight-kernels}, using the signed bracket conventions above,
gives the following complete list; omitted powers have coefficient zero.

\begingroup
\small
\renewcommand{\arraystretch}{1.08}
\begin{longtable}{@{}c >{\raggedright\arraybackslash}p{0.41\textwidth}
 >{\raggedright\arraybackslash}p{0.41\textwidth}@{}}
\caption{Nonnegative $d$-weight dimensions in $\fp$ and $\fh$}
\label{evii:tab:weight-dimensions}\\
\toprule
orbit(s)&$\mathcal P_e(q)$&$\mathcal H_e(q)$\\
\midrule
\endfirsthead
\toprule
orbit(s)&$\mathcal P_e(q)$&$\mathcal H_e(q)$\\
\midrule
\endhead
$1,2$   &$20+16q+q^2$&$47+16q$\\
$3,4$   &$2+16q+10q^2$&$47+16q$\\
$5$     &$18+16q+2q^2$&$31+16q+8q^2$\\
$6,7$   &$27q^2$&$79$\\
$8,9$   &$32+11q^2$&$47+16q^2$\\
$10$    &$30+12q^2$&$37+20q^2+q^4$\\
$11,12$ &$12+12q+5q^2+4q^3$&$21+12q+12q^2+4q^3+q^4$\\
$13,14$ &$2+8q+9q^2+8q^3+q^6$&$31+8q+8q^3+8q^4$\\
$15$    &$18+16q^2+2q^4$&$31+16q^2+8q^4$\\
$16,19$ &$16+10q^2+8q^4+q^6$&$31+16q^2+8q^4$\\
$17,18$ &$26q^2+q^6$&$47+16q^4$\\
$20$    &$14+10q^2+8q^4+2q^6$&$19+14q^2+9q^4+6q^6+q^8$\\
$21,22$ &$17q^2+9q^6+q^{10}$&$31+16q^4+8q^8$\\
\bottomrule
\end{longtable}
\endgroup
For every row the symmetry of the $d$-spectrum gives the checks
\[
 p_0(e)+2\sum_{j>0}p_j(e)=54,
 \qquad
 h_0(e)+2\sum_{j>0}h_j(e)=79.
\]
Thus the table accounts for all of $\fp$ and $\fh$, not merely their
highest spaces.

For an irreducible graded $\mathfrak{sl}_2$-module, write
$V_\ell^{\fh}$ or $V_\ell^{\fp}$ according as its highest line lies in
$\fh$ or $\fp$.  Lemma~\ref{evii:lem:graded-multiplicities} and
Table~\ref{evii:tab:weight-dimensions} give the following proposition.

\begin{proposition}[EVII orbit invariants]\label{evii:prop:EVII-orbit-invariants}
With Djokovi\'c's orbit numbering, the exact data are as follows.  In the
second column, $V_\ell$ denotes a $\fp$-parity highest module.
\begin{center}
\renewcommand{\arraystretch}{1.10}
\begin{tabular}{@{}c >{$}l<{$} r r@{}}
\toprule
orbit(s) & \multicolumn{1}{c}{$\fp$-highest modules} & $T(e)$ & $T(e)-54$\\
\midrule
$1,2$   & V_2+16V_1+20V_0                     & 16 & -38\\
$3,4$   & 10V_2+16V_1+2V_0                    & 16 & -38\\
$5$     & 2V_2+16V_1+10V_0                    & 32 & -22\\
$6,7$   & 27V_2                               & 0  & -54\\
$8,9$   & 11V_2+16V_0                         & 32 & -22\\
$10$    & 11V_2+10V_0                         & 44 & -10\\
$11,12$ & 4V_3+4V_2+8V_1                      & 48 & -6\\
$13,14$ & V_6+8V_3+V_2+2V_0                   & 52 & -2\\
$15$    & 2V_4+8V_2+2V_0                      & 60 & 6\\
$16,19$ & V_6+8V_4+2V_2                       & 44 & -10\\
$17,18$ & V_6+10V_2                           & 60 & 6\\
$20$    & V_6+2V_4+V_2                        & 84 & 30\\
$21,22$ & V_{10}+V_6+V_2                      & 84 & 30\\
\bottomrule
\end{tabular}
\end{center}
\end{proposition}

\begin{proof}
Apply \eqref{evii:eq:graded-difference} row by row to
Table~\ref{evii:tab:weight-dimensions}, and then use
\eqref{evii:eq:T-from-weights}.  For example, in orbit $15$,
\[
 \mathcal P_e=18+16q^2+2q^4,\qquad
 \mathcal H_e=31+16q^2+8q^4,
\]
so
\[
 (m_{0,\fp},m_{2,\fp},m_{4,\fp})=(2,8,2),\qquad
 (m_{0,\fh},m_{2,\fh},m_{4,\fh})=(15,14,8),
\]
and $T(e)=2\cdot14+4\cdot8=60$.  In orbit $20$ the same subtraction gives
\[
 (m_{2,\fp},m_{4,\fp},m_{6,\fp})=(1,2,1)
\]
and
\[
 (m_{0,\fh},m_{2,\fh},m_{4,\fh},m_{6,\fh},m_{8,\fh})
 =(9,6,7,6,1),
\]
whence $T(e)=2\cdot6+4\cdot7+6\cdot6+8=84$.  The remaining rows are identical
one-line subtractions and give the displayed table.
\end{proof}

For the six orbits with nonnegative defect, the same calculation gives the
$\fh$-parity highest modules
\begin{align}
15:&\quad 15V_0^{\fh}\oplus14V_2^{\fh}\oplus8V_4^{\fh},
\label{evii:eq:EVII-h15}\\
17,18:&\quad 21V_0^{\fh}\oplus15V_4^{\fh},
\label{evii:eq:EVII-h17}\\
20:&\quad 9V_0^{\fh}\oplus6V_2^{\fh}\oplus7V_4^{\fh}
       \oplus6V_6^{\fh}\oplus V_8^{\fh},
\label{evii:eq:EVII-h20}\\
21,22:&\quad 14V_0^{\fh}\oplus7V_4^{\fh}\oplus7V_8^{\fh}.
\label{evii:eq:EVII-h21}
\end{align}
Every highest weight in \eqref{evii:eq:EVII-h15}--\eqref{evii:eq:EVII-h21} is even.
The $\fp$-highest weights in the same cases are also even.  Consequently the
normal quotient lines are exactly the lowest lines attached to the displayed
$\fp$-highest modules; no additional line arises from an odd
$\fh$-highest module.

The zero orbit is excluded by Lemma~\ref{core:lem:origin-adapted}.  By
Proposition~\ref{evii:prop:EVII-orbit-invariants} and
Corollary~\ref{core:cor:strict-trace}, the nonzero orbits
\begin{equation}\label{evii:eq:EVII-strict-orbits}
 1\text{--}14,\qquad16,\qquad19
\end{equation}
are excluded.  It remains to treat
\begin{equation}\label{evii:eq:EVII-remaining-orbits}
 15,\qquad17,18,\qquad20,\qquad21,22.
\end{equation}

\subsection{Congruence obstructions}

\begin{proposition}\label{evii:prop:EVII-congruence}
The orbits $17,18,21,22$ contribute no adapted distribution.
\end{proposition}

\begin{proof}
For orbits $17,18$, the defect is six and
\[
 \fp^e\simeq V_6^{\fp}\oplus10V_2^{\fp}.
\]
The elementary degrees from Lemma~\ref{core:lem:resonance-general} are $8$ and
$4$.  Every holomorphic weighted degree is divisible by four, whereas the
required type-I degree is six.  Type II still requires one side of degree six.
Thus both types are impossible.

For orbits $21,22$, the defect is thirty and
\[
 \fp^e\simeq V_{10}^{\fp}\oplus V_6^{\fp}\oplus V_2^{\fp}.
\]
The elementary degrees are $12,8,4$, so every holomorphic degree is divisible
by four, whereas thirty is not.  Again neither adapted type occurs.
\end{proof}

\subsection{A central-torus refinement}

For the remaining orbits $15$ and $20$, the normal one-parameter subgroup
admits the numerical degree.  Let
\begin{equation}\label{evii:eq:EVII-triple-centralizer}
 \mathfrak c_{\fh}(e,d,f)
 =\{z\in\fh:[z,e]=[z,d]=[z,f]=0\}.
\end{equation}
The centralizer of an $\mathfrak{sl}_2$-triple is reductive, so the connected
center of \eqref{evii:eq:EVII-triple-centralizer} exponentiates to an algebraic
torus in $H_e$.

\begin{lemma}[No hidden modular twist]\label{evii:lem:EVII-modular-torus}
For each of the orbits $15$ and $20$, let $Z\simeq\C^\times$ be the connected
central torus of \eqref{evii:eq:EVII-triple-centralizer}.  Then
\begin{equation}\label{evii:eq:EVII-modular-Z}
 \Delta_{(H\times\C^\times)_e}|_Z=1.
\end{equation}
Thus, on $Z$, the Frobenius fiber has only the character contributed by the
normal tensor.
\end{lemma}

\begin{proof}
The decompositions \eqref{evii:eq:EVII-h15} and \eqref{evii:eq:EVII-h20} show that
every $\mathfrak{sl}_2$ highest weight occurring in $\fh^e$ is even.  Let
$\kappa$ be the Killing form of $\fg$.  On each even isotypic component, the
highest and lowest lines have the same $\theta$-parity.  Since
$\fh\perp\fp$ for $\kappa$, the restriction of $\kappa$ induces a
nondegenerate $Z$-invariant bilinear form on the multiplicity space of every
$V_\ell^{\fh}$.  Hence each multiplicity space is self-dual as a $Z$-module.
Its determinant character is equal to its inverse.  A connected algebraic
torus has no nontrivial character of order two, so the determinant character
is trivial.  Multiplying over the isotypic components gives
\[
 \det(\Ad(z)|_{\fh^e})=1.
\]
With the action convention $\rho(\lambda)e=\lambda^{-1}e$, the stabilizer
Lie algebra is
\[
 \Lie((H\times\C^\times)_e)
 =\{(X,s)\in\fh\oplus\C:[X,e]-se=0\}.
\]
It is the direct sum of $\fh^e$ and the line represented by $(d,2)$, because
$[d,e]=2e$.  The torus $Z$ centralizes $d$ and acts trivially on the external
$\C^\times$ factor, so it fixes this complementary line.  Together with
$\det(\Ad(z)|_{\fh^e})=1$, this proves that the full stabilizer modular
character is trivial on $Z$.
\end{proof}

\begin{lemma}\label{evii:lem:EVII-holomorphic-test}
Assume \eqref{evii:eq:EVII-modular-Z}.  If the holomorphic symmetric algebra of
$N_e$ has no $Z$-invariant monomial of weighted degree $T(e)-54$, then the
orbit $H\cdot e$ contributes no $B$-adapted distribution.
\end{lemma}

\begin{proof}
After complexifying the underlying real normal space, a tensor splits into
holomorphic and antiholomorphic factors.  A torus character on the two factors
has the form $z^a\bar z^b$.  If it is identically one on the real Lie group
$\C^\times$, then $a=b=0$: first restrict to positive real $z$, then to the
unit circle.  Thus invariance forces the two characters to vanish separately.
In type I, both weighted degrees equal $T(e)-54$; in type II, one of them still
equals $T(e)-54$.  The antiholomorphic statement is the conjugate of the
holomorphic one.
\end{proof}

\subsection{Orbit 20}

All root vectors below use Djokovi\'c's numbering and normalization
\cite[Tables~1 and~3 and Appendix]{Djokovic}.  For $r>0$ write
\begin{equation}\label{evii:eq:SA-notation}
 S_r=X_r+X_{-r},\qquad A_r=X_r-X_{-r}.
\end{equation}

\begin{lemma}[The triple centralizer for orbit 20]\label{evii:lem:EVII-orbit20-centralizer}
For orbit $20$, the Lie algebra \eqref{evii:eq:EVII-triple-centralizer} has dimension
nine and one-dimensional center.  A generator of its center is
\begin{equation}\label{evii:eq:EVII-z20}
 z_{20}=3(X_7+X_{-7})-2(X_5+X_{-5})
 +(X_3+X_{-3})-(X_{28}+X_{-28}).
\end{equation}
If
\[
 M_\ell^{\fp}=\ker(\ad e)\cap\ker(\ad d-\ell)\cap\fp,
\]
then
\[
 \dim(M_6^{\fp},M_4^{\fp},M_2^{\fp})=(1,2,1)
\]
and
\begin{equation}\label{evii:eq:EVII-z20-char}
\begin{array}{c|ccc}
\ell&6&4&2\\ \hline
\det(s-\ad z_{20}|_{M_\ell^{\fp}})&s&s^2+36&s.
\end{array}
\end{equation}
Consequently, the connected central torus acts trivially on
$M_6^{\fp}$ and $M_2^{\fp}$ and by reciprocal nontrivial characters on
$M_4^{\fp}$.
\end{lemma}

\begin{proof}
For this row
\[
 E_{20}=2X_1+\sqrt3\,X_6+\sqrt3\,X_{19}+2X_{37}
          +\sqrt3\,X_{40}+\sqrt3\,X_{41},
 \qquad
 H_{0,20}=(8,10,14,20,16,12,6),
\]
and $F_{20}$ is obtained by replacing every $X_r$ by $X_{-r}$.
Write
\[
 S_r=X_r+X_{-r},\qquad A_r=X_r-X_{-r}.
\]
Consider the exact linear map
\[
 \Phi_{20}:\fh\longrightarrow\fg^{\oplus3},\qquad
 x\longmapsto([x,e],[x,d],[x,f]).
\]
Substitution of the Cayley transform into the signed Chevalley bracket table
and row reduction give $\rk\Phi_{20}=70$.  Its nine-dimensional kernel has the
following basis:
\begin{align*}
 c_1&=H_2,& c_2&=X_2,&c_3&=X_{-2},\\
 c_4&=S_5-S_7,
 &c_5&=X_9-X_{22}-X_{-10}-X_{-11},\\
 c_6&=X_4+X_{17}-X_{-15}-X_{-16},
 &c_7&=X_{15}+X_{16}-X_{-4}-X_{-17},\\
 c_8&=-X_{10}-X_{11}+X_{-9}-X_{-22},
 &c_9&=S_3+S_5-S_{28}.
\end{align*}
Each displayed vector is fixed by $\theta$ and has zero bracket with $e,d,f$;
the rank computation proves that they exhaust the triple centralizer.
If $x=\sum_{r=1}^9a_rc_r$, the equations $[x,c_r]=0$ for all $r$ reduce to
\[
 a_1=a_2=a_3=a_5=a_6=a_7=a_8=0,
 \qquad a_4+3a_9=0.
\]
Thus the center is generated by $c_9-3c_4$, which is precisely the element
$z_{20}$ in \eqref{evii:eq:EVII-z20}.

It remains to compute the action on the three highest spaces.  For
$\ell\in\{6,4,2\}$ set
\[
 \Psi_\ell:\fp\longrightarrow\fg\oplus\fg,
 \qquad x\longmapsto([e,x],[d,x]-\ell x).
\]
The respective ranks are $53,52,53$.  Hence the kernels have dimensions
$1,2,1$.  The following vectors form convenient ordered bases:
\begingroup
\small
\begin{align*}
 a_6={}&4\sqrt3iH_1+2\sqrt3iH_3-2\sqrt3iH_5
        -4\sqrt3iH_6-2\sqrt3iH_7\\
 &+3\sqrt3A_1-3A_6-3A_{19}+3\sqrt3A_{37}-3A_{40}-3A_{41}\\
 &+3iA_{42}-3iA_{43}-3iA_{44}-3iA_{51}
   +\sqrt3A_{55}+\sqrt3A_{61},\\[2mm]
 a_4={}&-6iA_7-2\sqrt3A_{12}+2\sqrt3A_{13}-2\sqrt3A_{35}
        -\sqrt3iA_{38}-2\sqrt3A_{45}\\
 &\hspace{2.4em}-\sqrt3iA_{46}+\sqrt3iA_{47}-\sqrt3iA_{48},\\[2mm]
 b_4={}&-6iH_7+2\sqrt3A_6-2\sqrt3A_{19}
        +2\sqrt3A_{40}-2\sqrt3A_{41}\\
 &\hspace{2.4em}-\sqrt3iA_{42}+\sqrt3iA_{43}
        -\sqrt3iA_{44}-\sqrt3iA_{51},\\[2mm]
 a_2={}&8H_1+10H_2+14H_3+20H_4+16H_5+12H_6+6H_7\\
 &-2iA_1-\sqrt3i(A_6+A_{19}+A_{40}+A_{41})-2iA_{37}.
\end{align*}
\endgroup
Direct bracket substitution gives
\[
 [z_{20},a_6]=[z_{20},a_2]=0,
 \qquad
 [z_{20},a_4]=6b_4,
 \qquad
 [z_{20},b_4]=-6a_4.
\]
Therefore the matrices of $\ad z_{20}$ on the ordered bases
$(a_6)$, $(a_4,b_4)$, and $(a_2)$ are
\[
 (0),\qquad
 \begin{pmatrix}0&-6\\6&0\end{pmatrix},\qquad
 (0),
\]
which proves \eqref{evii:eq:EVII-z20-char}.  Since the center of a reductive
triple centralizer is toral, the two nonzero eigenvalues are the differentials
of reciprocal nontrivial algebraic characters.  The exclusion argument below
uses only this qualitative dichotomy---trivial versus reciprocal nontrivial
characters---not the numerical values $\pm6i$ themselves.
\end{proof}

\begin{proposition}\label{evii:prop:EVII-orbit20}
Orbit $20$ contributes no adapted distribution.
\end{proposition}

\begin{proof}
The duality \eqref{core:eq:normal-duality-general} inverts the $Z$-characters when
passing from the highest spaces $M_\ell^{\fp}\subset\fp^e$ to the normal
quotient.  In \eqref{evii:eq:EVII-z20-char} the nontrivial characters already occur
as a reciprocal pair, while the other characters are trivial; hence the same
character multiset occurs on the normal variables.  The defect is thirty.  The elementary degrees attached to
$\ell=6,4,2$ are $8,6,4$.  By Lemma~\ref{evii:lem:EVII-orbit20-centralizer}, the
degree-eight and degree-four variables are torus-trivial, while the two
degree-six variables have reciprocal nontrivial characters.  In a torus-invariant
monomial, the latter variables occur with equal total exponent.  Therefore every
invariant holomorphic weighted degree is
\[
 8a+4c+12b,
 \qquad a,b,c\in\Z_{\ge0},
\]
which is divisible by four.  It cannot equal thirty.  Lemmas
\ref{evii:lem:EVII-modular-torus} and \ref{evii:lem:EVII-holomorphic-test} conclude the
proof.
\end{proof}

\subsection{Orbit 15}

\begin{lemma}[The triple centralizer for orbit 15]\label{evii:lem:EVII-orbit15-centralizer}
For orbit $15$, the Lie algebra \eqref{evii:eq:EVII-triple-centralizer} has dimension
fifteen and one-dimensional center.  A generator is
\begin{equation}\label{evii:eq:EVII-z15}
 z_{15}=(X_7+X_{-7})-(X_{20}+X_{-20})-(X_{21}+X_{-21}).
\end{equation}
The $\fp$-highest coefficient spaces satisfy
\[
 \dim(M_4^{\fp},M_2^{\fp},M_0^{\fp})=(2,8,2),
\]
and
\begin{equation}\label{evii:eq:EVII-z15-char}
\begin{array}{c|ccc}
\ell&4&2&0\\ \hline
\det(s-\ad z_{15}|_{M_\ell^{\fp}})&s^2+4&s^8&s^2+4.
\end{array}
\end{equation}
Thus the connected central torus acts by reciprocal nontrivial characters on
$M_4^{\fp}$ and $M_0^{\fp}$ and trivially on $M_2^{\fp}$.
\end{lemma}

\begin{proof}
For orbit $15$ the representative and neutral element are
\[
 E_{15}=\sqrt2\,(X_6+X_{34}+X_{40}+X_{56}),
 \qquad
 H_{0,15}=(4,6,8,12,10,8,4),
\]
with $F_{15}$ obtained by replacing $X_r$ by $X_{-r}$.  Continue to write
$S_r=X_r+X_{-r}$ and $A_r=X_r-X_{-r}$.  For
\[
 \Phi_{15}:\fh\longrightarrow\fg^{\oplus3},\qquad
 x\longmapsto([x,e],[x,d],[x,f]),
\]
exact row reduction gives $\rk\Phi_{15}=64$.  Its kernel has the following
basis:
\begingroup
\small
\begin{align*}
 c_1&=H_3,&c_2&=H_4,&c_3&=X_3,&c_4&=X_4,\\
 c_5&=X_{10},&c_6&=X_{-3},&c_7&=X_{-4},&c_9&=X_{-10},\\
 c_8&=-X_{28}-X_{-2}+X_{-5},
 &c_{10}&=X_{22}-X_{-9}+X_{-11},\\
 c_{11}&=X_{15}-X_{17}+X_{-16},
 &c_{12}&=-X_{16}-X_{-15}+X_{-17},\\
 c_{13}&=-S_7+S_{20}+S_{21},
 &c_{14}&=-X_9+X_{11}+X_{-22},\\
 c_{15}&=X_2-X_5+X_{-28}.
\end{align*}
\endgroup
Substitution verifies that these vectors are fixed by $\theta$ and centralize
the triple.  If $x=\sum_{r=1}^{15}a_rc_r$, the center equations
$[x,c_r]=0$ reduce to
\[
 a_r=0\quad(r\neq13).
\]
Thus the center is the line spanned by $c_{13}=-z_{15}$, proving the first
part of the lemma.

For $\ell=4,2,0$, let
\[
 \Psi_\ell:\fp\longrightarrow\fg\oplus\fg,
 \qquad x\longmapsto([e,x],[d,x]-\ell x).
\]
Its ranks are $52,46,52$, respectively.  Hence the corresponding highest
spaces have dimensions $2,8,2$.  Convenient ordered bases are as follows:
\begingroup
\small
\begin{align*}
 a_4={}&-\sqrt2i(2H_1+H_2+2H_3+2H_4+H_5+2H_7+A_{58}+A_{59})\\
 &\hspace{2.5em}+2(A_6-A_{34}+A_{40}-A_{56}),\\
 b_4={}&2(A_{13}-A_{27}-A_{45}-A_{53})
 +\sqrt2i(-2A_7+A_{20}+A_{21}-A_{49}-A_{63}),\\[1mm]
 w_1={}&2iX_{14}-\sqrt2X_{41}-\sqrt2X_{50}
       +\sqrt2X_{-19}+2iX_{-26}-\sqrt2X_{-33},\\
 w_2={}&-2iX_8+\sqrt2X_{36}+\sqrt2X_{46}
       +\sqrt2X_{-25}+2iX_{-32}-\sqrt2X_{-38},\\
 w_3={}&-2iX_1-\sqrt2X_{30}-\sqrt2X_{43}
       +\sqrt2X_{-31}+2iX_{-37}-\sqrt2X_{-42},\\
 w_4={}&\sqrt2X_{31}+2iX_{37}-\sqrt2X_{42}
       -2iX_{-1}-\sqrt2X_{-30}-\sqrt2X_{-43},\\
 w_5={}&-\sqrt2X_{25}-2iX_{32}+\sqrt2X_{38}
       +2iX_{-8}-\sqrt2X_{-36}-\sqrt2X_{-46},\\
 w_6={}&\sqrt2X_{19}+2iX_{26}-\sqrt2X_{33}
       +2iX_{-14}-\sqrt2X_{-41}-\sqrt2X_{-50},\\
 w_7={}&\sqrt2X_{13}+2iX_{20}-2iX_{21}+\sqrt2X_{27}
       +\sqrt2X_{45}-\sqrt2X_{53}\\
 &+\sqrt2X_{-13}+2iX_{-20}-2iX_{-21}+\sqrt2X_{-27}
       +\sqrt2X_{-45}-\sqrt2X_{-53},\\
 w_8={}&4iH_1+6iH_2+8iH_3+12iH_4+10iH_5+8iH_6+4iH_7\\
 &+\sqrt2\,(X_6+X_{34}+X_{40}+X_{56}
          -X_{-6}-X_{-34}-X_{-40}-X_{-56}),\\[1mm]
 a_0={}&-2H_1-H_2-2H_3-2H_4-H_5+H_7,\\
 b_0={}&-(A_7+A_{20}+A_{21}).
\end{align*}
\endgroup
Thus $(a_4,b_4)$ is a basis of $M_4^{\fp}$,
$(w_1,\ldots,w_8)$ is a basis of $M_2^{\fp}$, and $(a_0,b_0)$ is a basis
of $M_0^{\fp}$.  The remaining bracket calculation is
\begin{align*}
 [z_{15},a_4]&=-2b_4,& [z_{15},b_4]&=2a_4,\\
 [z_{15},w_r]&=0\quad(1\leq r\leq8),\\
 [z_{15},a_0]&=2b_0,& [z_{15},b_0]&=-2a_0.
\end{align*}
Accordingly the three matrices are
\[
 \begin{pmatrix}0&2\\-2&0\end{pmatrix},
 \qquad 0_8,
 \qquad
 \begin{pmatrix}0&-2\\2&0\end{pmatrix},
\]
which proves \eqref{evii:eq:EVII-z15-char} and the asserted character
statement.  As for orbit $20$, only the resulting trivial/nontrivial reciprocal
character pattern is used below; the numerical eigenvalues $\pm2i$ themselves
are not needed.
\end{proof}

\begin{proposition}\label{evii:prop:EVII-orbit15}
Orbit $15$ contributes no adapted distribution.
\end{proposition}

\begin{proof}
As above, \eqref{core:eq:normal-duality-general} inverts the central-torus
characters on passage to the normal quotient.  The two nontrivial characters
in each of $M_4^{\fp}$ and $M_0^{\fp}$ are reciprocal and the $M_2^{\fp}$
characters are trivial, so the multiset in \eqref{evii:eq:EVII-z15-char} is
unchanged.  The defect is six.  The elementary degrees for $\ell=4,2,0$ are
$6,4,2$.  A holomorphic monomial of total weighted degree six has one of the
three patterns
\[
 6,
 \qquad
 4+2,
 \qquad
 2+2+2.
\]
In the first pattern the degree-six factor has a nontrivial torus character.
In the second, the degree-four factor is trivial but the degree-two factor is
nontrivial.  In the third, each degree-two factor has character $\nu$ or
$\nu^{-1}$ for a fixed nontrivial character $\nu$; the sum of three exponents
$\pm1$ is $\pm1$ or $\pm3$, never zero.  Thus there is no invariant
holomorphic monomial of weighted degree six.  Lemmas
\ref{evii:lem:EVII-modular-torus} and \ref{evii:lem:EVII-holomorphic-test} apply.
\end{proof}

\subsection{Completion of speciality}

\begin{proof}[Proof of Theorem~\ref{evii:thm:EVII-special}]
The zero orbit is excluded by Lemma~\ref{core:lem:origin-adapted}, and the
nonzero orbits in \eqref{evii:eq:EVII-strict-orbits} are excluded by the
strict trace inequality.  Proposition~\ref{evii:prop:EVII-congruence} excludes
orbits $17,18,21,22$, Proposition~\ref{evii:prop:EVII-orbit20} excludes orbit
$20$, and Proposition~\ref{evii:prop:EVII-orbit15} excludes orbit $15$.  Therefore
every conormal Frobenius space that could support a $B$-adapted invariant
distribution on the nilpotent cone is zero.  Apply
\cite[Theorem~2.5.6]{AG09} to the action of $H\times\C^\times$ separately
for each of the two adapted homogeneity characters.  The
Kostant--Sekiguchi list above is finite and exhaustive, so the theorem gives
vanishing of every adapted invariant distribution supported on the nilpotent
cone.

Apply Lemma~\ref{core:lem:speciality-criterion} to the one-summand decomposition
$Q(\fp)=\fp$ and the form \eqref{evii:eq:EVII-B}.  This proves speciality; weak
linear tameness and regularity follow.
\end{proof}

\begin{corollary}[EVII component implication]\label{cor:EVII-component}
The finite-component datum for EVII is componentwise regular, and every finite
effective enlargement arising from a central quotient is regular.
\end{corollary}

\begin{proof}
Theorem~\ref{evii:thm:EVII-special} proves vanishing of every connected-group
invariant distribution adapted to the hyperbolic form
\eqref{evii:eq:EVII-B}.  That form is the restriction of the invariant form on
the $\mathfrak e_7$ ideal and is preserved by every inner component.  Apply
Lemma~\ref{lem:special-component} and then
Proposition~\ref{prop:finite-assembly}.
\end{proof}

\section{Completion of the conjectures}\label{sec:completion}

\subsection{A compatible Chevalley anti-automorphism}

We first supply the GP1--GP2 equivalence needed for symmetric subgroups.  It is
enough to prove it on the canonical central cover, where the central torus and
the simply connected semisimple factors are separated.

\begin{proposition}[Compatible Chevalley involution]
\label{prop:compatible-chevalley}
Let
\[
 \widetilde G=Z\times G_{\mathrm{ss}}
\]
with $Z$ a complex torus and $G_{\mathrm{ss}}$ simply connected semisimple, and
let $\theta$ be an involution of $\widetilde G$.  There exists an involutive
algebraic automorphism $C$ such that
\begin{enumerate}[label=\textup{(\roman*)}]
\item $C$ commutes with $\theta$;
\item $C$ restricts to inversion on a maximal torus of $\widetilde G$;
\item the anti-automorphism $\sigma=C\circ\mathrm{inv}$ is
      $\Ad(\widetilde G)$-admissible and satisfies
      $\sigma(\widetilde G^\theta)=\widetilde G^\theta$.
\end{enumerate}
Consequently, for every irreducible Casselman--Wallach representation $E$,
\[
 \dim\Hom_{\widetilde G^\theta}(E,\C)
 =\dim\Hom_{\widetilde G^\theta}(E^\vee,\C).
\]
Thus GP1 and GP2 are equivalent for $(\widetilde G,\widetilde G^\theta)$.
\end{proposition}

\begin{proof}
On $Z$, take $C_Z(z)=z^{-1}$; inversion commutes with every algebraic
automorphism of a torus.  It remains to treat a $\theta$-orbit of simple factors.
For an orbit of size two, write it as $S\times S'$ and choose an isomorphism
$\alpha:S\to S'$ such that
\[
 \theta(x,y)=(\alpha^{-1}(y),\alpha(x)).
\]
Choose a pinning on $S$, transport it by $\alpha$ to $S'$, let $C_S$ be the
corresponding Chevalley involution, and put
$C_{S'}=\alpha C_S\alpha^{-1}$.  Then $C_S\times C_{S'}$ is involutive,
inverts the product torus, and
\[
 (C_S\times C_{S'})\theta(x,y)
 =\bigl(C_S\alpha^{-1}y,\alpha C_Sx\bigr)
 =\theta(C_Sx,C_{S'}y).
\]
Thus it commutes with $\theta$ on the exchanged factor orbit.

For a $\theta$-stable semisimple factor, choose a $\theta$-stable pair
$(B,T)$ and a pinning.  Let $\delta$ be the pinned diagram automorphism with
the same action as $\theta$ on the based root datum.  Then
$\delta^{-1}\theta$ preserves $(B,T)$ and acts trivially on the based root
datum, so the standard structure theorem for pinned automorphisms
\cite{Steinberg} makes it $\operatorname{Int}(t)$ for some $t\in T$.  Thus
\[
 \theta=\operatorname{Int}(t)\delta,
 \qquad t\delta(t)\in Z(G_{\mathrm{ss}}).
\]
Let $C_0$ be the Chevalley involution attached to the pinning.  It inverts $T$
and commutes with $\delta$.  Set
\[
 C=\operatorname{Int}(t)C_0.
\]
Since $C_0(t)=t^{-1}$, one has $C^2=1$.  Moreover
\[
 C\theta C^{-1}
 =\delta\operatorname{Int}(t^{-1})
 =\operatorname{Int}(\delta(t)^{-1})\delta
 =\operatorname{Int}(t)\delta=\theta,
\]
the last equality following from centrality of $t\delta(t)$.  Taking the product
over the factor orbits and with $C_Z$ gives (i) and (ii).

Put $\sigma=C\circ\mathrm{inv}$.  It is involutive.  It normalizes the
conjugation action because
\[
 \sigma\circ\Ad(g)\circ\sigma^{-1}=\Ad(\sigma(g^{-1})).
\]
Every closed conjugacy class is semisimple and meets $T$, while $\sigma$ fixes
$T$ pointwise; hence $\sigma$ preserves every closed conjugacy class.  Thus it
is $\Ad(\widetilde G)$-admissible.  Since $C$ commutes with $\theta$, both $C$
and $\sigma$ preserve the fixed subgroup.

The admissible-anti-automorphism theorem
\cite[Theorem~8.2.1 and Corollary~8.2.3]{AG09} identifies $E^\vee$ with the
$C$-twist of $E$.  Because $C$ preserves $\widetilde G^\theta$, twisting does
not change the dimension of fixed linear forms.  The displayed equality and
the equivalence GP1$\Leftrightarrow$GP2 follow.
\end{proof}

\subsection{Proof of regularity for all complex symmetric pairs}

The following table summarizes the irreducible reduction from Rubio's
Tables~2--4 and Proposition~6.3, together with the odd--odd Spin pleasantness
result of Proposition~\ref{spin:prop:pleasant}.  ``Pleasant'' and ``nice''
imply regularity by the results recalled in Section~\ref{sec:framework}; the
three classical residual families and EVII are proved in this paper.

\begingroup
\small
\renewcommand{\arraystretch}{1.08}
\begin{longtable}{@{}p{0.46\textwidth}p{0.46\textwidth}@{}}
\caption{Regularity of irreducible complex symmetric pairs}
\label{tab:irreducible-reduction}\\
\toprule
irreducible symmetric Lie algebra pair & reason for regularity \\
\midrule
\endfirsthead
\toprule
irreducible symmetric Lie algebra pair & reason for regularity \\
\midrule
\endhead
$(A_{2n},B_n)$ & nice (also pleasant) \\
$(A_{2n-1},D_n)$ & nice \\
$(A_{2n-1},C_n)$ & pleasant \\
$(A_{r+s+1},A_r+A_s+\C)$, $r\ne s$ & pleasant \\
$(A_{2r+1},A_r+A_r+\C)$ & nice \\
$(C_n,A_{n-1}+\C)$ & nice \\
$(C_{r+s},C_r+C_s)$, $r\ne s$ & pleasant \\
$(C_{2r},C_r+C_r)$ & balanced CII, Section~\ref{sec:CII-family} \\
$(D_r,A_{r-1}+\C)$ & DIII, Section~\ref{sec:DIII-family} \\
$(\mathfrak{so}_{r+s},\mathfrak{so}_r+\mathfrak{so}_s)$,
$|r-s|\le2$ & nice \\
$(\mathfrak{so}_{r+s},\mathfrak{so}_r+\mathfrak{so}_s)$,
$|r-s|>2$, $r,s$ odd & pleasant by Proposition~\ref{spin:prop:pleasant} \\
$(\mathfrak{so}_{r+s},\mathfrak{so}_r+\mathfrak{so}_s)$,
$|r-s|>2$, $rs$ even & residual Spin block, Section~\ref{sec:Spin-family} \\
$(G_2,A_1+A_1)$ & nice \\
$(F_4,B_4)$ & pleasant \\
$(F_4,C_3+A_1)$ & nice \\
$(E_6,C_4)$ & nice \\
$(E_6,A_5+A_1)$ & nice \\
$(E_6,F_4)$ & pleasant \\
$(E_6,D_5+\C)$ & pleasant \\
$(E_7,A_7)$ & nice \\
$(E_7,D_6+A_1)$ & pleasant \\
$(E_7,E_6+\C)$ & EVII, Section~\ref{sec:EVII-family} \\
$(E_8,D_8)$ & nice \\
$(E_8,E_7+A_1)$ & pleasant \\
\bottomrule
\end{longtable}
\endgroup

The low-rank accidental identifications
$A_1=B_1=C_1$, $B_2=C_2$, $D_2=A_1+A_1$, and $D_3=A_3$, as well as the
triality automorphisms of $D_4$, only identify entries already present in the
table and create no additional irreducible case.  In particular
$(D_4,A_3+\C)$ is the $r=4$ DIII entry.  At the group level, different
isogeny forms and the possible diagonal coupling of component groups are
handled by Section~\ref{sec:component-assembly}, rather than by adding new
Lie-algebra rows.

\begin{proof}[Proof of Theorem~\ref{thm:regularity-all}]
Pass to the canonical central cover of Lemma~\ref{lem:canonical-cover} and group
the simple factors by the orbits of the lifted involution.  A two-factor orbit
is a diagonal symmetric pair and is pleasant.  For a simple factor fixed by
the permutation induced by $\theta$, Table~\ref{tab:irreducible-reduction}
shows explicitly that it is pleasant or nice except for the four families of
Theorem~\ref{thm:four-inputs}; this is the reduction in
\cite[Tables~2--4 and Proposition~6.3]{Rubio}.  The DIII component implication
is Corollary~\ref{cor:DIII-component}, the balanced CII implication is
Corollary~\ref{cor:CII-component}, the Spin implication is
Corollary~\ref{cor:Spin-component}, and the EVII implication is
Corollary~\ref{cor:EVII-component}.  Pleasant factors have no obstruction
character, and nice factors are covered by Lemma~\ref{lem:special-component}.

Proposition~\ref{prop:component-product} makes the canonical-cover product
componentwise regular.  The central anti-invariant space is removed from
$Q(\fp)$ by Lemma~\ref{lem:fixed-vectors}.  Finally, the original group is a
finite $\theta$-stable central quotient of the cover; its full fixed subgroup is
represented, by Lemma~\ref{lem:exact-component-identification}, by the exact
intermediate component group $B_F$.  Proposition
\ref{prop:finite-assembly} proves regularity for that full fixed subgroup and
for every admissible element.  This proves the theorem.
\end{proof}

\subsection{Generalized descent and finite central descent}

\begin{proof}[Proof of Theorem~\ref{thm:gelfand-all}]
Let
\[
 m:\widetilde G=Z(G)^\circ\times G_{\mathrm{der}}^{\mathrm{sc}}\longrightarrow G
\]
be the canonical central cover and let $\widetilde\theta$ be the lifted
involution.  The centralizer of a semisimple element in a simply connected
semisimple complex group is connected by Steinberg's connected-centralizer
theorem \cite{Steinberg}; hence every
semisimple centralizer in $\widetilde G$ is connected.  Every descendant of
$(\widetilde G,\widetilde G^{\widetilde\theta})$ is therefore a connected
complex reductive symmetric pair and is regular by
Theorem~\ref{thm:regularity-all}.

Every connected complex symmetric pair is good
\cite[Corollary~7.1.7]{AG09}.  The generalized Harish--Chandra descent
criterion \cite[Theorem~7.4.5]{AG09} now shows that the cover pair is a
Gelfand--Kazhdan pair, hence has GP2.  Proposition
\ref{prop:compatible-chevalley} converts GP2 into GP1 on the cover.

Let $E$ be an irreducible Casselman--Wallach representation of $G$ and inflate
it along $m$.  Pullback through a finite covering preserves smoothness,
Fr\'echet moderate growth, admissibility, and finite length.  It also preserves
irreducibility: a closed $\widetilde G$-invariant subspace is $G$-invariant
because $m$ is surjective.  If
$\widetilde H=\widetilde G^{\widetilde\theta}$, then
$m(\widetilde H)\subset H$, and therefore
\[
 \Hom_H(E,\C)\subset
 \Hom_{m(\widetilde H)}(E,\C)
 \simeq\Hom_{\widetilde H}(E\circ m,\C).
\]
The right-hand space has dimension at most one.  Thus $(G,H)$ has GP1.  If
$\pi$ is an irreducible admissible unitary representation, its $K$-finite
vectors form an irreducible Harish--Chandra module; the Casselman--Wallach
globalization theorem identifies its smooth vectors with the irreducible
smooth Fr\'echet globalization of that module
\cite[Theorem~10.6]{BernsteinKrotz}.  Therefore GP1 implies the unitary
multiplicity-one property and proves van Dijk's conjecture.
\end{proof}

\begin{remark}
The central-cover argument is used only for the multiplicity-one conclusion.
Theorem~\ref{thm:regularity-all} itself is proved for the full fixed subgroup of
every connected central quotient by the finite-component assembly of
Section~\ref{sec:component-assembly}.
\end{remark}

\begin{thebibliography}{99}

\bibitem{AGSchwartz}
A.~Aizenbud and D.~Gourevitch,
\emph{Schwartz functions on Nash manifolds},
Int. Math. Res. Not. IMRN 2008, Art. ID rnm155, 37 pp.

\bibitem{AG09}
A.~Aizenbud and D.~Gourevitch,
\emph{Generalized Harish--Chandra descent, Gelfand pairs, and an Archimedean
analog of Jacquet--Rallis' theorem},
with an appendix by A.~Aizenbud, D.~Gourevitch, and E.~Sayag,
Duke Math. J. \textbf{149} (2009), no.~3, 509--567.

\bibitem{AGS08}
A.~Aizenbud, D.~Gourevitch, and E.~Sayag,
\emph{$(\GL_{n+1}(F),\GL_n(F))$ is a Gelfand pair for any local field $F$},
Compos. Math. \textbf{144} (2008), no.~6, 1504--1524.

\bibitem{Aiz13}
A.~Aizenbud,
\emph{A partial analog of the integrability theorem for distributions on
$p$-adic spaces and applications},
Israel J. Math. \textbf{193} (2013), no.~1, 233--262.

\bibitem{BernsteinKrotz}
J.~Bernstein and B.~Kr\"otz,
\emph{Smooth Fr\'echet globalizations of Harish--Chandra modules},
Israel J. Math. \textbf{199} (2014), no.~1, 45--111.

\bibitem{Borel}
A.~Borel,
\emph{Sous-groupes commutatifs et torsion des groupes de Lie compacts connexes},
Tohoku Math. J. (2) \textbf{13} (1961), 216--240.

\bibitem{Bulois}
M.~Bulois,
\emph{Composantes irr\'eductibles de la vari\'et\'e commutante nilpotente d'une
alg\`ebre de Lie sym\'etrique semi-simple},
Ann. Inst. Fourier (Grenoble) \textbf{59} (2009), no.~1, 37--80.

\bibitem{ChenSun}
Y.~Chen and B.~Sun,
\emph{Schwartz homologies of representations of almost linear Nash groups},
J. Funct. Anal. \textbf{280} (2021), no.~7, Paper No.~108817, 50 pp.

\bibitem{CM}
D.~H. Collingwood and W.~M. McGovern,
\emph{Nilpotent Orbits in Semisimple Lie Algebras},
Van Nostrand Reinhold, New York, 1993.

\bibitem{Djokovic}
D.~\v{Z}.~Djokovi\'c,
\emph{Explicit Cayley triples in real forms of $E_7$},
Pacific J. Math. \textbf{191} (1999), no.~1, 1--23.

\bibitem{JSZ}
D.~Jiang, B.~Sun, and C.-B.~Zhu,
\emph{Uniqueness of Ginzburg--Rallis models: the Archimedean case},
Trans. Amer. Math. Soc. \textbf{363} (2011), no.~5, 2763--2802.

\bibitem{KR}
B.~Kostant and S.~Rallis,
\emph{Orbits and representations associated with symmetric spaces},
Amer. J. Math. \textbf{93} (1971), 753--809.

\bibitem{Przebinda}
T.~Przebinda,
\emph{The oscillator duality correspondence for the pair
$O(2,2),\,Sp(2,\R)$},
Mem. Amer. Math. Soc. \textbf{79} (1989), no.~403.

\bibitem{Rubio}
R.~Rubio,
\emph{On the Gelfand property for complex symmetric pairs},
Trans. Amer. Math. Soc. \textbf{379} (2026), no.~3, 1505--1540.

\bibitem{Sekiguchi}
J.~Sekiguchi,
\emph{Invariant spherical hyperfunctions on the tangent space of a symmetric
space},
in \emph{Algebraic groups and related topics}, Adv. Stud. Pure Math.
\textbf{6}, North-Holland, Amsterdam, 1985, 83--126.

\bibitem{Steinberg}
R.~Steinberg,
\emph{Endomorphisms of Linear Algebraic Groups},
Memoirs of the American Mathematical Society, no.~80,
American Mathematical Society, Providence, RI, 1968.

\end{thebibliography}

\end{document}
